\documentclass[11pt]{article}
\usepackage{amsmath, amssymb, amsfonts, mathtools}
%\usepackage{hyperref}
\usepackage[margin=1in]{geometry}
%\usepackage{showlabels}
\newcommand{\tr}{\mathrm{tr}}
\newcommand{\ii}{\mathrm{i}}
\newcommand{\ep}{\epsilon}
\newcommand{\p}{\partial}
\newcommand{\dd}{\mathrm{d}}
\newcommand{\eps}{\varepsilon}
\newcommand{\mO}{\mathcal{O}}
\usepackage{
  jheppub
} % for details on the use of the package, please see the JINST-\
\newcommand{\nn}{\nonumber}
\usepackage{braket}
\newcommand{\lb}{\left(}
\newcommand{\rb}{\right)}
\usepackage{amsmath, amssymb}
\newcommand{\NL}[1]{\textcolor{blue}{#1}}
\usepackage{xcolor}

\newcommand{\JM}[1]{\textcolor{orange}{#1}}
\title{Relative entropy of mixtures of excited states in CFT}
\author{Nima Lashkari, Jignesh Mohanty to UIUC colleagues}
\affiliation{Department of Physics, Purdue University, West Lafayette, IN 47907, USA}
\emailAdd{nlashkari@purdue.edu}
%\abstract{We study the relative entropy of mixtures of excited states in conformal field theories in $1+1$ dimensions. We consider two energy eigenstates of a CFT $\ket{\mO_\alpha}$ and $\ket{\mO_\beta}$ corresponding to Virasoro primaries $\mO_{\alpha}$ and $\mO_{\beta}$. }


\begin{document}

\maketitle
\flushbottom
  \section{Introduction and Motivation}

  Consider a 2d CFT, and two energy eigenstates of $\ket{\mO_\alpha}$ and
  $\ket{\mO_\beta}$ corresponding to spinless Virasoro primaries $\Phi$ and
  $\Psi$. Denote by $\rho_\Phi$ and $\rho_\Psi$, the reduced density matrices of these states to a region $A$, and consider the mixture state $\rho_y=y\rho_\Phi+(1-y)\rho_\Psi$ for some $0\leq y\leq 1$. Motivated by connections to quantum error correction codes \cite{}, in this note, we study the relative entropies $S(\rho_\Phi\|\rho_y,A)$ in the small and large $|A|$ limits.
  
  
  % that for the mixture $(\rho_{\alpha}+\rho_{\beta})/2$ we
  % have the following asymptotic behavior of relative entropy
  % \begin{eqnarray}
  %   &&S(x):=S(\rho_\alpha|(\rho_\alpha+\rho_\beta)/2)\nn\\
  %   &&S(x)\simeq A_0 x^{\delta_0}, \qquad S(1-x)\simeq A_1(1-x)^{\delta_1},\nn\\
  %   &&\delta_1\leq \delta_0\ .
  % \end{eqnarray}
  % In this note, we check this conjecture. In particular, we will compute $\delta_1$ and $\delta_0$ in the special case of massless 2d free boson CFT and energy eigenstates created by the action of the vertex operator
  % states $V_{\alpha}=e^{i\alpha\phi}$ and $V_{\beta}= e^{i\beta\phi}$ at the
  % center of radial quantization. 

\section{Relative entropy of excited states}

\subsection{Setup and Replica Trick}

  Consider a conformal field theory (CFT) in $1+1$ dimensions on a cylinder
  $\mathcal{C}_{1}$ of circumference $R$. We focus on a spatial subsystem $A$ defined
  by the interval $(-l/2, l/2)$ at a fixed time slice. We simplify our notation
  by taking $A=(-\theta,\theta)$ on a cylinder of circumference $2\pi$ such that
  $\theta=\pi x$.

  Our states of interest are the reduced density matrices corresponding to
  energy eigenstates corresponding to the insertion of Virasoro primary fields at the center of radial quantization. These states are generated by the action of 
  % vertex operators $V_{\alpha}=:e^{i\alpha \phi}$
  % \JM{(
  operators $\Psi$ and $\Phi$ with conformal weights $(h_\Psi, \bar h_\Psi)$ and $(h_\Phi, \bar h_\Phi)$
  % )}: 
  % vertex operators $V_{\alpha}=:e^{i\alpha \phi}:$ 
  on the vacuum at past infinity in imaginary
  time on the Euclidean cylinder $\mathcal{C}_{1}$:
  % \begin{equation}
  %   |\alpha\rangle = \lim_{u\to i\infty}\frac{V_{\alpha}(u)|\Omega\rangle}{\braket{V_{-\alpha}(-u)V_\alpha(u)}^{1/2}}
  % \end{equation}
    \begin{equation}
    % \JM{
    |\Psi\rangle = \lim_{u\to i\infty}\frac{\Psi(u)|\Omega\rangle}{\braket{\Psi^*(-u)\Psi(u)}^{1/2}} \qquad |\Phi\rangle = \lim_{u\to i\infty}\frac{\Phi(u)|\Omega\rangle}{\braket{\Phi^*(-u)\Phi(u)}^{1/2}}
    % }
  \end{equation}
  % \JM{
  While considering the special case of massless free boson theory, we will replace these operators with Vertex operators.
  % }
  The reduced density matrix of the state   
  % $|\alpha\rangle$ 
  % \JM{(
  $|\Psi \rangle$
  % )}
  on the subsystem $A$
  is denoted by 
  % $\rho_{\alpha}$ \JM{(
  $\rho_\Psi$.
  % )}. 
  % A simple way to keep track of normalization of the
  % density matrix to define normalized operators
  % \begin{eqnarray}
  %   \tilde{V}_{\alpha}(u^+)\tilde{V}_{-\alpha}(u^-)=\frac{V_{\alpha}(u^{+})V_{-\alpha}(u^{-})}{\braket{V_{\alpha}(u^+)V_{-\alpha}(u^-)}_{\mathcal{C}_1}}
  % \end{eqnarray}
The replica trick that compute the relative entropy of excited states is
  % \begin{align}
  %   S\left(\rho_\alpha \Big\| \rho_\beta\right)&=\lb \partial_{n}
  %   \log \left[ \frac{\text{tr}(\rho_{\alpha}^{n})}{\text{tr}(\rho_{\alpha}\rho_{\beta}^{n-1})}
  %   \right]\rb_{n\to 1}\nn\\
  %   & =\lb\p_n\log\left[
  %   \frac{\braket{\prod_{i=1}^{n} V_{\alpha}(u^+_i)V_{-\alpha}(u^-_i)}_{\mathcal{C}_n}\braket{V_\beta(u^+)V_{-\beta}(u^-)}_{\mathcal{C}_1}^{n-1}}{\braket{V_\alpha(u^+_n)V_{-\alpha}(u^-_n)\mO^{(n-1)}_{\beta}}_{\mathcal{C}_n}\braket{V_\alpha(u^+)V_{-\alpha}(u^-)}_{\mathcal{C}_1}^{n-1}}\right]\rb_{n=1}\label{corrFncyl}\nn\\
  % \end{align}
  % where \JM{(this contains just the holomorphic part)}
  % \begin{eqnarray}
  %   &&\mO_{\beta}^{(n-1)}=
  %   V_{\beta}(u_1^+)V_{-\beta}(u_1^-)\cdots V_{\beta}(u_{n-1}^+)V_{-\beta}(u_{n-1}^-).
  % \end{eqnarray}
  \begin{align}\label{relEntropyReplica}
    S\left(\rho_\Psi \Big\| \rho_\Phi\right)&=\lb \partial_{n}
    \log \left[ \frac{\text{tr}(\rho_{\Psi}^{n})}{\text{tr}(\rho_{\Psi}\rho_{\Phi}^{n-1})}
    \right]\rb_{n\to 1}\nn\\
    % & =\lb\p_n\log\left[
    % \frac{\braket{\prod_{i=1}^{n} V_{\alpha}(u^+_{i,n}, \bar u_{i,n}^+)V_{-\alpha}(u^-_{i,n},\bar u_{i,n}^-)}_{\mathcal{C}_n}\braket{V_\beta(u^+,\bar u^+)V_{-\beta}(u^-,\bar u^-)}_{\mathcal{C}_1}^{n-1}}{\braket{V_\alpha(u^+_{n,n}, \bar u_{n,n}^+)V_{-\alpha}(u^-_{n,n},\bar u_{n,n}^-)\mO^{(n-1)}_{\beta}}_{\mathcal{C}_n}\braket{V_\alpha(u^+,\bar u^+)V_{-\alpha}(u^-,\bar u^-)}_{\mathcal{C}_1}^{n-1}}\right]\rb_{n=1}\label{corrFncyl}\nn\\
    % \JM{S(\rho_\Psi \| \rho_\Phi)}
    &=\lb
    % \JM{
    \p_n\log\left[
    \frac{\braket{\prod_{i=1}^{n} \Psi(u^+_{i,n}, \bar u_{i,n}^+)\Psi^*(u^-_{i,n},\bar u_{i,n}^-)}_{\mathcal{C}_n}\braket{\Phi(u^+,\bar u^+)\Phi^*(u^-,\bar u^-)}_{\mathcal{C}_1}^{n-1}}{\braket{\Psi(u^+_{n,n}, \bar u_{n,n}^+)\Psi^*(u^-_{n,n},\bar u_{n,n}^-)\mO^{(n-1)}_{\Phi}}_{\mathcal{C}_n}\braket{\Psi(u^+,\bar u^+)\Psi^*(u^-,\bar u^-)}_{\mathcal{C}_1}^{n-1}}\right]
    % }
    \rb_{n=1}\\
  \end{align}
  % \begin{eqnarray*}
  %   &&\mO_{\beta}^{(n-1)}=
  %   V_{\beta}(u_{1,n}^+,\bar u_{1,n}^+)V_{-\beta}(u_{1,n}^-, \bar u_{1,n}^-)\cdots V_{\beta}(u_{n-1,n}^+,\bar u_{n-1,n}^+)V_{-\beta}(u_{n-1,n}^-,\bar u_{n-1,n}^-).
  % \end{eqnarray*}
  \begin{eqnarray}\label{corrFncyl}
    % \JM{
    \mO_{\Phi}^{(n-1)}=
    \Phi(u_{1,n}^+,\bar u_{1,n}^+)\Phi^*(u_{1,n}^-, \bar u_{1,n}^-)\cdots \Phi(u_{n-1,n}^+,\bar u_{n-1,n}^+)\Phi^*(u_{n-1,n}^-,\bar u_{n-1,n}^-).
    % }
  \end{eqnarray}
   The conformal map
  \begin{eqnarray}
    z=\lb\frac{\sin((u-\theta)/2)}{\sin((u+\theta)/2)}\rb^{1/n}
  \end{eqnarray}
  sends the interval $A=(-\theta,\theta)$ to $A=(0,\infty)$ uniformizing
  $\mathcal{C}_{n}$ to the complex plane. Under this transformation Virasoro primaries
  transform as
  \begin{eqnarray}
    &&\mO(u,\bar{u})=\lb\frac{\partial z}{\partial u} \rb^{h_\mO}\lb\frac{\partial
    \bar{z}}{\partial\bar{u}}
    \rb^{\bar{h}_\mO}\mO(z,\bar{z})\nn\\
    % &&V_\alpha(u_{k,n},\bar{u}_{k,n})=\lb\frac{z_{k,n}\Lambda}{n}\rb^{h_\alpha}\lb\frac{\bar{z}_{k,n}\Lambda}{n}\rb^{\bar{h}_\alpha}V_\alpha(z_{k,n},\bar{z}_{k,n}), \nn \\
    % &&\JM{
    &&\Psi(u_{k,n},\bar{u}_{k,n})=\lb\frac{z_{k,n}\Lambda}{n}\rb^{h_\alpha}\lb\frac{\bar{z}_{k,n}\Lambda}{n}\rb^{\bar{h}_\alpha}\Psi(z_{k,n},\bar{z}_{k,n}),
    % }
  \end{eqnarray}
  where in our example
  \begin{eqnarray}
    z^\pm_{k,n}=e^{\frac{i\pi}{n}(2k\pm x)},\qquad \Lambda=\frac{\sin(\theta)}{\cos(\theta)-\cos(u)}
  \end{eqnarray}
  We define 
%   \begin{align}
%   \mathcal{J}_{n,\alpha}&=\prod_{k=1}^{n-1}\lb \frac{\partial z^+_{k,n}}{\partial u^+_{k,n}}\rb^{h_\alpha}\lb \frac{\partial \bar{z}^+_{k,n}}{\partial \bar{u}^+_{k,n}}\rb^{\bar{h}_\alpha}\lb \frac{\partial z^-_{k,n}}{\partial u^-_{k,n}}\rb^{h_\alpha}\lb \frac{\partial \bar{z}^-_{k,n}}{\partial \bar{u}^-_{k,n}}\rb^{\bar{h}_\alpha}\nn\\
%   &=\prod_{k=1}^{n-1}\lb \frac{z^+_{k,n}z^-_{k,n}\Lambda^2}{n^2}\rb^{h_\alpha}\lb \frac{\bar{z}^+_{k,n}\bar{z}^-_{k,n}\Lambda^2}{n^2}\rb^{\bar{h}_\alpha}\nn\\
%   &=\prod_{k=1}^{n-1} (\Lambda/n)^{2(h_\alpha+\bar{h}_\alpha)}e^{4\pi i (h_\alpha-\bar{h}_\alpha)k/n}\nn\\
% %  &=n^{-2(n-1)(h_\alpha+\bar{h}_\alpha)}e^{2\pi i (n-1)(h_\alpha-\bar{h}_\alpha)}e^{-2\pi i (n-1)\bar{h}_\alpha}\nn\\
%   &=(\Lambda/n)^{2(n-1)\Delta_\alpha}e^{2\pi i(n-1) s_\alpha}
%   %=n^{-2(n-1)(\Delta_\alpha+\Delta_\beta)}
%   \end{align}
  \begin{align}
  \mathcal{J}_{n,\Psi}&=\prod_{k=1}^{n-1}\lb \frac{\partial z^+_{k,n}}{\partial u^+_{k,n}}\rb^{h_\Psi}\lb \frac{\partial \bar{z}^+_{k,n}}{\partial \bar{u}^+_{k,n}}\rb^{\bar{h}_\Psi}\lb \frac{\partial z^-_{k,n}}{\partial u^-_{k,n}}\rb^{h_\Psi}\lb \frac{\partial \bar{z}^-_{k,n}}{\partial \bar{u}^-_{k,n}}\rb^{\bar{h}_\Psi}\nn\\
  &=\prod_{k=1}^{n-1}\lb \frac{z^+_{k,n}z^-_{k,n}\Lambda^2}{n^2}\rb^{h_\Psi}\lb \frac{\bar{z}^+_{k,n}\bar{z}^-_{k,n}\Lambda^2}{n^2}\rb^{\bar{h}_\Psi}\nn\\
  &=\prod_{k=1}^{n-1} (\Lambda/n)^{2(h_\Psi+\bar{h}_\Psi)}e^{4\pi i (h_\Psi-\bar{h}_\Psi})k/n\nn\\
%  &=n^{-2(n-1)(h_\alpha+\bar{h}_\alpha)}e^{2\pi i (n-1)(h_\alpha-\bar{h}_\alpha)}e^{-2\pi i (n-1)\bar{h}_\alpha}\nn\\
  &=(\Lambda/n)^{2(n-1)\Delta_\Psi}e^{2\pi i(n-1) s_\Psi}
  %=n^{-2(n-1)(\Delta_\alpha+\Delta_\beta)}
  \end{align}
The overall conformal factor that captures the transformation of the correlator in (\ref{corrFncyl}) from the cylinder to the complex plane is
% \begin{align}
% \mathcal{J}&=\frac{\mathcal{J}_{n,\alpha}}{\mathcal{J}_{n,\beta}}\lb z^+_{1,1}z^-_{1,1}\Lambda^2\rb^{(n-1)(h_\beta-h_\alpha)}\lb\bar{z}^+_{1,1}\bar{z}^-_{1,1}\Lambda^2\rb^{(n-1)(\bar{h}_\beta-\bar{h}_\alpha)}\nn\\
% &=n^{-2(n-1)(\Delta_\alpha-\Delta_\beta)}e^{2\pi i(n-1)(s_\alpha-s_\beta)}
% \end{align}
\begin{align}
\mathcal{J}&=\frac{\mathcal{J}_{n,\Psi}}{\mathcal{J}_{n,\Phi}}\lb z^+_{1,1}z^-_{1,1}\Lambda^2\rb^{(n-1)(h_\Phi-h_\Psi)}\lb\bar{z}^+_{1,1}\bar{z}^-_{1,1}\Lambda^2\rb^{(n-1)(\bar{h}_\Phi-\bar{h}_\Psi)}\nn\\
&=n^{-2(n-1)(\Delta_\Psi-\Delta_\Phi)}e^{2\pi i(n-1)(s_\Psi-s_\Phi)}
\end{align}


 % where in the last step we have used the fact that $s_\alpha-s_\beta\in \mathbb{Z}$.
Therefore, mapping the correlator to the complex plane, we have 
% \begin{align}
%   S(\rho_\alpha\|\rho_\beta)=&\lb\p_n\log\mathcal{F}_n\rb_{n=1}\nn\\
%   \mathcal{F}_n&=\left[\frac{\braket{\prod_{i=1}^{n} V_{\alpha}(z^+_i)V_{-\alpha}(z^-_i)}\braket{V_\beta(z^+)V_{-\beta}(z^-)}^{n-1}}{\braket{V_\alpha(z^+_n)V_{-\alpha}(z^-_n)\tilde{\mO}^{(n-1)}_{\beta}}\braket{V_\alpha(z^+)V_{-\alpha}(z^-)}^{n-1}}\right]\lb \frac{e^{2\pi i (s_\alpha-s_\beta)}}{n^{2(\Delta_\alpha-\Delta_\beta)}}\rb^{n-1}\label{fncorrspin}\\
%   \tilde{\mO}^{(n-1)}_\beta&= \prod_{k=1}^{n-1} V_{\beta}(z^+_k)V_{-\beta}(z^-_k)
% \end{align}
\begin{align}
  S(\rho_\Psi\|\rho_\Phi)=&\lb\p_n\log\mathcal{F}_n\rb_{n=1}\nn\\
  \mathcal{F}_n&=\left[\frac{\braket{\prod_{i=1}^{n} \Psi (z^+_i)\Psi^(z^-_i)}\braket{\Phi(z^+)\Phi^*(z^-)}^{n-1}}{\braket{\Psi(z^+_n)\Psi^*(z^-_n)\tilde{\mO}^{(n-1)}_{\Phi}}\braket{\Psi(z^+)\Psi^*(z^-)}^{n-1}}\right]\lb \frac{e^{2\pi i (s_\Psi-s_\Phi)}}{n^{2(\Delta_\Psi-\Delta_\Phi)}}\rb^{n-1}\label{fncorrspin}\\
  \tilde{\mO}^{(n-1)}_\Phi&= \prod_{k=1}^{n-1} \Phi(z^+_k)\Phi^*(z^-_k)
\end{align}
% \NL{Add both $z$ and $\bar{z}$ dependence}
% \begin{align*}
%   S(\rho_\alpha\|\rho_\beta)=&\lb\p_n\log\mathcal{F}_n\rb_{n=1}\nn\\
%   \mathcal{F}_n&=\left[\frac{\braket{\prod_{i=1}^{n} V_{\alpha}(z^+_{i,n},\bar z^+_{i,n})V_{-\alpha}(z^-_{i,n},\bar z^-_{i,n})}\braket{V_\beta(z^+,\bar z^+)V_{-\beta}(z^-,\bar z^-)}^{n-1}}{\braket{V_\alpha(z^+_{n,n}, \bar z^+_{n,n})V_{-\alpha}(z^-_{n,n}, \bar z^-_{n,n})\tilde{\mO}^{(n-1)}_{\beta}}\braket{V_\alpha(z^+,\bar z^+)V_{-\alpha}(z^-,\bar z^-)}^{n-1}}\right]\lb \frac{e^{2\pi i (s_\alpha-s_\beta)}}{n^{2(\Delta_\alpha-\Delta_\beta)}}\rb^{n-1}\label{fncorrspin}\\
%   \tilde{\mO}^{(n-1)}_\beta&= \prod_{k=1}^{n-1} V_{\beta}(z^+_{k,n},\bar z^+_{k,n})V_{-\beta}(z^-_{k,n},\bar z^-_{k,n})
% \end{align*}
\begin{align}
  S(\rho_\Psi\|\rho_\Phi)=&\lb\p_n\log\mathcal{F}_n\rb_{n=1}\nn\\
  \mathcal{F}_n&=\left[\frac{\braket{\prod_{i=1}^{n} \Psi(z^+_{i,n},\bar z^+_{i,n})\Psi^*(z^-_{i,n},\bar z^-_{i,n})}\braket{\Phi(z^+,\bar z^+)\Phi^*(z^-,\bar z^-)}^{n-1}}{\braket{\Psi(z^+_{n,n}, \bar z^+_{n,n})\Psi^*(z^-_{n,n}, \bar z^-_{n,n})\tilde{\mO}^{(n-1)}_{\Phi}}\braket{\Psi(z^+,\bar z^+)\Psi^*(z^-,\bar z^-)}^{n-1}}\right]\lb \frac{e^{2\pi i (s_\Psi-s_\Phi)}}{n^{2(\Delta_\Psi-\Delta_\Phi)}}\rb^{n-1}\label{fncorrspin}\\
  \tilde{\mO}^{(n-1)}_\Phi&= \prod_{k=1}^{n-1} \Phi(z^+_{k,n},\bar z^+_{k,n})\Phi^*(z^-_{k,n},\bar z^-_{k,n})
\end{align}
Since $\Psi \Psi^*$ and $\Phi \Phi^*$ OPEs are invariant under a rotation of $\pi$ about its midpoint \cite{calabrese2011entanglement}, the total conformal spin should be an even integer. This requires $2s_\Psi, 2s_\Phi \in 2\mathbb{Z}$. Thus we have, $s_\Psi, s_\Phi \in \mathbb{Z}$. 
Since 
% $s_\alpha-s_\beta$ (\JM{
$s_\Psi-s_\Phi$
% })
is an integer, we drop the 
% $e^{2\pi i (s_\alpha-s_\beta)}$ (\JM{
$e^{2\pi i (s_\Psi-s_\Phi)}$ 
% }) 
term. \NL{Is this statement correct?}
% where we have used
% \begin{align}
%   \lb\p_n\log\left[n^{-2(n-1)(\Delta_\alpha+\Delta_\beta)}\right]\rb_{n\to 1}=0.
% \end{align}


% \NL{HERE}
%   Therefore, for normalized operators $\tilde{V}$ the contribution of $\Lambda$
%   terms cancel and we find\footnote{For scalar (spinless) operators in $2d$ we have $h_{\alpha}=\bar{h}_{\alpha}$
%   and we find
%   $\tilde{V}_{\alpha}(u_{k},\bar{u}_{k})=n^{-2h_\alpha}  \tilde{V}_{\alpha}(z_{k,n}
%     ,\bar{z}_{k,n})$.}
%   \begin{eqnarray}
%     \tilde{V}_\alpha(u^\pm_k,\bar{u}^\pm_k)&&=\lb\frac{z_{k,n}}{n z_{1,1}}\rb^{h_\alpha}\lb\frac{\bar{z}_{k,n}}{n
%     z_{1,1}}\rb^{\bar{h}_\alpha}\tilde{V}_\alpha(z_{k,n},\bar{z}_{k,n})\nn\\
%     &&=\lb\frac{e^{i\pi(2k\pm x)(1-n)/n}}{n }\rb^{h_\alpha}\lb\frac{e^{-i\pi(2k\pm
%     x)(1-n)/n}}{n}\rb^{\bar{h}_\alpha}\tilde{V}_\alpha(z^\pm_{k,n},\bar{z}^\pm_{k,n})
%   \end{eqnarray}
%   Therefore,
%   \begin{align}
%     \tilde{V}_\alpha(u^+_k,\bar{u}^+_k)\tilde{V}_{\alpha}(u^-_k,\bar{u}^-_k)=n^{-2\Delta_\alpha}e^{4\pi i k s_\alpha(n-1)/n}\tilde{V}_\alpha(z^+_{k,n},\bar{z}^+_{k,n})\tilde{V}_\alpha(z^-_{k,n},\bar{z}^-_{k,n})
%   \end{align}
% The relative entropy replica trick in terms of the correlator on the complex plane gives
%   \begin{align}\label{CorrFn}
%     S(\rho_\alpha\|\rho_\beta)&=\lb \partial_n\log\mathcal{F}_n(\alpha|\beta)\rb_{n\to 1}\nn\\
%     \mathcal{F}_n(\alpha|\beta)&=
%     \frac{\braket{\prod_{k=1}^{n} \tilde{V}_{\alpha}(z^+_{k,n})\tilde{V}_{-\alpha}(z^-_{k,n})}}{\braket{\tilde{V}_\alpha(z^+_{n,n})\tilde{V}_{-\alpha}(z^-_{n,n})\tilde{\mO}^{(n-1)}_{\beta}}}\nn\\
%     \tilde{\mO}_{\beta}^{(n-1)}&=n^{2(n-1)(\Delta_\alpha-\Delta_\beta)}\prod_{k=1}^{n-1}\lb e^{4\pi i k (s_\beta-s_\alpha)(n-1)/n}\frac{V_{\beta}(z_{k,n}^{+})V_{-\beta}(z_{k,n}^{-})}{\braket{V_{\beta}(z^+_{1,1})V_{-\beta}(z^-_{1,1})}}\rb\nn\\
%     &=n^{2(n-1)(\Delta_\alpha-\Delta_\beta)}e^{2\pi i (s_\beta-s_\alpha)(n-1)^2}\prod_{k=1}^{n-1}\frac{V_{\beta}(z_{k,n}^{+})V_{-\beta}(z_{k,n}^{-})}{\braket{V_{\beta}(z^+_{1,1})V_{-\beta}(z^-_{1,1})}}.
%   \end{align}
%   Note that the term $e^{2\pi i (s_\alpha-s_\beta)(n-1)^2}$ term does not contribute to relative entropy because
%   \begin{align}
%       &-\p_n\lb 2\pi i (s_\alpha-s_\beta)(n-1)^2\rb_{n\to 1}=0
%   \end{align}
%   As a result, we can drop this term from the denominator. 
 % We will specialize to the case where $s_\alpha=s_\beta$ so that we can drop the contribution of the spin.

  % \subsection{Small interval limit $x\to 0^+$}
\subsection{Small interval limit \texorpdfstring{$x\to 0^+$}{x -> 0+}}
Label by $\mO_L$, $\mO_{L''}$, and $\mO_{L'}$ the lightest Virasoro primaries that appear, respectively, in the OPE of $\Psi \Psi^*$, $\Phi\Phi^*$, and $\Psi\Phi^*$\footnote{Otherwise the lowest order
  term would come from stress tensor.}. 
  To simplify our notation, we define
   % \JM{(
   $C_{\Psi} = C^{L}_{\Psi\Psi^*}$, $C_{\Phi} = C^{L''}_{\Phi\Phi^*}$ and $C_{\Psi \Phi} = C^{L'}_{\Psi \Phi^*}=C^{L'}_{\Psi^*\Phi}$\footnote{Note that $C^{L'}_{\Psi\Phi^*}=(C^{L'}_{\Psi^*\Phi})^*$, but in a unitary CFT, the OPE coefficients are real.}.
   % )
   % } 
   % $C_{\omega}=C^{L}_{\omega(-\omega)}$ for $\omega=\alpha,\beta$ and $C_{\alpha\beta}=C^L_{\alpha\beta}$ 
   In the small $x$ limit, we can expand in the OPE
  channel
  % \begin{align}
  %   \frac{V_{w_k}(z_{k}^{+})V_{-w_k}(z_{k}^{-})}{\braket{V_{w_k}(z^+)V_{-w_k}(z^-)}}&=e^{-4\pi i k s_{w_k}/n}\lb
  %   \frac{\sin(\pi x)}{\sin(\pi x/n)}\rb^{2\Delta_{w_k}}\times\nn\\
  %   &\quad\lb 1+ C_{w_k}(2\sin(\pi x/n))^{\Delta_L}e^{\pi i s_L(2k/n+1/2)}\mO_L(e^{2\pi i k/n})+\cdots\rb
  % \end{align}
  \begin{align}
    \frac{\Psi(z_{k}^{+})\Psi^*(z_{k}^{-})}{\braket{\Psi(z^+)\Psi^*(z^-)}}&=e^{-4\pi i k s_{\Psi}/n}\lb
    \frac{\sin(\pi x)}{\sin(\pi x/n)}\rb^{2\Delta_{\Psi}}\times\nn \\
    &\quad\lb  1+ C_{\Psi}(2\sin(\pi x/n))^{\Delta_L}e^{\pi i s_L(2k/n+1/2)}\mO_L(e^{2\pi i k/n})+\rb  \\
    \frac{\Phi(z_{k}^{+})\Phi^*(z_{k}^{-})}{\braket{\Phi(z^+)\Phi^*(z^-)}}&=e^{-4\pi i k s_{\Phi}/n}\lb
    \frac{\sin(\pi x)}{\sin(\pi x/n)}\rb^{2\Delta_{\Phi}}\times\nn\\
    &\quad\lb  1+ C_{\Phi}(2\sin(\pi x/n))^{\Delta_{L''}}e^{\pi i s_{L''}(2k/n+1/2)}\mO_{L''}(e^{2\pi i k/n})+\cdots\rb 
  \end{align}
  % \begin{align*}
  %   \JM{\frac{V_{w_k}(z_{k,n}^{+},\bar z_{k,n}^+)V_{-w_k}(z_{k,n}^{-},\bar z_{k,n}^-)}{\braket{V_{w_k}(z^+)V_{-w_k}(z^-)}}}&=e^{-4\pi i k s_{w_k}/n}\lb
  %   \frac{\sin(\pi x)}{\sin(\pi x/n)}\rb^{2\Delta_{w_k}}\times\nn\\
  %   &\quad\lb 1+ C_{w_k}(2\sin(\pi x/n))^{\Delta_L}e^{\pi i s_L(2k/n+1/2)}\mO_L(e^{2\pi i k/n})+\cdots\rb
  % \end{align*}
  In our notation, the Virasoro primary $\mO_{L}$ has
  dimensions $(h_{L},\bar{h}_{L})$, weight $\Delta_{L}=h_{L}+\bar{h}_{L}<2$ and spin $s_{L}=h_{L}-\bar{h}_{L}$. Depending on the theory, there might be several operators contributing at the lowest order. In this case, we will have to sum over all possible operators.
We simplify our notation by keeping the sum over all possible operators implicit. Note that the factors coming from spins of $V_\alpha$ and $V_\beta$ becomes
% \begin{align}
%  \prod_{k=1}^{n-1} e^{-4\pi i k (s_\alpha-s_\beta)/n}=e^{-2\pi i (n-1)(s_\alpha-s_\beta)}
% \end{align}
\begin{align}
 \prod_{k=1}^{n-1} e^{-4\pi i k (s_\Psi-s_\Phi)/n}=e^{-2\pi i (n-1)(s_\Psi-s_\Phi)}
\end{align}
which precisely cancels out the that of equation (\ref{fncorrspin}).
Therefore, the analytic correlator we need to compute to find the relative entropy is
  \begin{align}
    \label{eq:F_n}\mathcal{F}_{n} & =\frac{\lb\,1+(2\sin(\pi x/n))^{2\Delta_L}f(n)+\cdots\,\rb}{\lb D_n(x)^{n-1}\left[\,1+(2\sin(\pi x/n))^{2\Delta_L}g(n)+\cdots\,\right] \rb}\\
    \label{Dn}
    D_{n}(x)&=\left(\frac{n\sin(\pi x/n)}{\sin(\pi x)}\right)^{2(\Delta_\Psi-\Delta_\Phi)}
    ,\qquad D_{1}(x)=1,
  \end{align}
  where the leading correction $f(n)$ and $g(n)$ are given by the sums in (\ref{fnexpression1}) and (\ref{gnw1}), respectively, over OPE coefficients and light operator correlators on the replica manifold. 
  First, observe that $D_n(x)^{n-1}$ term does not contribute to relative entropy:
  \begin{align}
  &-\p_n\lb (n-1)\log D_n(x)\rb_{n\to 1}=0,
  \end{align}
Hence, we drop it from the denominator, and
  \begin{align}
    % S(\rho_\alpha\|\rho_\beta)\to
    S(\rho_\Psi  \| \rho_\Phi)&=\p_n\left[\log \frac{1+\xi^{2\Delta_L} f(n)+\cdots}{1+\xi^{2\Delta_L} g(n)+\cdots}\right]_{n=1}=\xi^{2\Delta_L}\p_n\left[f(n)-g(n)\right]_{n=1}+\cdots\nn\\
    \xi&=\frac{2\pi x}{n}\ll 1.
  \end{align}
where
  % \begin{align}\label{fnexpression1}
  %   f(n) & =e^{i\pi s_L} \frac{C_{\alpha}^{2}}{2}\sum_{k\neq k'=0}^{n-1}e^{2\pi i s_L (k+k')/n}\langle \mO_{L}(e^{2\pi i k/n})\mO_{L}(e^{2\pi i k'/n})\rangle \nn    \\
  %   &=\frac{e^{i\pi s_L} C_\alpha^2}{2}\sum_{k\neq k'=0}^{n-1}(2i\sin(\pi(k-k')/n))^{-2h_L}(-2i\sin(\pi(k-k')/n))^{-2\bar{h}_L}\nn\\
  %        & =\frac{C_{\alpha}^{2}}{2}\sum_{k\neq k'=0}^{n-1}\frac{1}{(2\sin(\pi |k-k'|/n))^{2\Delta_L}}\ .
  % \end{align}
  \begin{align}\label{fnexpression1}
    f(n) & =e^{i\pi s_L} \frac{C_{\Psi}^{2}}{2}\sum_{k\neq k'=0}^{n-1}e^{2\pi i s_L (k+k')/n}\langle \mO_{L}(e^{2\pi i k/n})\mO_{L}(e^{2\pi i k'/n})\rangle \nn    \\
    &=\frac{e^{i\pi s_L} C_{\Psi}^2}{2}\sum_{k\neq k'=0}^{n-1}e^{-i\pi s_L}(2i\sin(\pi(k-k')/n))^{-2h_L}(-2i\sin(\pi(k-k')/n))^{-2\bar{h}_L}\nn\\
         & =\frac{C_{\Psi}^2}{2}\sum_{k\neq k'=0}^{n-1}\frac{1}{(2\sin(\pi |k-k'|/n))^{2\Delta_L}}\ .
  \end{align}
  % Rewrite the sum above by fixing the separation $m:=k-k' \ (\mathrm{mod}\ n)$. For
  % each $m\in\{1,2,\dots,n-1\}$, set $k'=k-m$ (mod $n$). Then
  % \[
  %   e^{-2\pi i s_L (k+k')/n}= e^{-2\pi i s_L (2k-m)/n}= e^{+2\pi i s_L m/n}\,e^{-4\pi
  %   i s_L k/n},
  % \]
  % and the denominator depends only on $m$. Hence the sum factorizes:
  % \begin{equation}
  %   \sum_{m=1}^{n-1}\frac{e^{2\pi i s_L m/n}}{\bigl(2\sin\frac{\pi m}{n}\bigr)^{2\Delta_L}}
  %   \underbrace{\sum_{k=0}^{n-1} e^{-4\pi i s_L k/n}}_{=:G(n)}.
  % \end{equation}
  % The inner sum $G(n)$ is a geometric series with ratio $q=e^{-4\pi i s_L/n}$:
  % \[
  %   G(n)=\sum_{k=0}^{n-1}q^{k}=
  %   \begin{cases}
  %     n, & q=1,     \\
  %     0, & q\neq 1.
  %   \end{cases}
  % \]
  % Therefore $S(n)=0$ for integer $n$ unless $q=1$, i.e.
  % \[
  %   e^{-4\pi i s_L/n}=1 \quad\Longleftrightarrow\quad \frac{4s_{L}}{n}\in\mathbb{Z}
  %   .
  % \]
  % Now, since $4s_{L}/n$ is not an integer, the numerator $e^{2\pi i s_L m/n}$ is
  % at most a sign. For $s_L=0$, we are left with
  The analytic continuation
  \begin{align}
     & S_{2}(n,\Delta_L)=\frac{n}{2}\sum_{k=1}^{n-1}(2\sin(\pi k/n))^{-2\Delta_L},\qquad S_{2}(1,\Delta_L)=0,
  \end{align}
  was discussed in
  \cite{calabrese2011entanglement}\footnote{Note that our notation has an extra factor $2^{-2\Delta_L}$.}
  \begin{align}
    S'_{2}(1,\Delta)=\frac{2^{-2\Delta}\sqrt{\pi}\Gamma(\Delta+1)}{4\Gamma(\Delta+3/2)}\ .
  \end{align}
  Therefore,
  \begin{align}
    f'(1)=C^2_\Psi \frac{2^{-2\Delta}\sqrt{\pi}\Gamma(\Delta+1)}{4\Gamma(\Delta+3/2)}
  \end{align}
  In a general CFT, there could be several inequivalent choices of the lightest primary
  operator $\mO_{L}$ that contribute at the same order. In this case, the leading
  correction $f(n)$ would be given by the sum over all such operators, that
  following \cite{calabrese2011entanglement} we denote by
  \begin{align}
    f(n) & = C_\Psi^2 S_{2}(n,\Delta) 
  \end{align}
  % \begin{align*}
  %     \JM{f(n) = C_\alpha^2 S_2(n)}
  % \end{align*}
The denominator in $\mathcal{F}_n$ depends on how $\Delta_L$ and $\Delta_{L''}$ compare. We first consider the case, $\Delta_L=\Delta_{L''}$. Then,
  % \begin{align}\label{gnw1}
  %   g(n)= & C_{\alpha}C_\beta\sum_{k=1}^{n-1}\frac{1}{(2\sin(\pi k/n))^{2\Delta_L}}+\frac{C_\beta^2}{2}\sum_{k\neq k'=1}^{n-1}\frac{1}{(2\sin(\pi |k-k'|/n))^{2\Delta_L}}, \nn\\% \nn\\
  %   =&(C_\alpha-C_\beta)C_\beta\sum_{k=1}^{n-1}\frac{1}{(2\sin(\pi k/n))^{2\Delta_L}}+C_\beta^2 S_2(n)
  %   % &=C_\alpha C_\omega \frac{2 S_2(n)}{n}+\sum_{k\neq k'=1}^{n-1}
  %   % \frac{C_{w_k}C_{w_{k'}}}{(2\sin(\pi |k-k'|/n))^{4h_L}}
  % \end{align}
  \begin{align}\label{gnw1}
    g(n)= & C_{\Psi}C_\Phi\sum_{k=1}^{n-1}\frac{1}{(2\sin(\pi k/n))^{2\Delta_L}}+\frac{C_\Phi^2}{2}\sum_{k\neq k'=1}^{n-1}\frac{1}{(2\sin(\pi |k-k'|/n))^{2\Delta_{L}}}, \nn\\% \nn\\
    =&(C_\Psi -C_\Phi)C_\Phi\sum_{k=1}^{n-1}\frac{1}{(2\sin(\pi k/n))^{2\Delta_L}}+C_\Phi^2 S_2(n,\Delta_{L})
    % &=C_\alpha C_\omega \frac{2 S_2(n)}{n}+\sum_{k\neq k'=1}^{n-1}
    % \frac{C_{w_k}C_{w_{k'}}}{(2\sin(\pi |k-k'|/n))^{4h_L}}
  \end{align}
Therefore, expanding in small $x$ we find
% \begin{align}
%   &\log\mathcal{F}_n(x)=(2\pi x/n)^{2\Delta_L}\lb(C_\alpha^2-C_\beta^2)S_2(n)-(C_\alpha-C_\beta)C_\beta \frac{2S_2(n)}{n}\rb+\cdot
% \end{align}
\begin{align}
  &\log\mathcal{F}_n(x)=(2\pi x/n)^{2\Delta_L}\lb(C_\Psi^2-C_\Phi^2)S_2(n)-(C_\Psi-C_\Phi)C_\Phi \frac{2S_2(n,\Delta_L)}{n}\rb+\cdot
\end{align}
Since $S_2(1,\Delta)=0$ we find
% \begin{align}
% S(\rho_\alpha\|\rho_\beta)=&(C_\alpha-C_\beta)^2 S'_2(1)(2\pi x)^{2\Delta_L}+\cdots \nn\\
% &=(C_\alpha-C_\beta)^2 \frac{\sqrt{\pi}\Gamma(\Delta_{L}+1)}{4\Gamma(\Delta_{L}+3/2)}(\pi x)^{2\Delta_L}+\cdots
% \end{align}
\begin{align}
% \JM{
S(\rho_\Psi\|\rho_\Phi)
% }
=&
% \JM{
(C_\Psi-C_\Phi)^2 S'_2(1,\Delta_L)(2\pi x)^{2\Delta_L}+\cdots \nn\\
&=(C_\Psi-C_\Phi)^2\frac{\sqrt{\pi}\Gamma(\Delta_{L}+1)}{4\Gamma(\Delta_{L}+3/2)}(\pi x)^{2\Delta_L}+\cdots
\end{align}

Next, consider the case $\Delta_{L''}<\Delta_L$. In this case, only the denominator of $\mathcal{F}_n$ contributes to the lowest order, and we have
  \begin{align}\label{gnw1p}
    g(n)= & \frac{C_\Phi^2}{2}\sum_{k\neq k'=1}^{n-1}\frac{1}{(2\sin(\pi |k-k'|/n))^{2\Delta_{L''}}}, \nn\\% \nn\\
    =&-C_\Phi^2\sum_{k=1}^{n-1}\frac{1}{(2\sin(\pi k/n))^{2\Delta_{L''}}}+C_\Phi^2 S_2(n,\Delta_{L''})
    % &=C_\alpha C_\omega \frac{2 S_2(n)}{n}+\sum_{k\neq k'=1}^{n-1}
    % \frac{C_{w_k}C_{w_{k'}}}{(2\sin(\pi |k-k'|/n))^{4h_L}}
  \end{align}
and since we assumed $\Delta_{L''}$
\begin{align}
  \log\mathcal{F}_n(x)&=(2\pi x/n)^{2\Delta_{L''}}\lb -C_\Phi^2 S_2(n,\Delta_{L''})+C_\Phi^2 \frac{2S_2(n,\Delta_{L''})}{n}\rb+\cdot\nn\\
% \JM{
S(\rho_\Psi\|\rho_\Phi)
% }
=&
% \JM{
C_\Phi^2 S'_2(1,\Delta_{L''})(2\pi x)^{2\Delta_{L''}}+\cdots \nn\\
&=C_\Phi^2\frac{\sqrt{\pi}\Gamma(\Delta_{L''}+1)}{4\Gamma(\Delta_{L''}+3/2)}(\pi x)^{2\Delta_{L''}}+\cdots
\end{align}

Finally, consider the case $\Delta_{L''}>\Delta_L$, then only the numerator contributes to the lowest order to $\mathcal{F}_n$, and we have
\begin{align}
% \JM{
S(\rho_\Psi\|\rho_\Phi)
% }
=&
% \JM{
C_\Psi^2 S'_2(1,\Delta_{L})(2\pi x)^{2\Delta_{L}}+\cdots \nn\\
&=C_\Psi^2\frac{\sqrt{\pi}\Gamma(\Delta_{L}+1)}{4\Gamma(\Delta_{L}+3/2)}(\pi x)^{2\Delta_{L}}+\cdots
\end{align}
The lesson to learn here is that when $\Delta_{L''}<\Delta_{L}$, to the lowest order the calculation is the same as $S(\rho_0\|\rho_\Phi)$, and when $\Delta_{L''}>\Delta_L$, to the lowest order, it is the same as $S(\rho_\Psi\|\rho_0)$, with $\rho_0$ the vacuum reduced state. As a result, in the remainder of this draft, with no loss of generality, we focus on the case $\Delta_L=\Delta_{L''}$ to simplify our notation. 

Stress tensors $T$ and $\bar{T}$ contribute to the OPE expansion with weights $(2,0)$ and $(0,2)$. In theories with conserved currents $J$ and $\bar{J}$ (e.g. Kac-Moody algebras), these currents have weights $(1,0)$ and $(0,1)$. Therefore, to the lowest order in small $x$ expansion, in any 2d CFT with a relevant deformation $\Delta_L<1$, we can ignore these corrections. 

 % \subsection{Relative entropy $S(\rho_\alpha\|\rho_\beta)$ in the limit $x\to 0^+$}

  In the small interval limit, $x\to 0^+$, the relative entropy of primary states of dimension 
  % $(h_\alpha,\bar{h}_\alpha)$ and $(h_\beta,\bar{h}_\beta)$
  $(h_\Psi,\bar{h}_\Psi)$ and $(h_\Phi,\bar{h}_\Phi)$
  is universal and given by \ref{Ssmallx}
  % \begin{align}
  % S(\rho_\alpha\|\rho_\beta)=\begin{cases}
  % \frac{\sqrt{\pi}\Gamma(\Delta+1)}{4\Gamma(\Delta+3/2)}(C^L_{\alpha\alpha}-C^L_{\beta\beta})^2 (\pi x)^{2\Delta_L}+\cdots&\text{if }\Delta_L\leq 2\\
  % \frac{16}{5c}(h_\alpha-h_\beta)^2(\pi x)^4+\cdots &\text{if }\Delta_L>2
  % \end{cases}
  % \end{align}
  % \begin{align*}
  % S(\rho_\Psi\|\rho_\Phi)=\begin{cases}
  % \frac{\sqrt{\pi}\Gamma(\Delta+1)}{4\Gamma(\Delta+3/2)}(C^L_{\Psi\Psi^*}-C^L_{\Phi\Phi^*})^2 (\pi x)^{2\Delta_L}+\cdots&\text{if }\Delta_L< 2\\
  % \frac{16}{5c}(h_\alpha-h_\beta)^2(\pi x)^4+\cdots &\text{if }\Delta_L>2
  % \end{cases}
  % \end{align*}
  % \NL{Look into the second line of the result above}
  where $\Delta_L$ is the scaling dimension of the lightest spinless primary in the spectrum that appears in the OPE

If we consider the case $\Delta_L>2$, then we will have the contribution from $T$ and $\bar T$ in the OPE which has scaling dimension 2. So we need the following OPEs,
\begin{align}
\frac{\Psi(z_{k}^{+})\Psi^*(z_{k}^{-})}
     {\braket{\Psi(z^+)\Psi^*(z^-)}}
&=
e^{-4\pi i k s_{\Psi}/n}
\left(
\frac{\sin(\pi x)}{\sin(\pi x/n)}
\right)^{2\Delta_{\Psi}}
\nonumber\\
&\quad\times
\Bigg(
1
+\frac{2h_\Psi}{c}(2\sin(\pi x/n))^{2}
 e^{2\pi i(2k/n+1/2)}
 T(e^{2\pi i k/n})
\nonumber\\
&\qquad \qquad
+\frac{2\bar h_\Psi}{c}(2\sin(\pi x/n))^{2}
 e^{-2\pi i(2k/n+1/2)}
 \bar T(e^{2\pi i k/n})
+\cdots
\Bigg),
\\[1ex]
\frac{\Phi(z_{k}^{+})\Phi^*(z_{k}^{-})}
     {\braket{\Phi(z^+)\Phi^*(z^-)}}
&=
e^{-4\pi i k s_{\Phi}/n}
\left(
\frac{\sin(\pi x)}{\sin(\pi x/n)}
\right)^{2\Delta_{\Phi}}
\nonumber\\
&\quad\times
\Bigg(
1
+\frac{2h_\Phi}{c}(2\sin(\pi x/n))^{2}
 e^{2\pi i(2k/n+1/2)}
 T(e^{2\pi i k/n})
\nonumber\\
&\qquad \qquad
+\frac{2\bar h_\Phi}{c}(2\sin(\pi x/n))^{2}
 e^{-2\pi i(2k/n+1/2)}
 \bar T(e^{2\pi i k/n})
+\cdots
\Bigg).
\end{align}
The identity contribution from the two point function of stress tensor is given by,
\begin{align}
    \langle T (z_{k'})T(z_k) \rangle &= e^{-4\pi i ((k'+k)/n + 1/2)}\frac{c}{2(2\sin(\pi (k'-k)/n))^4} \\
    \langle \bar T (\bar z_{k'})\bar T(\bar z_k) \rangle &= e^{4\pi i ((k'+k)/n + 1/2)}\frac{c}{2(2\sin(\pi (k'-k)/n))^4} 
\end{align}
It is clear that the phase term from the OPE gets cancelled from the phase in the two point function. Thus, we can now just substitute $C_\Psi = \frac{2h_\Psi}{c}$ and $C_\Phi = \frac{2h_\Phi}{c}$ in the final answer, with a factor of $c/2$ coming from the two point function of stress tensor. By adding the contribution from both the holomorphic and anti-holomorphic part, we get,
\begin{align}
    S(\rho_\Psi \| \rho_\Phi) &= (\pi x)^4\frac{c}{2} \frac{\sqrt{\pi}\Gamma(3)}{4\Gamma(7/2)}\left[\lb \frac{2h_\Psi}{c}- \frac{2h_\Phi}{c} \rb ^2 + \lb \frac{2\bar h_\Psi}{c}- \frac{2\bar h_\Phi}{c} \rb ^2  \right]+\cdots \nn\\
    &= (\pi x)^4\frac{8}{15c}\lb (h_\Psi - h_\Phi)^2 + (\bar h_\Psi - \bar h_\Phi)^2 \rb + \cdots
\end{align}
  In the small interval limit, $x\to 0^+$, the relative entropy of primary states of dimension 
  % $(h_\alpha,\bar{h}_\alpha)$ and $(h_\beta,\bar{h}_\beta)$
  $(h_\Psi,\bar{h}_\Psi)$ and $(h_\Phi,\bar{h}_\Phi)$
  is universal and given by  
%   \begin{align}
% \boxed{  S(\rho_\Psi\|\rho_\Phi)=\begin{cases}
%   \frac{\sqrt{\pi}\Gamma(\Delta+1)}{4\Gamma(\Delta+3/2)}(C^L_{\Psi\Psi^*}-C^L_{\Phi\Phi^*})^2 (\pi x)^{2\Delta_L}+\cdots&\text{if }\Delta_L\leq 2\\
%   \frac{8}{15c}\lb (h_\Psi - h_\Phi)^2 + (\bar h_\Psi - \bar h_\Phi)^2 \rb(\pi x)^4 + \cdots &\text{if }\Delta_L>2
%   \end{cases}}
%   \end{align}
\begin{align}\label{Ssmallx}
\boxed{
S(\rho_\Psi\|\rho_\Phi)
=
\begin{cases}
\begin{aligned}
&\frac{\sqrt{\pi}\Gamma(\Delta_L+1)}
{4\Gamma(\Delta_L+\frac32)}
\sum_L
(C^L_{\Psi\Psi^*}-C^L_{\Phi\Phi^*})^2 
% &\qquad\qquad\qquad \qquad
(\pi x)^{2\Delta_L}
+\cdots
\end{aligned}
&
\text{if }\Delta_L\le 2,
\\[2ex]
\begin{aligned}
&\frac{8}{15c}
\Big[
(h_\Psi-h_\Phi)^2 
% &\qquad\qquad
+(\bar h_\Psi-\bar h_\Phi)^2
\Big]
(\pi x)^4
+\cdots
\end{aligned}
&
\text{if }\Delta_L>2 .
\end{cases}
}
\end{align}
Note that we have included a sum over $L$ because there can be multiple scalar primaries of the same lowest dimension $\Delta_L$.

  If we include the three-point function, it turns out that it is also independent of spin.
The three-point function contribution from the numerator is given by,
\begin{align}
    \frac{(\pi x)^{3\Delta_L}}{6}\sum_{k_1\neq k_2\neq k_3 =0}^{n-1} \frac{C_\Psi^3}{\lb \sin(\pi (k_3 - k_2)/n) \sin (\pi (k_2 - k_1)/n) \sin(\pi (k_3-k_1)/n)\rb^{\Delta_L}}
\end{align}
The total contribution from the denominator consists of two classes of three-point functions. The first arises from inserting the lightest primary in each of the three $\Phi \Phi^*$ OPEs, and therefore is proportional to $C_\Phi^3$. 
  \begin{align}
     \sum_{k_1 \neq k_2 \neq k_3 = 1}^{n-1} C_\Phi^3(\pi x)^{3\Delta_L}& e^{i\pi s_L(2(k_1 + k_2+k_3)/n + 3/2) }\braket{ \mO_L(z_{k_3})\mO_L(z_{k_2})\mO_L(z_{k_1})} \\
            & = C_\Phi^3\frac{(\pi x)^{3\Delta_L}}{6}\sum_{k_1 \neq k_2 \neq k_3= 1}^{n-1}\frac{1}{\lb \sin(\pi (k_3 - k_2)/n) \sin (\pi (k_2 - k_1)/n) \sin(\pi (k_3-k_1)/n)\rb^{\Delta_L}}
  \end{align}
  Another one arises from inserting the lightest primary in two $\Phi \Phi^*$ OPEs and one from $\Psi \Psi^*$ OPE, which is proportional to $C_\Psi C_\Phi^2$. 
    \begin{align}
       \sum_{k_1 \neq k_2 = 1}^{n-1}C_\Psi C_\Phi^2 (\pi x)^{3\Delta_L}& e^{i\pi s_L(2(k_1 + k_2)/n + 3/2) }\braket{ \mO_L(z_{k_2})\mO_L(z_{k_1})\mO_L(z_{0})} \nn \\ 
       & = C_\Psi C_\Phi^2\frac{(\pi x)^{2\Delta_L}}{2}\sum_{k_1 \neq k_2 = 1}^{n-1}\frac{\zeta_n^{\Delta_L}}{\lb \sin(\pi (k_2 - k_1)/n) \sin (\pi k_2/n) \sin(\pi k_1/n)\rb^{\Delta_L}}
  \end{align}
These contributions can be combined in the following form,
\begin{align}
    \frac{(\pi x)^{3\Delta_L}}{6}\sum_{k_1\neq k_2\neq k_3 =0}^{n-1} \frac{\delta_{k_1k_2k_3,0}\lb C_\Psi C_\Phi^2 - C_\Phi^3 \rb  + C_\Phi^3}{\lb \sin(\pi (k_3 - k_2)/n) \sin (\pi (k_2 - k_1)/n) \sin(\pi (k_3-k_1)/n)\rb^{\Delta_L}}
\end{align}
The net contribution to relative entropy from the three-point function would be,
\begin{align}
    S(\rho_\Psi \| \rho_\Phi) &= \frac{\sqrt{\pi}\Gamma(\Delta+1)}{4\Gamma(\Delta+3/2)}(C^L_{\Psi\Psi^*}-C^L_{\Phi\Phi^*})^2 (\pi x)^{2\Delta_L}+(\p_n \mathcal{T}_n )|_{n=1}(\pi x)^{3\Delta_L} +\cdots  \\
    \mathcal{T}_n &= \frac{1}{6}\sum_{k_1\neq k_2 \neq k_3 = 0}^{n-1} \frac{C_\Psi^3 - C_\Phi^3 - \delta_{k_1k_2k_3,0}(C_\Psi C_\Phi^2 - C_\Phi^3)}{\lb \sin(\pi (k_3 - k_2)/n) \sin (\pi (k_2 - k_1)/n) \sin(\pi (k_3-k_1)/n)\rb^{\Delta_L}}
\end{align}

\subsubsection{Some examples}
  \paragraph{Chiral vertex operators}

  In massless free boson theory, chiral coherent states $V_\alpha$ and $V_{i\partial\phi}$ are eigenstates of the Hamiltonian with eigenvalues $(\alpha^2/2,0)$ and $(\beta^2/2,0)$ respectively. The lowest dimension scalar primary field that appears in the OPE of a pair of $V_\alpha$ is $J=i(\p\varphi)$ with dimenion $(1,0)$, and $C^{i\p\varphi}_{\alpha\alpha}=-\alpha$. Therefore,
  \begin{align}
    S(\rho_\alpha\|\rho_\beta)=\frac{(\alpha-\beta)^2(\pi x)^2}{3}+\cdots
  \end{align}
  which matches the expansion of 
  \begin{align}
    S(\rho_\alpha\|\rho_\beta)=(\alpha-\beta)^2(1-\pi x\cot(\pi x))=\frac{(\alpha-\beta)^2(\pi x)^2}{3}+\cdots
  \end{align}

  \paragraph{Non-chiral vertex operators}

  The algebra of free massless boson is a tensor product of the $U(1)$ algebra $U(1)\times U(1)$, with generators $(J,\bar{J})$ and charges $(h_\alpha,\bar{h}_\alpha)$. A general vertex operator is
  \begin{align}
    V_{\alpha,\bar{\alpha}}(z,\bar{z})=(e^{i\alpha \varphi(z)}e^{i\bar{\alpha}\bar{\varphi}(\bar{z})})(z,\bar{z})
  \end{align}
  with dimesnion $(\alpha^2/2,\bar{\alpha}^2/2)$. The $n$-point function of these correlators factor 
  \begin{align}
    \braket{V_{\alpha_1,\bar{\alpha}_1}\cdots V_{\alpha_n,\bar{\alpha}_n}}=\braket{V_{\alpha_1}\cdots V_{\alpha_n}}\braket{V_{\bar{\alpha}_1}\cdots V_{\bar{\alpha}_n}}
  \end{align}
  and as a result, their relative entropies factor
  \begin{align}
  S(\rho_{\alpha,\bar{\alpha}}\|\rho_{\beta,\bar{\beta}})=S(\rho_{\alpha}\|\rho_{\beta})+S(\rho_{\bar{\alpha}}\|\rho_{\bar{\beta}})
  \end{align}
  Therefore, in the small interval expansion, we have
  \begin{align}\label{smallxCoherentstates}
\boxed{S(\rho_{\alpha,\bar{\alpha}}\|\rho_{\beta,\bar{\beta}})=((\alpha-\beta)^2+(\bar{\alpha}-\bar{\beta})^2)(1-\pi x\cot(\pi x))=\frac{((\alpha-\beta)^2+(\bar{\alpha}-\bar{\beta})^2)(\pi x)^2}{3}+\cdots}
  \end{align}
  Spinless vertex operators of massless boson have charge $h_\alpha=\bar{h}_\alpha=\alpha$ and dimension $\Delta=\alpha^2$. Therefore,
  \begin{align}
  S(\rho_{\alpha,\alpha}\|\rho_{\beta,\beta})=(2(\alpha-\beta)^2)(1-\pi x\cot(\pi x))=\frac{2(\alpha-\beta)^2(\pi x)^2}{3}+\cdots
  \end{align}

  \paragraph{Current operators}

  The relative entropy of the primary state created by $J=i\p\varphi$ with respect to a general chiral primary operator is
  \begin{align}
  S(\rho_{i\p\varphi}\|\rho_\beta)=S(\rho_{i\p\varphi}\|\rho_0)+S(\rho_0\|\rho_\beta)
  \end{align}
At small $x$ we have
  \begin{align*}
      S(\rho_0\| \rho_\beta) = \beta^2\lb \frac{(\pi x)^2}{3} + \frac{(\pi x)^4}{45}+\mathcal{O}(x^6)\cdots\rb
  \end{align*}
  and
  \begin{align}
      S(\rho_{i\partial\phi}\|\rho_0)&=2(1-\pi x \cot(\pi x))+2\lb \log(2\sin(\pi x))+\psi_0(\csc(\pi x)/2)+\sin(\pi x)\rb\nn\\
      &\simeq_{x \to 0^+} \frac{8}{15}(\pi x)^4 +\mO(x^6) 
  \end{align}
  where we used the large $z$ expansion of $\psi_0$,
  \begin{align}
      \psi_0(z) &= \log z - \frac{1}{2z} - \frac{1}{12z^2} + \frac{1}{120z^4} + \cdots, \\ 
      z &= \frac{1}{2\sin(\pi x)}
  \end{align}
  Adding the above two gives,
  \begin{align*}
      S(\rho_{i\p \varphi}\| \rho_\beta) \simeq_{x \to 0^+} \beta^2\lb \frac{(\pi x)^2 }{3}+ \frac{5(\pi x)^4}{9}+\mathcal{O}(x^6)+\cdots\rb
  \end{align*}
  % $:(V^\dagger_\alpha V_\beta)(0):$ which is a coherent operator of charge $\alpha-\beta$, and dimension $h_L+\bar{h}_L=(\alpha-\beta)^2$.
  % This relative entropy has been discussed extensively in the small $x$ limit, and for coherent eigenstates $V_\alpha$ and $V_{i\partial\phi}$ states of the free boson CFT. 
  % Here, we first go over the less discussed $x\to 1^-$ limit of the relative entropy, and show that universally relative entropy is expected to diverge as
  % \begin{align}
  %   S(\rho_\alpha\|\rho_\beta)=
  % \end{align}
  

% \subsection{Large interval, \texorpdfstring{$x\to 1^-$ limit}{x -> 1-}}
% \subsection{Large interval, \texorpdfstring{$x\to 1^-$ limit}{x -> 1^- limit}}
\subsection{Large interval, \texorpdfstring{$x\to 1^-$ limit}{Large interval, x to 1 minus limit}}
% Before we start the replica trick, we note that for any pair of reduced states
% \begin{align}
% &S(\psi_A\|\chi_A)=\tr((\psi_A-\chi_A)K_{\chi_A})-\Delta S(A)\nn\\
% &S(\psi_{\bar{A}}\|\chi_{\bar{A}})=\tr\lb(\psi_{\bar{A}}-\chi_{\bar{A}})K_{\chi_\bar{A}}\rb -\Delta S(\bar{A})\nn\\
% &\Delta S(A)=S(\psi_A)-S(\chi_A),\qquad K_\rho\equiv -\log\rho .
% \end{align}
Before we start the replica trick, we note that for any pair of reduced states
\begin{align}
S(\psi_A\|\chi_A)
&=\operatorname{tr}\bigl((\psi_A-\chi_A)K_{\chi_A}\bigr)-\Delta S(A)\nonumber\\
S(\psi_{\bar A}\|\chi_{\bar A})
&=\operatorname{tr}\bigl((\psi_{\bar A}-\chi_{\bar A})K_{\chi_{\bar A}}\bigr)
-\Delta S(\bar A)\nonumber\\
\Delta S(A)
&=S(\psi_A)-S(\chi_A),\qquad K_\rho\equiv -\log\rho .
\end{align}
Since for pure global states we have $S(A)=S(\bar{A})$ we define the global modular Hamiltonian $\hat{K}_A=K_A-K_{\bar{A}}$ and find
\begin{align}\label{Khatformula}
   S(\psi_{\bar{A}}\|\chi_{\bar{A}})-S(\psi_A\|\chi_A)= -\braket{\Psi|\hat{K}_{\chi,A}|\Psi}\ .
\end{align}
Therefore, in the small $1-x$ limit, we should expect a term that is order $\ep^{2\Delta_L}$ coming from $S(\rho_\Psi\|\rho_\Phi)$, and then a correction term due to the modular Hamiltonian $\hat{K}_{\chi,A}$. As we see below, the key observation is that the $\hat{K}$ term diverges as $x\to 1$ with corrections starting at order $1/\ep$. Note that further corrections to $\hat{K}$ can appear at lower order than $\ep^{2\Delta_L}$, as well. 


%\NL{Change this paragraph from the notation $V_\alpha$ to $\Phi$ and $\Psi$}
In the large interval limit, we expand the correlator in (\ref{relEntropyReplica}) in the limit of small $\ep\equiv 1-x$. In this limit, we have to expand in a different OPE channel (at $z_k^+$ and $z_{k+1}^-$) where operators from sheet $k$ are close to operators on sheet $k+1$. The correlators with $n=1$ (normalizations of density matrices) are symmetric under $x\to 1-x$. 
% \begin{align}
%  &\frac{V_{w_k}(z_k^{+})V_{-w_{k+1}}(z_{k+1}^{-})}{\sqrt{\braket{V_{w_k}(z^+)V_{-w_k}(z^-)}\braket{V_{w_{k+1}}(z^+)V_{-w_{k+1}}(z^-)}}}=\sum_p C^p_{\omega_k\omega_{k+1}}\zeta_1^{\Delta_p}\lb \frac{\zeta_n}{\zeta_1}\rb^{\Delta_p-\Delta_{\omega_k}-\Delta_{\omega_{k+1}}}e^{i\pi\Phi_k^{(p)}}\mO_p(e^{i\pi(2k+1)/n})\nn\\
%  &\zeta_n=2\sin(\pi\ep/n),\qquad \Phi^{(p)}_k=s_p(\theta_k-1/2)-(s_{\omega_k}+s_{\omega_{k+1}})\lb \theta_k-1\rb,\qquad \theta_k=\frac{2k+1}{n}
% \end{align}
\begin{align}
 &\frac{\Psi(z_k^{+})\Phi^*(z_{k+1}^{-})}{\sqrt{\braket{\Psi(z^+)\Psi^*(z^-)}\braket{\Phi(z^+)\Phi^*(z^-)}}}=\sum_p C^p_{\Psi \Phi^*}\zeta_1^{\Delta_p}\lb \frac{\zeta_n}{\zeta_1}\rb^{\Delta_p-\Delta_{\Psi}-\Delta_{\Phi}}e^{i\pi\Omega_k^{(p)}}\mO_p(e^{i\pi(2k+1)/n})\nn\\
 &\zeta_n=2\sin(\pi\ep/n),\qquad \Omega^{(p)}_k=s_p(\theta_k-1/2)-(s_{\Psi}+s_{\Phi})\lb \theta_k-1\rb,\qquad \theta_k=\frac{2k+1}{n}
\end{align}
For the $n$-dependent correlators, at small $1-x$ the OPE expansion of the numerator remains the same as small $x$ with $x\to 1-x$. However, the denominator changes because for $k=0$ and $k=n-1$ we now have the OPE of $\Psi$ and $\Phi^*$ and $\Phi$ and $\Psi^*$ respectively.  For $h_\Psi\neq h_\Phi$, there cannot be an identity contribution in the OPE of $\Psi$ and $\Phi$. Therefore, in the denominator, to the lowest order, we need the following OPEs 
 %  \begin{align}
 %    % \frac{V_\beta(z_k^+)V_{-\beta}(z_{k+1}^{-})}{\braket{V_\beta(z^+)V_{-\beta}(z^-)}}&=\lb\frac{\zeta_1}{\zeta_n}\rb^{2\Delta_{\beta}}  e^{-2i\pi s_\beta(\theta_k-1)}\lb 1+ C_\beta\zeta_n^{\Delta_L} e^{i\pi s_L(\theta_k-1/2)}\mO_L(e^{i\pi\theta_k})\cdots\rb\nn\\
 %    \frac{\Phi^*(z_{k+1}^{-})\Phi(z_k^+)}{\braket{\Phi(z^+)\Phi^*(z^-)}}&=\lb\frac{\zeta_1}{\zeta_n}\rb^{2\Delta_{\Phi}}  e^{-2i\pi s_\Phi(\theta_k)}\lb 1+ C_\Phi\zeta_n^{\Delta_L} e^{i\pi s_L(\theta_k+1/2)}\mO_L(z_k^+)\cdots\rb\nn\\
 %   % \frac{V_{\alpha}(z_0^{+})V_{-\beta}(z_{1}^{-})}{\sqrt{\braket{V_{\alpha}(z^+)V_{-\alpha}(z^-)}\braket{V_{\beta}(z^+)V_{-\beta}(z^-)}}}&=\lb\frac{\zeta_1}{\zeta_n}\rb^{\Delta_\alpha+\Delta_{\beta}}\zeta_n^{\Delta_{L'}} C^{L'}_{\alpha\beta} e^{i\pi(s_L(1/n-1/2)-(s_\alpha+s_\beta)(1/n-1))}\mO_{L'}(e^{i\pi /n})+\cdots\nn\\
 %    \frac{\Phi^*(z_{1}^{-})\Psi(z_0^{+})}{\sqrt{\braket{\Psi(z^+)\Psi^*(z^-)}\braket{\Phi(z^+)\Phi^*(z^-)}}}&=\lb\frac{\zeta_1}{\zeta_n}\rb^{\Delta_\Psi + \Delta_\Phi} \zeta_n^{\Delta_{L'}}C_{\Psi\Phi}e^{i\pi(s_{L'} - s_\Psi - s_\Phi)(1/n+1/2)}e^{i\pi (s_\Psi + s_\Phi)/2}\mO_{L'}(z_0^+) + \cdots\nn \\
 %    % \frac{V_{\beta}(z_{n-1}^{+})V_{-\alpha}(z_0^{-})}{\sqrt{\braket{V_{\alpha}(z^+)V_{-\alpha}(z^-)}\braket{V_{\beta}(z^+)V_{-\beta}(z^-)}}}&=\lb\frac{\zeta_1}{\zeta_n}\rb^{\Delta_\beta+\Delta_{\alpha}}\zeta_n^{\Delta_{L'}} C^{L'}_{\alpha\beta} e^{ i\pi(s_L(3/2-1/n)-(s_\alpha+s_\beta)(1-1/n))}\mO_{L'}(e^{-i\pi /n})+\cdots \nn\\
 %    \frac{\Phi(z_{n-1}^{+})\Psi^*(z_0^{-})}{\sqrt{\braket{\Psi(z^+)\Psi^*(z^-)}\braket{\Phi(z^+)\Phi^*(z^-)}}}&= \lb\frac{\zeta_1}{\zeta_n}\rb^{\Delta_\Psi + \Delta_\Phi} \zeta_n^{\Delta_{L'}}C_{\Psi\Phi}e^{i\pi(s_{L'} -s_\Psi - s_\Phi)(-1/n + 5/2)}e^{i \pi (s_\Psi + s_\Phi)/2}\mO_{L'}(z_0^-) + \cdots
 % %   \delta s&=s_L-s_{\alpha}-s_{\beta}
 %  \end{align}
  \begin{align}\label{eq: general OPE with identity}
      \frac{\Phi^*(z_{k+1}^{-})\Phi(z_k^+)}{\braket{\Phi(z^+)\Phi^*(z^-)}}&=\lb\frac{\zeta_1}{\zeta_n}\rb^{2\Delta_{\Phi}}  e^{-2i\pi s_\Phi(\theta_k)}\lb 1+ \sum_j C^{j}_{\Phi}\zeta_n^{\Delta_j} e^{i\pi s_j(\theta_k+1/2)}\mO_{j}(e^{i\pi \theta_k})\cdots\rb\nn\\
      \frac{\Phi^*(z_{1}^{-})\Psi(z_0^{+})}{\sqrt{\braket{\Psi(z^+)\Psi^*(z^-)}\braket{\Phi(z^+)\Phi^*(z^-)}}}&=\lb\frac{\zeta_1}{\zeta_n}\rb^{\Delta_\Psi + \Delta_\Phi} \big( \delta_{\Delta_{\Phi}\Delta_{\Psi}} e^{-2\pi i s_\Phi/n}1 +   \nn \\ &\qquad \sum_j C^j_{\Phi\Psi}\zeta_n^{\Delta'_j}e^{i\pi(s'_j - s_\Psi - s_\Phi)(1/n+1/2)}e^{i\pi (s_\Psi + s_\Phi)/2}\mO'_j(e^{i\pi/n}) + \cdots\big)\nn \\
      \frac{\Phi(z_{n-1}^{+})\Psi^*(z_0^{-})}{\sqrt{\braket{\Psi(z^+)\Psi^*(z^-)}\braket{\Phi(z^+)\Phi^*(z^-)}}}&= \lb\frac{\zeta_1}{\zeta_n}\rb^{\Delta_\Psi + \Delta_\Phi} \big( \delta_{\Delta_{\Phi} \Delta_{\Psi}} e^{-2\pi i s_{\Phi}(-1/n +2)}1 + \nn\\
      &\qquad \sum_j C^j_{\Psi\Phi}\zeta_n^{\Delta'_j}e^{i\pi(s'_j -s_\Psi - s_\Phi)(-1/n + 5/2)}e^{i \pi (s_\Psi + s_\Phi)/2}\mO'_j(e^{i(2n-1)\pi/n}) + \cdots\big)
  \end{align}
  where $\mO_j$ is the $j^{th}$ lowest dimension operator in the OPE $\Psi\Psi^*$ and  $\Phi\Phi^*$,\footnote{As before, we continue to assume, without loss of generality, that the primaries in the OPEs of $\Psi \Psi^*$ and $\Phi \Phi^*$ have the same dimension.} and $\mO'_j$ is the $j^{th}$ lowest dimension primary in the OPE $\Psi\Phi^*$ (and similarly $(\mO'_j)^*$ is the $j^{th}$ lowest dimension primary in the OPE of $\Psi^*\Phi$) which need not be the same as $\mO_L$ which is the lowest dimension operator in the OPE of $\Phi$ and $\Phi^*$. To match the notation from previous section, we use $\Delta_L\equiv \Delta_1$, $C_\Phi=C_\Phi^{(1)}$, $C_{\Psi\Phi}=C^{(1)}_{\Psi\Phi}$. 

During computation of the contribution to $\mathcal{F}_n$, we want to keep only those terms which are invariant under a rotation of $\pi$ about its midpoint. The numerator would contribute only if $s_\Psi \in \mathbb{Z}$. This implies that the denominator will contribute only if $s_\Phi \in \mathbb{Z}$, since we would require $s_\Psi + s_\Phi \in \mathbb{Z}$, and hence $s_\Phi \in \mathbb{Z}$.


In the numerator of $\mathcal{F}_n$, the phases due to the identity piece is given by:
% \begin{align}
%   \prod_{k=0}^{n-1}e^{-2\pi i s_\alpha(\theta_k-1)}=1
% \end{align}
% or should we take the product over $z_{k+1}^-- z_{k}^+$, instead of $z_k^+-z_{k+1}^-$, which would give,
\begin{align}
  \prod_{k=0}^{n-1}e^{-2\pi i s_\Psi(\theta_k)}=e^{-2\pi i ns_\Psi}
\end{align}
and to the lowest order $\zeta_n^{2\Delta_L}$ we have 
% \begin{align}
%   &\lb \frac{\zeta_1}{\zeta_n}\rb^{2n\Delta_\alpha}\lb 1+\zeta_n^{2\Delta_L}\frac{C_\alpha^2}{2}\sum_{k\neq k'}e^{i\pi s_L(\theta_k+\theta_{k'}-1)}\braket{\mO_L(e^{i\pi\theta_k})\mO_L(e^{i\pi\theta_{k'}})}+\cdots \rb\nn\\
%   &=\lb \frac{\zeta_1}{\zeta_n}\rb^{2n\Delta_\alpha}\lb 1+\zeta_n^{2\Delta_L}\frac{C_\alpha^2}{2}\sum_{k\neq k'}\frac{1}{(2\sin(\pi|k-k'|/n))^{2\Delta_L}}+\cdots \rb\nn\\
%   &=\lb \frac{\zeta_1}{\zeta_n}\rb^{2n\Delta_\alpha}\lb 1+\zeta_n^{2\Delta_L}C_\alpha^2 S_2(n)+\cdots \rb
% \end{align}
\begin{align}
  &\lb \frac{\zeta_1}{\zeta_n}\rb^{2n\Delta_\Psi}\lb 1+\zeta_n^{2\Delta_L}\frac{C_\Psi^2}{2}\sum_{k\neq k'}e^{i\pi s_L(\theta_k+\theta_{k'}+1)}\braket{\mO_L(e^{i\pi\theta_k})\mO_L(e^{i\pi\theta_{k'}})}+\cdots \rb\nn\\
  &=\lb \frac{\zeta_1}{\zeta_n}\rb^{2n{\Delta_\Psi}}\lb 1+\zeta_n^{2\Delta_L}\frac{C^2_{\Psi}}{2}\sum_{k\neq k'}\frac{1}{(2\sin(\pi|k-k'|/n))^{2\Delta_L}}+\cdots \rb\nn\\
  &=\lb \frac{\zeta_1}{\zeta_n}\rb^{2n{\Delta_\Psi}}\lb 1+\zeta_n^{2\Delta_L}C^2_{{\Psi}} S_2(n)+\cdots \rb
\end{align}

To the lowest order, in the denominator, we have $n-2$ copies of the OPE of $\Phi$ and $\Phi^*$ and the OPEs of $\Psi\Phi^*$ and $\Psi^*\Phi$. 
% Once again, the phases that depend on $s_\Psi$ and $s_\Phi$ cancel \JM{(phases from the identity piece cancel), i.e.,
The phase from the identity piece in the denominator is given by,
% \begin{align*}
%   \prod_{k=1}^{n-2}e^{-2\pi i s_\beta(\theta_k-1)}=1
% \end{align*}
% (however, if we consider again $z_{k+1}^- - z_k^+$, we get,
\begin{align}
  \prod_{k=1}^{n-2}e^{-2\pi i s_\Phi\theta_k}=e^{-2\pi i s_\Phi(n-2)}
\end{align}
and to the lowest order we have
% \begin{align}
%   &\lb \frac{\zeta_1}{\zeta_n}\rb^{2((n-1)\Delta_\beta+\Delta_\alpha)}\zeta_n^{2\Delta_{L'}}C_{\alpha\beta}^2 e^{i\pi s_{L'}}\braket{\mO_{L'}(e^{i\pi/n})\mO_{L'}(e^{-i\pi/n})}\nn\\
%   &=\lb \frac{\zeta_1}{\zeta_n}\rb^{2((n-1)\Delta_\beta+\Delta_\alpha)}\zeta_n^{2\Delta_{L'}}C_{\alpha\beta}^2(2\sin(\pi/n))^{-2\Delta_{L'}}
% \end{align}
\begin{align}
  &\lb \frac{\zeta_1}{\zeta_n}\rb^{2((n-1)\Delta_\Phi+\Delta_\Psi)}\zeta_n^{2\Delta_{L'}}C_{\Psi\Phi}^2 e^{-2\pi i (n-2)s_\Phi} e^{i\pi s_{L'}}e^{2 \pi i (s_{L'} - s_\Psi - s_\Phi)}\braket{\mO_{L'}(1)\mO_{L'}(e^{\pi/n})}\nn\\
  &=\lb \frac{\zeta_1}{\zeta_n}\rb^{2((n-1)\Delta_\Phi+\Delta_\Psi)}C_{\Psi\Phi}^2e^{-2\pi i s_\Phi(n-2)}e^{-2 \pi i( s_\Psi + s_\Phi)}\xi_n^{2\Delta_{L'}}\nn\\
  \xi_n&=\frac{\sin(\pi \ep/n)}{\sin(\pi/n)}
\end{align}
The $e^{2\pi i s_{L'}}$ gets cancelled from the phase of the two-point function.
   Therefore, the new expression for $\mathcal{F}_n$ is
  % Since the overall factors $(2\sin(\pi \ep/n))$ come with powers that are additive in the operator dimensions, the new expression for $\mathcal{F}_n$  takes the form
\begin{align}
    \label{eq:F_n_1_minus_x}
    % \mathcal{F}_{n}&=\frac{\lb \,1+\zeta_n^{2\Delta_L}C_\alpha^2 S_2(n)+\cdots\,\rb}{\lb (D_n(\ep)e^{2\pi i (s_\alpha-s_\beta)})^{(n-1)}\zeta_n^{2\Delta_{L'}}C_{\alpha\beta}^2(2\sin(\pi/n))^{-2\Delta_{L'}}+\cdots \rb}.\\
    \mathcal{F}_{n}&=\frac{e^{-2\pi i n s_\Psi}\lb \,1+\zeta_n^{2\Delta_L}C_\Psi^2 S_2(n)+\cdots\,\rb}{\lb \lb D_n(\ep)e^{-2\pi i (s_\Psi-s_\Phi)} \rb^{(n-1)}e^{-2\pi i (n-2)s_\Phi}e^{-2 \pi i( s_\Psi + s_\Phi)}\xi_n^{2\Delta_{L'}}(C^{L'}_{\Psi\Phi})^2 +\cdots\rb}.
    %l(\ep)&=\sum_{ijk}C_\Phi^{(k)}\lb \frac{C_{\Psi\Phi}^{(i)}C_{\Psi\Phi}^{(j)}}{(C^{(1)}_{\Psi\Phi})^2}\rb 
\end{align}
At this point, the leading term in the OPE of $\Psi\Phi^*$ results in a divergent relative entropy because $\p_n\xi_n|_{n\to 1}=\infty$. The origin of this divergence can be traced back to the fact that the two-point function of $\mO_{L'}\mO_{L'}$ is evaluated at distance $(2\sin(\pi/n))$ which vanishes at $n\to 1$. In fact, the limit $\ep\to 0$ and $n\to 1$ do not commute here. Note that the corrections to the denominator can come from summing over $L'$, namely the 2-point function of $\mO_{L'}\mO_{L'}$ coming from two pairs of OPEs of $\Phi\Psi^*$, and the three-point function coming from $\mO_{L'}\mO_{L''}\mO_L$, where $\mO_{L'}$ and $\mO_{L''}$ come from the OPEs of $\Phi\Psi^*$, and $\mO_L$ comes from the OPE of $\Phi\Phi$. There could also be higher point functions. Each correction from the two-point function also carries the same divergence! Going back to the same $x$ expansion of relative entropy in (\ref{fnexpression1}), we note that if we only pick one of the terms with fixed $k\neq k'$, the analytic continuation gives the same divergence.

Let us explain this point more explicitly. The difference between the numerator and the denominator of $\mathcal{F}_n$ is the switching of one copy $\Phi\Phi^*$ to a copy of $\Psi\Phi^*$ in the 2n-point function. Define the formal primary operator $(\Psi\Phi^{-1})$ with the property:
\begin{align}
    \lim_{\delta\to 0}\delta^{\Delta_{(\Psi\Phi^{-1})}}(\Psi\Phi^{-1})(\delta)\Phi(0)=\Psi(0)
\end{align}
so that in the denominator, we have the correlator
\begin{align}
  \lim_{\delta\to 0}  \braket{(\Psi\Phi^{-1})(z_0^+-\delta)(\Psi\Phi^{-1})^*(z_0^-+\delta)\prod_{k=0}^{n-1}\Phi(z_k^+)\Phi^*(z_k^-)}
\end{align}
Note that in general, such a primary might not exist. The intuition originates from the example of coherent states of $U(1)$ current algebra, the primary   $: e^{i(\alpha-\beta)\phi}:$ plays the role of $(e^{i\alpha \phi}e^{-i\beta\psi})$.
We can expand the normalized version of this correlator by taking the OPE of $(\Psi\Phi^{-1})(\Psi\Phi^{-1})^*$ with distance $(2\sin(\pi(1-x)/n)$, and taking pair-wise OPEs of $\Phi\Phi^*$ with distance $(2\sin(\pi \ep/n))$ so that we have
\begin{align}
    \lb \frac{\sin(\pi \ep/n)}{\sin(\pi (1-\ep)/n)}\rb^{-2\Delta_{(\Psi\Phi^{-1})}}\lb 1+\sum_p C^p_{(\Psi\Phi^{-1})(\Psi\Phi^{-1})} C^p_{\Phi\Phi}\lb \frac{\sin(\pi \ep/n)}{\sin(\pi (1-\ep)/n)}\rb^{-2\Delta_p}\sum_{k=0}^n\frac{1}{(2\sin(\pi (k+1/2)/n)^{2\Delta_p}}+\cdots\rb
\end{align}
Note that at small $\ep$ and integer $n$ this is a sensible expansion, and higher order terms are suppressed; however, the small $\ep$ does not commute with the small $n-1$. If we first expand in $\ep$ at integer $n$ we find
\begin{align}
    -2\Delta_{(\Psi\Phi^{-1})}\p_n\lb \frac{\sin(\pi \ep/n)}{\sin(\pi (1-\ep)/n)}\rb_{n\to 1}=2\Delta_{(\Psi\Phi^{-1})}\cot(\pi\ep)
\end{align}
which diverges as $\ep\to 0$. However, keeping $\ep$ small and fixed, and taking the $n\to 1$ limit we find that 
\begin{align}
    \frac{\sin(\pi \ep/n)}{\sin(\pi (1-\ep)/n)}=1-\pi\cot(\pi \ep)(n-1)+\cdots
\end{align}
which means that we are not allowed to truncate the series expansion, as higher order $\Delta_p$ also contribute at the same $1/\ep$ order.


To better understand the relative entropy in the limit $x\to 1^-$ we consider special cases.

\subsubsection{Special cases:}
\paragraph{Relative entropy $S(\rho_\Phi\|\rho_0)$ at $x\to 1^-$:}
%\NL{Change notation to $\Phi$ and $\Psi$}
The replica trick for this relative entropy in the $x\to 1^-$ limit is
% \begin{align}
%   \mathcal{F}_n=\frac{1+C_\beta^2(\sin(\pi \ep/n))^{2\Delta_L}f(n)+\cdots }{D_n(\ep)^{(n-1)}(2\sin(\pi \ep/n))^{4h_\beta}\braket{V_\beta(z_{n-1,n}^+)V_{-\beta}(z_{1,n}^-)}}
% \end{align}
\begin{align}
  \mathcal{F}_n&=\frac{1+C_\Phi^2(2\sin(\pi \ep/n))^{2\Delta_L}f(n)+\cdots }{\lb \frac{n \sin (\pi \ep /n)}{\sin \pi \ep} \rb^{2\Delta_\Phi (n-1)}\frac{e^{2\pi i ns_\Phi}e^{i\pi s_\Phi}}{e^{2\pi i (n-1)s_\Phi}}
  % D_n(\ep)^{(n-1)}
  \lb2\sin(\pi \ep/n)\rb^{2\Delta_\Phi}\braket{\Phi(z_{0,n}^+)\Phi^*(z_{0,n}^-)}} \nn \\
  &= \frac{1+C_\Phi^2(2\sin(\pi \ep/n))^{2\Delta_L}f(n)+\cdots }{D_n(\ep)^{n-1} \lb \frac{\sin(\pi \ep/n)}{\sin(\pi (1-\ep)/n)}\rb^{2\Delta_\Phi}
  % D_n(\ep)^{(n-1)}
  }
\end{align}
Therefore,
% \begin{align}
%   \log\mathcal{F}_n=&(\pi \ep/n)^{2\Delta_L}C_\beta^2 f(n)-(n-1)\log D_n(\ep)-4h_\beta\log\lb\frac{2\sin(\pi \ep/n)}{2\sin(\pi (1+\ep)/n)}\rb+\cdots
% \end{align}
\begin{align}
  \log\mathcal{F}_n=&(2\pi \ep/n)^{2\Delta_L}C_\Phi^2 f(n)-(n-1) 2\Delta_\Phi \log \lb \frac{n \sin (\pi \ep/n )}{\sin (\pi \ep)} \rb -2\Delta_\Phi\log\lb\frac{\sin(\pi \ep/n)}{\sin(\pi (1-\ep)/n)}\rb+\cdots\nn\\
  &=(2\pi \ep/n)^{2\Delta_L}C_\Phi^2 f(n)- 2n\Delta_\Phi \log \lb \frac{n \sin (\pi \ep/n )}{\sin (\pi \ep)} \rb -2\Delta_\Phi\log\lb\frac{\sin(\pi \ep)}{n\sin(\pi (1-\ep)/n)}\rb+\cdots
    \end{align}
We have
% \begin{align}
%   S(\rho_\alpha\|\rho_\beta)&=\p_n\left[\log \frac{1+\zeta^{2\Delta_L} f(n)+\cdots}{\zeta^{2\Delta_{L'}} g_2(n)+\cdots}\right]_{n=1}=\zeta^{2\Delta_L}\p_n\left[f(n)-g_2(n)\right]_{n=1}+\cdots\nn\\
%   &=(C_\alpha-C_{\alpha\beta})^2 S'_2(1)(2\pi (1-x))^{2\Delta_L}+\cdots \nn\\
% &=(C_\alpha-C_{\alpha\beta})^2 \frac{\sqrt{\pi}\Gamma(\Delta_{L}+1)}{4\Gamma(\Delta_{L}+3/2)}(\pi (1-x))^{2\Delta_L}+\cdots
% \end{align}
\begin{align}
  S(\rho_\Phi\|\rho_0)&=(\pi \ep)^{2\Delta_L}C_{\Phi}^2\frac{\sqrt{\pi}\Gamma(\Delta_L+1)}{4\Gamma(\Delta_L+3/2)} + 2\Delta_\Phi \pi \cot(\pi \ep) + \cdots \nn \\ 
&\simeq (\pi \ep)^{2\Delta_L}C_{\Phi}^2\frac{\sqrt{\pi}\Gamma(\Delta_L+1)}{4\Gamma(\Delta_L+3/2)}  +2\Delta_\Phi\lb \frac{1}{\ep} - \frac{\pi^2 \ep}{3} - \frac{\pi^4 \ep^3}{45} - \frac{2\pi^6 \ep^5}{945} + \cdots\rb 
% (C_\alpha-C_{\alpha\beta})^2 \frac{\sqrt{\pi}\Gamma(\Delta_{L}+1)}{4\Gamma(\Delta_{L}+3/2)}(\pi (1-x))^{2\Delta_L}+\cdots
\end{align}
We would like to highlight the appearance of the term $2\Delta_\Phi\pi \cot(\pi \ep)$. This is the expectation value of the full modular Hamiltonian of the vacuum in the state created by $\mO_{L}$:
\begin{align}
    \braket{\mO_\Phi|\hat{K}_0|\mO_\Phi}=2\Delta_\Phi\pi\cot(\pi \ep)\ .
\end{align}

    
% Once again, the divergent piece in the $x\to 1^-$ comes from the last term, while the numerator that computes the vacuum subtracted entanglement entropy is symmetric under $x\to 1-x$:
% % \begin{align}
% %   S(\rho_\beta\|\rho_0)&=4h_\beta\lb\frac{1}{\ep}+\frac{\pi^2\ep}{3}+\cdots\rb+C_\beta^2\ep^{2\Delta_L}\lb -2\Delta_L S_2(1)+S_2'(1)\rb+\cdots\nn\\
% %   &=4h_\beta\lb\frac{1}{\ep}+\frac{\pi^2\ep}{3}+\cdots\rb+C_\beta^2\ep^{2\Delta_L}\lb \frac{2^{-2\Delta_L}\sqrt{\pi}\Gamma(\Delta_L+1)}{4\Gamma(\Delta_L+3/2)}\rb+\cdots 
% % \end{align}
% \begin{align}
%   S(\rho_\Phi\|\rho_0)&=2\Delta_\Phi \lb\frac{1}{\ep}-\frac{\pi^2\ep}{3}+\cdots\rb+C_\Phi^2(2\pi \ep)^{2\Delta_L}\lb -2\Delta_L S_2(1)+S_2'(1)\rb+\cdots\nn\\
%   &=2\Delta_\Phi\lb\frac{1}{\ep}-\frac{\pi^2\ep}{3}+\cdots\rb +C_\Phi^2(\pi\ep)^{2\Delta_L}\lb \frac{\sqrt{\pi}\Gamma(\Delta_L+1)}{4\Gamma(\Delta_L+3/2)}\rb+\cdots 
% \end{align}
% Keeping the terms from the expansion in denominator, which are atleast as big as the numerator, we get
% \begin{align}
%     S(\rho_\Phi\|\rho_0) &= 2\Delta_{\Phi} \lb \frac{1}{\ep} - \sum_{k=1}^{K} \frac{2^{2k}|B_{2k}|} {(2k)!}\pi^{2k}\ep^{2k-1}+\mO\lb\ep^{2\Delta_L + 2}\rb\rb + (\pi \ep)^{2\Delta_L}C_{\Phi}^2\frac{\sqrt{\pi}\Gamma(\Delta_L+1)}{4\Gamma(\Delta_L+3/2)}  \nn \\
%     K &= \lfloor \Delta_L + 1/2\rfloor,
% \end{align}
% A better way to write this is


% \JM{Using the OPE given in eq.~[\ref{eq: general OPE with identity}], we obtain}
% \begin{align}
%     \JM{S(\rho_\Phi\|\rho_0)} &=\JM{
%     % (1-\delta_{\Psi \Phi}) 
%     2\Delta_{\Phi} \lb \frac{1}{\ep} - \frac{\ep}{3} + \cdots\rb } 
% \end{align}
% \pragraph{Relative entropy $S(\rho_\alpha\|\rho_\beta)$ at $x\to 1^-$:}

% As we saw above, the divergence in the relative entropy $S(\rho_\alpha\|\rho_\beta)$ in the $x\to 1^-$ limit comes from the denominator, because the numerator is symmetric under $x\to 1-x$. In this case, we have
% \begin{align}-\braket{\tilde{\mO}_{\alpha|\beta}^{(n-1)}}&=\lb \frac{n \sin(\pi\ep/n)}{\sin(\pi\ep)}\rb^{4(n-1)(h_\alpha-h_\beta)}(2\sin(\pi\ep/n))^{-4(n-1)(h_\alpha-h_\beta)}\times\nn\\
% &\Big(\braket{V_\alpha(z^+_{0,n})V_{-\alpha}(z^-_{0,n})}\nn\\
% &\quad+C_\beta \sum_{k=1}^{n-1} \lb 2\sin\lb\frac{\pi\epsilon}{n}\rb\rb^{\Delta_L}\braket{V_\alpha(z^+_{0,n})V_{-\alpha}(z^-_{0,n})\mO_L(e^{2\pi i k/n})}+\cdots\Big)\nn\\
% &=D_n(\ep)^{n-1}\lb\frac{\sin(\pi\ep/n)}{\sin(\pi(1+\ep)/n)}\rb^{4h_\beta}\times\nn\\
% &\Big(1+C_\beta^2\lb\frac{\sin(\pi\epsilon/n)}{\sin(\pi(1+\epsilon)/n)}\rb^{\Delta_L}\sum_{k=0}^{n-2} \lb 2\cos(\pi(1+\ep)/n)-2\cos(2\pi k/n) \rb^{-\Delta_L} +\cdots\Big)
% \end{align}



%\NL{Change $\alpha$ notation to $\Phi$ and $\Psi$}
{\bf Relative Entropy $S(\rho_0\|\rho_\Phi)$ at $x\to 1^-$:}
We would like to study the $x\to 1^-$ limit of relative entropy $S(\rho_0\|\rho_\Phi)$. In this limit, the replica trick requires us to compute the $2(n-1)$-point function of $\Phi$ in the OPE channel where $\Phi(z_k^+)$ approaches $\Phi^*(z_{k+1}^-)$. In this limit, the lowest order term in $\mathcal{F}_n$ comes from the two-point function at $z_{n-1,n}^+=\exp(\frac{2i\pi}{n}(n-1+x/2))$ and $z_{1,n}^-=\exp(\frac{2i\pi}{n}(1-x/2))$ and its OPE with $\mO_L$ inserted at $z_3 = \exp (i\pi 2k/n)$ for $k = 1,2, \cdots, n-2.$:
\begin{align}
  % &-\log\mathcal{F}_n=-\log\braket{\tilde{\mO_{\beta}^{(n-1)}}}\nn\\
  &\log\mathcal{F}_n=-\log\braket{\tilde{\mO}_{\Phi}^{(n-1)}}\nn \\
% &\braket{\tilde{\mO}_{\beta}^{(n-1)}}=\lb \frac{n \sin(\pi\ep/n)}{\sin(\pi\ep)}\rb^{-4(n-1)h_\beta}(2\sin(\pi\ep/n))^{4h_\beta}\times\nn\\
% &\lb\braket{V_\beta(z^+_{n-1,n})V_{-\beta}(z^-_{1,n})}+C_\beta\sum_{k=1}^{n-2} \lb 2\sin\lb\frac{\pi\epsilon}{n}\rb\rb^{\Delta_L}\braket{V_\beta(z^+_{n-1,n})V_{-\beta}(z^-_{1,n})\mO_L(e^{2\pi i k/n})}+\cdots\rb \nn \\
&\braket{\tilde{\mO}_{\Phi}^{(n-1)}}=\lb \frac{n \sin(\pi\ep/n)}{\sin(\pi\ep)}\rb^{-4(n-1)h_\Phi}(2\sin(\pi\ep/n))^{4h_\Phi}e^{2\pi i(n-1)s_\Phi  }e^{-2\pi i(n-2) s_\Phi}e^{i\pi s_\Phi}\times \nn\\
&\lb\braket{\Phi(z^-_{1,n})\Phi^*(z^+_{n-1,n})}+C_\Phi\sum_{k=1}^{n-2} \lb 2\sin\lb\frac{\pi\epsilon}{n}\rb\rb^{\Delta_L}e^{i\pi s_L((2k+1)/n + 1/2)}\braket{\Phi(z^-_{1,n})\Phi^*(z^+_{n-1,n})\mO_L(e^{ i\pi (2k+1)/n}})+\cdots\rb
% &\JM{\braket{\tilde{\mO}_{\beta}^{(n-1)}}=\lb \frac{n \sin(\pi\ep/n)}{\sin(\pi\ep)}\rb^{-4(n-1)h_\beta}\lb\frac{\sin(\pi\ep/n)}{\sin (\pi \ep)}\rb^{4h_\beta}\times}\nn\\
% &\lb\braket{V_\beta(z^+_{n-1,n})V_{-\beta}(z^-_{1,n})}+C_\beta\sum_{k=1}^{n-2} \lb 2\sin\lb\frac{\pi\epsilon}{n}\rb\rb^{\Delta_L}\braket{V_\beta(z^+_{n-1,n})V_{-\beta}(z^-_{1,n})\mO_L(e^{2\pi i k/n})}+\cdots\rb \nn .
%&=\lb 2\sin\lb\frac{\pi+\epsilon}{2 n}\rb\rb^{-4h_\beta}\lb 1+C_\beta \lb 4\sin\lb \frac{\pi(2k+1)+\epsilon/2}{2n}\rb\sin\lb\frac{\pi(2k-1)-\epsilon/2}{2n}\rb\rb^{-\Delta_L}+\cdots \rb\nn\\
\end{align}

We consider the three-point function  $\braket{\Phi(z^+_{n-1,n})\Phi^*(z^-_{1,n}) \mO_L(e^{i\pi (2k+1)/n})}$ as the correction term, where we sum over the indices $k = 1, \cdots , n-2$. It is computed by taking the two-point function of the operators $\mO(e^{i \pi (2k+1)/n})$ and the operator from the OPE of the pair $\Phi^*(z^+_{n-1,n})$ and $\Phi(z^+_{1,n})$. We have, 
\begin{align}
  % &-\log\mathcal{F}_n=-\log\braket{\tilde{\mO_{\beta}^{(n-1)}}}\nn\\
  &\log\mathcal{F}_n=-\log\braket{\tilde{\mO}_{\Phi}^{(n-1)}}\nn \\
% &\braket{\tilde{\mO}_{\beta}^{(n-1)}}=\lb \frac{n \sin(\pi\ep/n)}{\sin(\pi\ep)}\rb^{-4(n-1)h_\beta}(2\sin(\pi\ep/n))^{4h_\beta}\times\nn\\
% &\lb\braket{V_\beta(z^+_{n-1,n})V_{-\beta}(z^-_{1,n})}+C_\beta\sum_{k=1}^{n-2} \lb 2\sin\lb\frac{\pi\epsilon}{n}\rb\rb^{\Delta_L}\braket{V_\beta(z^+_{n-1,n})V_{-\beta}(z^-_{1,n})\mO_L(e^{2\pi i k/n})}+\cdots\rb \nn \\
&\braket{\tilde{\mO}_{\Phi}^{(n-1)}}=\lb \frac{n \sin(\pi\ep/n)}{\sin(\pi\ep)}\rb^{-4(n-1)h_\Phi}(2\sin(\pi\ep/n))^{4h_\Phi}e^{2\pi i(n-1)s_\Phi  }e^{-2\pi i(n-2) s_\Phi}e^{i\pi s_\Phi}\times \nn\\
&\lb\braket{\Phi(z^+_{n-1,n})\Phi^*(z^-_{1,n})}+C_\Phi^2 \sum_{k=1}^{n-2} \lb 2\sin\lb\frac{\pi\epsilon}{n}\rb\rb^{\Delta_L}e^{i\pi s_L((2k+1)/n + 1/2)}\braket{\Phi(z^+_{n-1,n})\Phi^*(z^-_{1,n}) \mO_L(e^{i\pi (2k+1)/n})})+\cdots\rb
% &\JM{\braket{\tilde{\mO}_{\beta}^{(n-1)}}=\lb \frac{n \sin(\pi\ep/n)}{\sin(\pi\ep)}\rb^{-4(n-1)h_\beta}\lb\frac{\sin(\pi\ep/n)}{\sin (\pi \ep)}\rb^{4h_\beta}\times}\nn\\
% &\lb\braket{V_\beta(z^+_{n-1,n})V_{-\beta}(z^-_{1,n})}+C_\beta\sum_{k=1}^{n-2} \lb 2\sin\lb\frac{\pi\epsilon}{n}\rb\rb^{\Delta_L}\braket{V_\beta(z^+_{n-1,n})V_{-\beta}(z^-_{1,n})\mO_L(e^{2\pi i k/n})}+\cdots\rb \nn .
%&=\lb 2\sin\lb\frac{\pi+\epsilon}{2 n}\rb\rb^{-4h_\beta}\lb 1+C_\beta \lb 4\sin\lb \frac{\pi(2k+1)+\epsilon/2}{2n}\rb\sin\lb\frac{\pi(2k-1)-\epsilon/2}{2n}\rb\rb^{-\Delta_L}+\cdots \rb\nn\\
\end{align}
If we naively truncate the denominator at the lowest order, and analytically continue in $n$ we find that this term results in a large negative relative entropy! This clearly means that our OPE expansion is invalid. What has gone wrong is that operators $\Phi(z^+_{n-1},n)\Phi^*(z^-_{1,n})$ are separated by angle $2\pi(1+x)/n$ which in the limit $n\to 1$ is larger than $2\pi$! 



\subsubsection{Some examples}

For the massless free boson theory, we consider two types of conformal primaries, the coherent vertex operator $V_{\alpha,\alpha}$, and the primary $i(\p\varphi+\bar{\p}\bar{\varphi})=i\p_t\phi$. The relative entropies in these cases are computed in \cite{}. They diverge in large subsystem limit $\ep=(1-x)\to 0^+$ as
\begin{align}\label{eq: free massless boson result large x limit}
&S(\rho_\alpha\|\rho_\beta)=2(\alpha-\beta)^2 (1-\pi x \cot(\pi x))\simeq_{x\to 1^-} 2(\alpha-\beta)^2\lb \frac{1}{\ep}-\frac{\pi^2\ep}{3}+\frac{\pi^2\ep^2}{3}+O(\ep^3)\rb\\
&S(\rho_{i\partial\phi}\|\rho_0)=2(1-\pi x \cot(\pi x))+2\lb \log(2\sin(\pi x))+\psi_0(\csc(\pi x)/2)+\sin(\pi x)\rb\\
&\simeq_{x\to 1^-} \frac{2}{\ep}-\frac{2\pi^2}{3}\ep+\cdots\\
&S(\rho_{i\partial\phi}\|\rho_\beta)=S(\rho_{i\partial\phi}\|\rho_0)+S(\rho_0\|\rho_\beta)\simeq_{x\to 1^-} (2+\beta^2)\lb \frac{1}{\ep}-\frac{\pi^2}{3}\ep+\cdots\rb 
\end{align}
The goal of this subsection is to reproduce these divergences from the replica trick.
% Note that the divergent piece of the relative entropy in the $x\to 1^-$ limit agrees with what we found above.




This term is precisely the divergent term we get from the contribution from just the two point function of the mixing OPE. There should be a factor in the denominator of $\mathcal{F}_n$ from the contribution of higher order terms, that will exactly cancel the numerator, and thus reproducing the massless free boson result (once we replace the general operators with Vertex operators. 

For non-chiral Vertex operators, the OPEs are given by,
\begin{align}
    V_\alpha(z_{k+1}^-)V_{-\alpha}(z_k^+)&= \zeta_n^{-2\alpha^2} [1 + \alpha \zeta_n J(e^{i\pi \theta_k}) + \cdots],\\
    V_\alpha(z_{k+1}^-)V_{-\beta}(z_k^+) &= \zeta_n^{-2\alpha\beta}V_{\alpha-\beta}(e^{i\pi \theta_k})\left[ 1+ \alpha \zeta_n J(e^{i\pi \theta_k}) + \cdots \right].
\end{align}
These OPEs are carried out at the midpoint of the two operators. This is also equivalent to the OPE,
\begin{align}
    V_\alpha(z_{k+1}^-)V_{-\alpha}(z_k^+) &= \zeta_n^{-2\alpha^2}[1 + \alpha \zeta_n J(z_k^+) + \cdots],\\
    V_\alpha(z_{k+1}^-)V_{-\beta}(z_k^+) &= \zeta_n^{-2\alpha\beta}V_{\alpha-\beta}(z_{k}^+)\left[ 1+ \alpha \zeta_n J(z_k^+) + \cdots \right].
\end{align}
The second OPE can be further approximated by,
\begin{align}
    V_\alpha(z_{k+1}^-)V_{-\beta}(z_k^+) &= \zeta_n^{-2\alpha\beta}V_{\alpha-\beta}(z_{k}^+)\left[ 1+ \alpha \zeta_n J(e^{i \pi \theta_k}) + \cdots \right]
\end{align}
where we Taylor expanded the current $J(z_k^+)\sim J(e^{i\pi \theta_k}) + \frac{\zeta_n}{2}\p J(e^{i\pi\theta_k}) + \cdots$, and we keep only terms of the order of $\zeta_n$.

Using this, the denominator of $\mathcal{F}_n$ gets modified, 
\begin{align}
    D_n(\ep)^{n-1}\zeta_n^{2(\alpha-\beta)^2}V_{\alpha- \beta}(e^{i\pi/n}) &(1 + \beta \zeta_n J(e^{i\pi /n}))\times \\
    &V_{\beta - \alpha} (e^{i\pi (2n-1)/n })(1 +\alpha \zeta_n J(e^{i\pi (2n-1)/n}))\prod_{k=1}^{n-2} (1 + \beta J(e^{i \pi \theta_k}))
\end{align}
Note, the term $\zeta_n^{2(\alpha - \beta)^2}$ follows directly from the OPEs we consider in the large interval limit, given in eq.~[\ref{eq: general OPE with identity}]. 
At the lowest order in $\zeta_n$, we have a two-point function of vertex operators of opposite charge $(\alpha-\beta)$. At the next order, we have to include a $J$ in this correlator; however, by conservation of charge, this three-point function vanishes. At the next order, we find  
\begin{align}
    D_n(\ep)^{n-1}
    \zeta_n^{2(\alpha-\beta)^2} \braket{V_{\alpha- \beta}(e^{i\pi/n})V_{\beta - \alpha} (e^{i\pi (2n-1)/n })} \lb 1 + \zeta_n^{2}g(n)\rb
\end{align}
where $g(n)$ is given in eq.~[\ref{gnw1}]. 
\begin{align}\label{eq: coherent states rel entr}
    \p_n \log \mathcal{F}_n &= \p_n \log \left[\frac{1+\zeta_n^{2} f(n)+\cdots}{D_n(\ep)^{n-1}\zeta_n^{2(\alpha-\beta)^2} (2\sin(\pi (1-\ep)/n))^{-2(\alpha-\beta)^2}\lb 1+ \zeta_n^2 g(n) +\cdots\rb}\right] \nn \\
    &= 2(\alpha - \beta)^2 \lb \frac{1}{\ep} - \frac{\pi^2 \ep}{3}+\mO(\ep^3)\rb + 2(\alpha- \beta)^2\frac{\lb \pi \ep \rb^2}{3} + \mO(\ep^4)\cdots\nn \\
    &= 2(\alpha - \beta)^2 \lb \frac{1}{\ep} - \frac{\pi^2\ep}{3}+ \frac{\pi^2\ep^2}{3}+ \mO(\ep^3)\rb .
\end{align}
This matches with the massless free boson result given in eq.~[\ref{eq: free massless boson result large x limit}]. This shows that we have to keep primaries in the mixing OPE until the dimension of which matches with that of numerator, in order to obtain the exact result. 

This calculation can be repeated for the coherent states to compute the relative entropy with respect to vacuum $S(\rho_0 \| \rho_\Phi)$.
\begin{align}
    V_{-\beta}(z_{1,n}^-) &= V_{-\beta}(z_{0,n}^+) - 2\zeta_n\p V_{-\beta}(z_{0,n}^+) + \cdots \\
    &=V_{-\beta}(z_{0,n}^+) + 2\zeta_n \beta J  V_\beta(z_{0,n}^+) + \cdots
\end{align}
where we used $\p V_{-\beta} = -\beta J V_\beta$ in the last line. 
% Similarly, for the other vertex operator sitting close to identity,
% \begin{align}
%     V_{-\beta}(z_{1,n}^-) &= V_{-\beta}(z_{0,n}^+) + 2\zeta_n\p V_{-\beta}(z_{0,n}^+) + \cdots \\
%     &=V_{-\beta}(z_{0,n}^-) + 2\zeta_n \beta J  V_{-\beta}(z_{0,n}^-) + \cdots
% \end{align}
Thus, in the denominator we have the following,
\begin{align}
    D_n(\ep)^{n-1} \zeta_n^{2\beta^2} V_{-\beta}(z_{0,n}^+)\lb V_\beta(z_{0,n}^-) (1+\beta \zeta_nJ(e^{i\pi (2n-1)/n}))\rb \prod_{k=1}^{n-2}(1+\beta J(e^{i\pi \theta_k}))
\end{align}
\JM{Note the asymetry we considered for the shifting of the Vertex operator sitting close to $\theta_1 = e^{i\pi/n}$}.
We can substitute $\alpha = 0$ corresponding to the identity operator, we get in the denominator,
\begin{align}
    \frac{\beta^2}{2}\sum_{k\neq k'=1}^{n-1}\frac{1}{(2\sin(\pi(k-k')/n))^{2\beta^2}} = \beta^2 S_2(n,\beta^2) - \frac{2S_2(n,\beta^2)}{n}
\end{align}
The derivative of which is given by,
\begin{align}
    -C_\Phi^2 S_2'(2,\beta^2).
\end{align}
And thus reproduce the relative entropy of vacuum with respect to coherent states. 
\begin{align}
    S(\rho_0 \| \rho_\Phi) = 2\beta^2\lb \frac{1}{\ep} - \frac{\pi^2\ep}{3} + \frac{\pi^2\ep^2}{3} + \cdots \rb
\end{align}
This reduces to a particular case of the general relative entropy we calculated previously for coherent states, namely $S(\rho_\alpha \| \rho_\beta)$. 

\JM{We want to compute the three point function, to look for the contribution of spin of primaries $\mO_L$ and $\mO_{L'}$. We will compute $\braket{\mO_{L'}(z_0^-) \mO_{L'}(z_0^+) \mO_L(z_k^+)}$ and take the sum over $k =1,\cdots , n-2$. The extra phase term we have to take into account would come from the primary of the OPE of $\Phi^*(z_{k+1}^-) \Phi(z_k^+)$, which is $e^{i \pi s_L(\theta_k + 1/2)}$. The common factors and phases that carry spin are given by,}
\begin{align}
    &\JM{C_{\Psi\Phi}^2C_\Phi e^{3\pi is_{L'}}e^{i\pi s_L/2}\lb\frac{\zeta_1}{\zeta_n}\rb^{2((n-1)\Delta_\Phi + \Delta_\Psi)}\zeta_n^{2\Delta_{L'} + \Delta_L}\times} 
    % &\qquad \qquad \qquad 
    \JM{\sum_{k=1}^{n-2}\braket{\mO_{L'}(z_0^-) \mO_{L'}(z_0^+) \mO_L(z_k^+)} e^{i \pi s_L \theta_k} }
\end{align}
\JM{We can then write the three point function,}
\begin{align}
    \JM{\braket{\mO_L(z_k^+) \mO_{L'}(z_0^+)\mO_{L'}(z_0^-) }} &= \JM{C_{\Psi\Phi}^2C_\Phi\frac{e^{-i3\pi (2s_{L'} -s_L)/2}}{\lb2\sin(\pi(1-\ep)/n)\rb^{2\Delta_{L'}-\Delta_L}}\times}\nn \\ 
    & \qquad \qquad \qquad \qquad 
    \JM{\sum_{k=1}^{n-2}\Bigg( \frac{e^{-i\pi s_L(k/n+3/2)}}{\lb2\sin(\pi(k+1-\ep)/n)\rb^{\Delta_L}} \frac{e^{-i\pi s_L((k+1-\ep)/n+1/2)}}{\lb 2\sin(\pi k/n)\rb^{\Delta_L}}\Bigg)} \nn \\
    &= \JM{C_{\Psi\Phi}^2C_\Phi e^{-3\pi is_{L'} } e^{-i\pi s_L/2}\lb 2\sin (\pi (1-\ep)/n)\rb^{\Delta_L-2\Delta_{L'}}\times }\nn \\
    &\qquad \qquad \qquad \qquad \JM{\sum_{k=1}^{n-2}\frac{e^{-i\pi s_L (2k+1-\ep)/n}}{\lb  4\sin(\pi (k+1-\ep)/n)\sin(\pi k /n)\rb^{\Delta_L}}} 
\end{align}
% \JM{where in the last step we multiplied the phase $e^{i \pi s_L(\theta_k + 1/2)}$ sitting outside. 
\JM{Hence we get,}
\begin{align}
    &\JM{\mathcal{F}_n = \frac{1+ \zeta_n^{2\Delta_L}f(n)}{(D_n(\ep))^{n-1}\lb \frac{\sin(\pi \ep/n)}{\sin (\pi(1-\ep)/n)}\rb^{2\Delta_{L'}}C_{\Psi \Phi}^2
    \lb 1 + C_\Phi \, \mathcal{T}(n)\rb} } \\
    &\JM{\mathcal{T}(n)} = \JM{\lb \sin(\pi \ep/n) \sin(\pi (1-\ep)/n) \rb^{\Delta_L}\sum_{k=1}^{n-2} \lb \frac{1}{\sin(\pi (k+1-\ep)/n) \sin (\pi k /n)}\rb^{-\Delta_L}}
\end{align}
\JM{We can compute the relative entropy $S(\rho_\Psi \| \rho_\Phi)$,}
\begin{align}
    \JM{S(\rho_\Psi \| \rho_\Phi) = \p_n \log \mathcal{F}_n \big|_{n=1}} & \simeq  \JM{2\Delta_{L'}\lb  \frac{1}{\ep} - \sum_{k=1}^{K} \frac{2^{2k}|B_{2k}|} {(2k)!}\ep^{2k-1}
    % - \frac{\pi^2 \ep}{3}
    \rb + (\pi \ep)^{2\Delta_L}C_{\Psi}^2\frac{\sqrt{\pi}\Gamma(\Delta_L+1)}{4\Gamma(\Delta_L+3/2)} } \nn \\ &   \JM{- C_\Phi  \p_n \mathcal{T}(n)|_{n=1} + \cdots  }
\end{align}
\JM{After using the general OPE with the identity piece given in eq.~[\ref{eq: general OPE with identity}], we get}
\begin{align}
    \JM{S(\rho_\Psi \| \rho_\Phi) = \p_n \log \mathcal{F}_n \big|_{n=1}} & \simeq  \JM{
    % (1-\delta_{\Psi \Phi}) 
    \lb 2\Delta_{L'}\lb  \frac{1}{\ep} - \frac{\ep}{3} + \cdots
    % - \frac{\pi^2 \ep}{3}
    \rb  - C_\Phi  \p_n \mathcal{T}(n)|_{n=1} + \cdots \rb }
\end{align}
% ----
%  In the denominator the first two insertions are $\alpha$ and the rest are $\beta$, therefore the product $\prod_{k=0}^{n-1}\delta_{\omega_k\omega_{k+1}}\sim\delta^2_{\alpha\beta}=0$. Since there is no identity operator in the OPE of $V_\alpha$ and $V_{-\beta}$, the only non-vanishing contribution to the denominator comes from the case where we take the OPE of $V_\alpha$ and $V_{-\beta}$ at $z_0^+$ and $z_1^-$ and $V_\beta$ and $V_{-\alpha}$ at $z_0^-$ and $z_1^+$. Therefore,
%   \begin{align}
%     g_2(n)=C_{\alpha\beta}^L \lb 2\sin(\pi \ep/n)\rb^{-(\Delta_\alpha+\Delta_\beta)}e^{-2i\pi (\delta s)}
% (2\sin(\pi \ep/n))^{2\Delta_L}\braket{O_L(0)\mO_L(e^{2\pi i/n})}+\cdots
% \nn
%   \end{align}

% ----

% Here, we establish that in the limit of a large subsystem $x\to 1^-$, relative entropy of excited states is universally given by
% \begin{align}
%   S(\rho_\alpha\|\rho_\beta)=
% \end{align}


  \section{Relative entropy of mixed excited states}

  

  We aim to compute the relative entropy $S(\rho || \sigma_{y})$ where
  $\rho = \rho_{\alpha}$ and the reference state is the mixture
  $\sigma_{y}= y\rho_{\alpha}+(1-y) \rho_{\beta}$ for $0\leq y\leq 1$. The standard
  replica trick for relative entropy relates this quantity to the analytic continuation
  of trace functionals that break replica symmetry. Specifically, for the mixed state
  case, we utilize the following identity derived from the generalized replica
  trick for normalized density matrices:
  % \begin{eqnarray}
  %   S\left(\rho_\alpha \Big\| y\rho_\alpha+(1-y)\rho_\beta\right)&&=\lb \partial_{n}
  %   \log \left[ \frac{\text{tr}(\rho_{\alpha}^{n})}{\text{tr}(\rho_{\alpha}(y\rho_{\alpha}+(1-y)\rho_{\beta})^{n-1})}
  %   \right]\rb_{n\to 1}\nn\\
  %   && =\lb\partial_n\log\left[
  %   \frac{\braket{\prod_{i=1}^{n} \tilde{V}_{\alpha}(u^+_i)\tilde{V}_{-\alpha}(u^-_i)}_{\mathcal{C}_n}}{\braket{\tilde{V}_\alpha(u^+_n)\tilde{V}_{-\alpha}(u^-_n)\mO^{(n-1)}_{\alpha|\beta}}_{\mathcal{C}_n}}\right]\rb_{n=1}
  % \end{eqnarray}
\begin{eqnarray}
    S\left(\rho_\Psi \Big\| y\rho_\Psi+(1-y)\rho_\Phi\right)&&=\lb \partial_{n}
    \log \left[ \frac{\text{tr}(\rho_{\Psi}^{n})}{\text{tr}(\rho_{\Psi}(y\rho_{\Psi}+(1-y)\rho_{\Phi})^{n-1})}
    \right]\rb_{n\to 1}\nn\\
    && =\lb\partial_n\log\left[
    \frac{\braket{\prod_{i=1}^{n} \tilde\Psi(u^+_i)\tilde\Psi^*(u^-_i)}_{\mathcal{C}_n}}{\braket{\tilde\Psi(u^+_n)\tilde\Psi^*(u^-_n)\mO^{(n-1)}_{\Psi|\Phi}}_{\mathcal{C}_n}}\right]\rb_{n=1}
  \end{eqnarray}
  where
  % \begin{eqnarray}
  %   &&\mO_{\alpha|\beta}^{(n-1)}=\sum_{p=0}^{n-1}y^{n-1-p}(1-y)^p\sum_{\substack{w\in\{\alpha,\beta\}^{n-1}\\ N_\beta(w)=p}}
  %   \tilde{V}_{w_1}(u_1^+)\tilde{V}_{-w_1}(u_1^-)\cdots \tilde{V}_{w_{n-1}}(u_{n-1}^+)\tilde{V}_{-w_{n-1}}(u_{n-1}^-).\nn\\
  % \end{eqnarray}
  \begin{align}
    \mO_{\Psi|\Phi}^{(n-1)}=\sum_{p=0}^{n-1}y^{n-1-p}(1-y)^p\sum_{\substack{w\in\{\Psi,\Phi\}^{n-1}\\ N_\beta(w)=p}}
    \tilde{\mO}_{w_1}(u_1^+)\tilde{\mO}_{w_1^*}(u_1^-)\cdots \tilde{\mO}_{w_{n-1}}(u_{n-1}^+)\tilde{\mO}_{w_{n-1}^*}(u_{n-1}^-).
  \end{align}
  where $\tilde\mO_{\omega_k} = \omega_k$, $\mO_{\omega^*_k} = \omega_k^*$.
  
  For compactness of notation, we suppress the anti-holomorphic part and write only the holomorphic vertex operators. Throughout the calculation, however, all correlation functions and OPE coefficients are understood to include both holomorphic and anti-holomorphic sectors. In particular, whenever 
  % $\tilde{V}_\omega(z)$ 
  $\tilde{\mO}_\omega(z)$ appears, it should interpreted as the full local primary 
  % $\tilde{V}_\omega(z,\bar z)$ 
  $\tilde{\mO}_\omega(z,\bar z)$, and all conformal dimensions are understood as the full scaling dimensions $\Delta_\omega = h_\omega + \bar h_\omega$.
  Here $w=(w_{1},\dots,w_{n-1})$ is an ordered word and $N_{\Phi}(w)$ counts how
  many letters equal $\Phi$, and the points are on the $n$-cover of cylinder
  $\mathcal{C}_{n}$ at $\pm i\infty$ on each sheet.

 


  Further assuming that $h_{\alpha}=h_{-\alpha}$ the relative entropy replica
  trick in terms of the correlator on the complex plane gives
  % \begin{eqnarray}\label{CorrFn_mixed}
  %   &&S(\rho_\alpha\|y\rho_\alpha+(1-y)\rho_\beta)=\lb \partial_n\log\mathcal{F}_n(\alpha|\beta;y)\rb_{n\to 1}\nn\\
  %   &&\mathcal{F}_n(\alpha|\beta;y)\equiv 
  %   \frac{\braket{\prod_{k=1}^{n} \tilde{V}_{\alpha}(z^+_{k,n})\tilde{V}_{-\alpha}(z^-_{k,n})}}{\braket{\tilde{V}_\alpha(z^+_{n,n})\tilde{V}_{-\alpha}(z^-_{n,n})\tilde{\mO}^{(n-1)}_{\alpha|\beta}}}\nn\\
  %   &&\tilde{\mO}_{\alpha|\beta;y}^{(n-1)}=\sum_{p=0}^{n-1}y^{n-1-p}(1-y)^p\: n^{4p(h_\alpha-h_\beta)}
  %   \sum_{\substack{w\in\{\alpha,\beta\}^{n-1}\\ N_\beta(w)=p}} \prod_{k=1}^{n-1}\frac{V_{w_k}(z_{k,n}^{+})V_{-w_k}(z_{k,n}^{-})}{\braket{V_{w_k}(z^+_{1,1})V_{-w_k}(z^-_{1,1})}}\nn\\
  % \end{eqnarray}
\begin{align}\label{CorrFn_mixed}
    S(\rho_\Psi\|y\rho_\Psi+(1-y)\rho_\Phi)&=\lb \partial_n\log\mathcal{F}_n(\Psi|\Phi;y)\rb_{n\to 1}\nn\\
    \mathcal{F}_n(\Psi|\Phi;y)&\equiv 
    \frac{\braket{\prod_{k=1}^{n} \tilde\Psi(z^+_{k,n})\tilde\Psi^*(z^-_{k,n})}}{\braket{\tilde\Psi(z^+_{n,n})\tilde\Psi^*(z^-_{n,n})\tilde{\mO}^{(n-1)}_{\Psi|\Phi}}}\nn\\
    \tilde{\mO}_{\Psi|\Phi;y}^{(n-1)}&=\sum_{p=0}^{n-1}y^{n-1-p}(1-y)^p\: n^{4p(h_\Psi-h_\Phi)}
    \sum_{\substack{w\in\{\Psi,\Phi\}^{n-1}\\ N_\beta(w)=p}} \prod_{k=1}^{n-1}\frac{\mO_{w_k}(z_{k,n}^{+})\mO_{w_k^*}(z_{k,n}^{-})}{\braket{\mO_{w_k}(z^+_{1,1})\mO_{w_k^*}(z^-_{1,1})}}\nn\\
  \end{align}
where $\mO_{\omega_k} = \omega_k$ and $\mO_{\omega_k^*} = \omega_k^*$.

For simplicity, we often supppress the notation $(\alpha|\beta;y)$ in $\mathcal{F}_n$.
The analytic continuation of the numerator in $\mathcal{F}_n$ after the $n$-derivative and the limit $n\to 1$ contributes as the vacuum subtracted entropy
% \begin{align}
%   \left[\partial_n\log(\tr(\rho_\alpha^n))-\log(\tr(\rho_0^n))\right]_{n\to 1}=S(\rho_0)-S(\rho_\alpha)
% \end{align}
\begin{align}
  \left[\partial_n\log(\tr(\rho_\Psi^n))- \partial_n\log(\tr(\rho_0^n))\right]_{n\to 1}=S(\rho_0)-S(\rho_\Psi)
\end{align}
whereas the analytic continuation of the denominator of $\mathcal{F}_n$ contributes as 
\begin{align}
-\left[\partial_n\log(\tr(\rho_\alpha \sigma_y^{n-1}))-\log(\tr(\rho_0^n)) \right]_{n\to 1}&=\tr(\rho_\alpha\log\sigma_y)+S(\rho_0)\nn\\
&=\tr(\rho_\alpha(\log\sigma_y-\log \rho_0))-\tr((\rho_0-\rho_\alpha)\log\rho_0)\nn\\
&=\braket{V_\alpha|(H_0-H_{\sigma_y})|V_\alpha}-\tr((\rho_\alpha-\rho_0)H_0)
\end{align}
\begin{align}
\left[\partial_n\log(\tr(\rho_\Psi \sigma_y^{n-1}))-\partial_n\log(\tr(\rho_0^n)) \right]_{n\to 1}&=\tr(\rho_\Psi\log\sigma_y)+S(\rho_0)\nn\\
&=\tr(\rho_\Psi(\log\sigma_y-\log \rho_0))-\tr((\rho_0-\rho_\Psi)\log\rho_0)\nn\\
&=\braket{\Psi|(H_0-H_{\sigma_y})|\Psi}-\tr((\rho_\Psi-\rho_0)H_0)
\end{align}
where 
% $\sigma_y=y\rho_\alpha+(1-y)\rho_\beta$ 
$\sigma_y=y\rho_\Psi+(1-y)\rho_\Phi$), and we have the modular Hamiltonians
\begin{align}
  H_{\sigma_y}&:=-\log\sigma_y; \qquad H_0=-\log\rho_0.
\end{align}
The reason we split the terms this way is that in vacuum 2d CFTs $H_0$ has a universal form that is a local integral over stress tensor. The modular Hamiltonian of the vacuum on a cylinder of circumference $2\pi$ in the interval $(-x,x)$ is given by
\begin{align}
H_0&=\frac{(2\pi)^2}{2}\int_{-x}^x\frac{\cos(\pi y)-\cos(\pi x)}{\sin(\pi x)}T_{00}(x).
\end{align}
In a state with uniform energy density such as an exact energy eigenstate with dimension $h_\alpha$ on a cylinder of circumference $2\pi$ we have $\Delta T_{00}(y)=h/\pi$, and as a result
% \begin{align}
% \Delta H_\sigma\equiv \text{tr}((\rho_\alpha-\rho_0)H_0)=4h_\alpha(1-\pi x\cot(\pi x))
% \end{align}
\begin{align}
\Delta H_\sigma\equiv \text{tr}((\rho_\alpha-\rho_0)H_0)=4h_{\Psi}(1-\pi x\cot(\pi x))
\end{align}
This means that in computing the contribution of the exact energy eigenkets $\ket{V_\alpha}$ we should always expect a term of the form
% \begin{align}
%   -4h_\alpha(1-\pi x\cot(\pi x)).
% \end{align}
\begin{align}
  -4h_{\Psi}(1-\pi x\cot(\pi x)).
\end{align}
In other words,
\begin{align}
S(\rho_\alpha\|\sigma_y)=S(0)-S(\rho_\alpha)+\braket{V_\alpha|(H_0-H_{\sigma_y})|V_\alpha}-4h_\alpha(1-\pi x\cot(\pi x)).
\end{align}
\begin{align}
\JM{S(\rho_\Psi\|\sigma_y)=S(0)-S(\rho_\Psi)-\braket{\Psi|(H_0-H_{\sigma_y})|\Psi}+4h_\Psi(1-\pi x\cot(\pi x)).}
\end{align}
This way of splitting the relative entropy has the advantage that it makes the analysis of regions that are larger than half, i.e., $x>1/2$, clearer. Note that the $\Delta S$ term is symmetric under $x\to 1-x$. The term proportional to $h_\alpha$ at $x\to 0$ starts at order $x^2$, whereas at $x\to 1^-$ it has a simple pole $(1-x)^{-1}$. The remainder term is the vacuum subtracted modular Hamiltonian whose $x\to 1-x$ behavior is to be studied. 

As we will see, in free massless boson, the coherent states of charge $\alpha$ have the property that $\rho_\alpha=U_\alpha \rho_0 U_\alpha^\dagger$ for some unitary operator $U_\alpha$. In this case, the relative entropy simplifies to
\begin{align}
S(\rho_\alpha\|y\rho_\alpha+(1-y)\rho_\beta)=S(\rho_0\|y\rho_0+(1-y)\rho_{\beta-\alpha}).
\end{align}



Another special case is $\rho_0$ the vacuum reduced density matrix. In this case, the numerator drops out, and the denominator becomes
\begin{align}
S(\rho_0\|\sigma_y)&=-\tr(\rho_0(H_{\sigma_y}-H_0))=\braket{H_{\sigma_y}-H_0}\nn\\
\end{align}
\begin{align*}
S(\rho_0\|\sigma_y)&=\JM{\tr(\rho_0(H_{\sigma_y}-H_0))=\braket{H_{\sigma_y}-H_0}}\nn\\
\end{align*}





  % The most non-trivial component of this calculation is the term
  % $\text{tr}\!\left (\rho_{\alpha}\,(y\rho_{\alpha}+(1-y)\rho_{\beta})^{n-1}\right
  % )$. Since the operators do not commute in general, one should expand the \emph{power}
  % as a sum over all words of length $n-1$ in the alphabet $\{\rho_{\alpha},\rho_{\beta}
  % \}$ rather than using a naive binomial theorem. In the case $\alpha=0$ we
  % write $\rho_{0}$ for the vacuum reduced density matrix and obtain the precise
  % finite-$n$ identity
  % \begin{equation}
  %   \label{eq:word-expansion-alpha0}\tr\!\left(\rho_{0}\,(y\rho_{0}+(1-y)\rho_{\beta}
  %   )^{n-1}\right) = \sum_{p=0}^{n-1}y^{\,n-1-p}(1-y)^{p}\sum_{\substack{w\in\{0,\beta\}^{n-1}\\ N_\beta(w)=p}}
  %   \tr\!\left(\rho_{0}\,\rho_{w_1}\rho_{w_2}\cdots \rho_{w_{n-1}}\right),
  % \end{equation}
  % where $w=(w_{1},\dots,w_{n-1})$ is an ordered word and $N_{\beta}(w)$ counts
  % how many letters equal $\beta$.

  % Equivalently, one may separate the ``naive binomial'' term (in which all
  % $\rho_{\beta}$ are taken to sit to the right) plus an explicit remainder built
  % from commutators. A convenient way to package this remainder is via the
  % adjoint action $\mathrm{ad}_{\rho_0}(X)\equiv[\rho_{0},X]$:
  % \begin{equation}
  %   \label{eq:binomial-plus-comm-alpha0}\tr\!\left(\rho_{0}\,(y\rho_{0}+(1-y)\rho
  %   _{\beta})^{n-1}\right) = \sum_{p=0}^{n-1}\binom{n-1}{p}y^{\,n-1-p}(1-y)^{p}\,
  %   \tr\!\left(\rho_{0}^{\,n-p}\rho_{\beta}^{\,p}\right) +\mathcal{C}_{n}(y),
  % \end{equation}
  % with the commutator correction
  % \begin{equation}
  %   \label{eq:Cn-def}\mathcal{C}_{n}(y) \equiv \sum_{p=0}^{n-1}y^{\,n-1-p}(1-y)^{p}
  %   \sum_{\substack{w\in\{0,\beta\}^{n-1}\\ N_\beta(w)=p}}\tr\!\left(\rho_{0}\,\Delta
  %   _{w}\right),\qquad \Delta_{w}\equiv \rho_{w_1}\cdots\rho_{w_{n-1}}-\rho_{0}^{\,n-1-p}
  %   \rho_{\beta}^{\,p}.
  % \end{equation}
  % Each $\Delta_{w}$ can be reduced to a finite sum of nested commutators by repeatedly
  % commuting $\rho_{0}$ past $\rho_{\beta}$ (so $\mathcal{C}_{n}(y)$ vanishes
  % whenever $[\rho_{0},\rho_{\beta}]=0$; in particular, for the free boson vertex-operator
  % states considered below the ordering dependence drops out).

  % Using the operator-state correspondence, the trace
  % $\text{tr}(\rho_{\alpha}^{n-p}\rho_{\beta}^{p})$ corresponds to a path
  % integral on an $n$-sheeted Riemann surface with specific operator insertions. In
  % the limit $n \to 1$, we can identify this with the correlation function of $p$
  % operators of type $\beta$ and $q=n-p$ operators of type $\alpha$:
  % \begin{equation}
  %   \label{fullcorr}C_{q,p}\equiv \langle \mathcal{O}_{\alpha}^{(q)}\mathcal{O}_{\beta}
  %   ^{(p)}\rangle
  % \end{equation}
  % As was pointed out, in the case of free bosons, the order of operators will not
  % matter, and the commutator terms vanish.
  % For operators $h_\alpha=\bar{h}_\alpha$, and $h_\beta=\bar{h}_\beta$, we have
  % \begin{eqnarray}
  %   \tilde{\mO}_{\alpha|\beta}^{(n-1)}=\sum_{p=0}^{n-1}y^{n-1-p}(1-y)^p\lb \frac{\Lambda^2}{n^2}\rb^{2p(h_\beta-h_\alpha)}\sum_{\substack{w\in\{\alpha,\beta\}^{n-1}\\ N_\beta(w)=p}}
  %   \prod_{k=1}^{n-1} \frac{V_{w_k}(z_k^+)V_{-w_k}(z_k^-)}{\braket{V_{w_k}(z^+)V_{-w_k}(z^-)}}
  % \end{eqnarray}

  % \subsection{Small $x$ limit}
\subsection{\texorpdfstring{Small $x$ limit}{Small x limit}}
  
  The analytic correlator we need to compute to find the relative entropy is
  \begin{align}
    \label{eq:F_n_mixed}\mathcal{F}_{n} & =\frac{\,1+(2\sin(\pi x/n))^{2\Delta_L}f(n)+\cdots\,}{\displaystyle \sum_{p=0}^{\,n-1}y^{\,n-1-p}((1-y)D_n(x))^{p}\sum_{\substack{w\in\{\alpha,\beta\}^{\,n-1}\\ N_\beta(w)=p}}\left[\,1+(2\sin(\pi x/n))^{2\Delta_L}g_{w}(n)+\cdots\,\right] }\\
    \label{Dn_mixed}
    D_{n}(x)&=\left(\frac{n\sin(\pi x/n)}{\sin(\pi x)}\right)^{4(h_\Psi-h_\Phi)}\simeq 1+\frac{2(h_\Psi-h_\Phi)(\pi x)^2}{3}\lb 1-\frac{1}{n^2}\rb+\cdots 
    ,\qquad D_{1}(x)=1,
    % D_{n}(x)&=\left(\frac{n\sin(\pi x/n)}{\sin(\pi x)}\right)^{4(h_\alpha-h_\beta)}\simeq 1+\frac{2p(h_\alpha-h_\beta)(\pi x)^2}{3}\lb 1-\frac{1}{n}\rb+\cdots 
    % ,\qquad D_{1}(x)=1,
  \end{align}
  % \begin{align}
  %   D_{n}(x)&=\left(\frac{n\sin(\pi x/n)}{\sin(\pi x)}\right)^{4(h_\Psi-h_\Phi)}\simeq 1+\frac{2(h_\Psi-h_\Phi)(\pi x)^2}{3}\lb 1-\frac{1}{n^2}\rb+\cdots 
  %   ,\qquad D_{1}(x)=1,
  % \end{align}
  where the leading correction $f(n)$ and $g_w(n)$ are given by the sums in (\ref{fnexpression}) and (\ref{gnw}), respectively, over OPE coefficients and light operator correlators on the replica manifold. It can be shown that the phase terms also get cancelled exactly, one from the conformal factor and the other from $\zeta_n$ from each OPE. So, there is no dependency on spin, which showed up from these phases. First observation is that if $\Delta_L> 1$, the dominant contibution to $\mathcal{F}_n$ comes from the $D_n(x)$ term and the denominator becomes 
  % \begin{align}
  %   &\sum_{p=0}^{n-1} y^{n-1-p}(1-y)^p\binom{n-1}{p}\lb 1+\frac{2p(1-1/n)}{3}(h_\alpha-h_\beta)(\pi x)^2+\cdots\rb\nn\\
  %   &=1+\frac{2(n-1)^2(h_\alpha-h_\beta)(\pi x)^2(1-y)}{3n}+\cdots 
  % \end{align}
  \begin{align}
    &\sum_{p=0}^{n-1} y^{n-1-p}(1-y)^p\binom{n-1}{p}\lb 1+\frac{2p(1-1/n^2)}{3}(h_\Psi-h_\Phi)(\pi x)^2+\cdots\rb\nn\\
    &=1+\frac{2(n-1)^2(n+1)(h_\Psi-h_\Phi)(\pi x)^2(1-y)}{3n^2}+\cdots 
  \end{align}
  where we have used 
  \begin{align}
  \sum_{p=0}^{\,n-1}y^{\,n-1-p}(1-y)^{p}p = (n-1)(1-y)
  % \sum_{\substack{w\in\{\alpha,\beta\}^{\,n-1}\\ N_\beta(w)=p}}1=\binom{n-1}{p}.
  \end{align}
  % \begin{align*}
  % \JM{\sum_{\substack{w\in\{\alpha,\beta\}^{\,n-1}\\ N_\beta(w)=p}}1=\binom{n-1}{p}.}
  % \end{align*}
Taking the logarithm and the $n$-derivative at $n=1$ we find that the $x^2$ term does not contribute to the relative entropy
% \begin{align}
% \partial_n\log \mathcal{F}_n(x)=-\frac{2}{3}(h_\alpha-h_\beta)(1-y)(\pi x)^2\partial_n \lb \partial_n \frac{(n-1)^2}{n}\rb_{n\to 1}+O(x^{2\Delta_L})=O(x^{2\Delta_L})
% \end{align}
\begin{align}
\partial_n\log \mathcal{F}_n(x)=-\frac{2}{3}(h_\Psi-h_\Phi)(1-y)(\pi x)^2\partial_n \lb  \frac{(n-1)^2 (n+1)}{n^2}\rb_{n\to 1}+O(x^{2\Delta_L})=O(x^{2\Delta_L})
\end{align}
Therefore, we can ignore the contribution of the $D_n(x)$ term in $\mathcal{F}_n$. For future convenience, we keep $D_n(x)$ in the calculation for now, because it does not play an important role in the analytic continuation discussed below, and it will help us generalize to general $x$ later. 


The expression for $f(n)$ is
  % \begin{align}\label{fnexpression}
  %   f(n) & =C_{\alpha}^{2}\sum_{k\neq k'=1}^{n}\langle \mO_{L}(e^{2\pi i k/n})\mO_{L}(e^{2\pi i k'/n})\rangle \nn    \\
  %        & =C_{\alpha}^{2}\sum_{k\neq k'=1}^{n-1}\frac{e^{-2\pi i s_L(k+k')/n}}{(2\sin(\pi |k-k'|/n))^{2\Delta_L}}\ .
  % \end{align}
  \begin{align}\label{fnexpression}
    f(n) & =\frac{C_{\Psi}^{2}}{2}\sum_{k\neq k'=0}^{n}\langle \mO_{L}(e^{2\pi i k/n})\mO_{L}(e^{2\pi i k'/n})\rangle \nn    \\
         & =\frac{C_{\Psi}^{2}}{2}\sum_{k\neq k'=0}^{n-1}\frac{e^{-2\pi i s_L(k+k')/n}}{(2\sin(\pi |k-k'|/n))^{2\Delta_L}}\ .
  \end{align}
  Rewrite the sum above by fixing the separation $m:=k-k' \ (\mathrm{mod}\ n)$. For
  each $m\in\{1,2,\dots,n-1\}$, set $k'=k-m$ (mod $n$). Then
  \[
    e^{-2\pi i s_L (k+k')/n}= e^{-2\pi i s_L (2k-m)/n}= e^{+2\pi i s_L m/n}\,e^{-4\pi
    i s_L k/n},
  \]
  and the denominator depends only on $m$. Hence the sum factorizes:
  \begin{equation}
    \sum_{m=1}^{n-1}\frac{e^{2\pi i s_L m/n}}{\bigl(2\sin\frac{\pi m}{n}\bigr)^{2\Delta_L}}
    \underbrace{\sum_{k=0}^{n-1} e^{-4\pi i s_L k/n}}_{=:G(n)}.
  \end{equation}
  The inner sum $G(n)$ is a geometric series with ratio $q=e^{-4\pi i s_L/n}$:
  \[
    G(n)=\sum_{k=0}^{n-1}q^{k}=
    \begin{cases}
      n, & q=1,     \\
      0, & q\neq 1.
    \end{cases}
  \]
  Therefore $S(n)=0$ for integer $n$ unless $q=1$, i.e.
  \[
    e^{-4\pi i s_L/n}=1 \quad\Longleftrightarrow\quad \frac{4s_{L}}{n}\in\mathbb{Z}
    .
  \]
  Now, since $4s_{L}/n$ is not an integer, the numerator $e^{2\pi i s_L m/n}$ is
  at most a sign. For $s_L=0$, we are left with
  \begin{align}
     & S_{2}(n)=\frac{n}{2}\sum_{k=1}^{n-1}(2\sin(\pi k/n))^{-2\Delta_L},\qquad S_{2}(1)=0,
  \end{align}
  whose analytic continuation of $S_{2}(n)$ was discussed in
  \cite{calabrese2011entanglement} (note that our notation has an extra factor $2^{-2\Delta_L}$.)
  \begin{align}
    S'_{2}(1)=\frac{2^{-2\Delta}\sqrt{\pi}\Gamma(\Delta_{L}+1)}{4\Gamma(\Delta_{L}+3/2)}\ .
  \end{align}
  In a general CFT, there could be several inequivalent choices of the lightest primary
  operator $\mO_{L}$ that contribute at the same order. In this case, the leading
  correction $f(n)$ would be given by the sum over all such operators, that
  following \cite{calabrese2011entanglement} we denote by
  \begin{align}
    f(n) & =S_{2}(n)
  \end{align}
  Similarly, we have
  % \begin{align}\label{gnw}
  %   g_{w}(n)= & C_{\alpha}\sum_{k=1}^{n-1}\frac{C_{w_k}}{(2\sin(\pi k/n))^{2\Delta_L}}+\sum_{k\neq k'=1}^{n-1}\frac{C_{w_k}C_{w_{k'}}}{(2\sin(\pi |k-k'|/n))^{2\Delta_L}}, % \nn\\
  %   % &=C_\alpha C_\omega \frac{2 S_2(n)}{n}+\sum_{k\neq k'=1}^{n-1}
  %   % \frac{C_{w_k}C_{w_{k'}}}{(2\sin(\pi |k-k'|/n))^{4h_L}}
  % \end{align}
  \begin{align}\label{gnw}
    g_{w}(n)= & C_{\Psi}\sum_{k=1}^{n-1}\frac{C_{w_k}}{(2\sin(\pi k/n))^{2\Delta_L}}+\frac{1}{2}\sum_{k\neq k'=1}^{n-1}\frac{C_{w_k}C_{w_{k'}}}{(2\sin(\pi |k-k'|/n))^{2\Delta_L}}, % \nn\\
    % &=C_\alpha C_\omega \frac{2 S_2(n)}{n}+\sum_{k\neq k'=1}^{n-1}
    % \frac{C_{w_k}C_{w_{k'}}}{(2\sin(\pi |k-k'|/n))^{4h_L}}
  \end{align}

  We can simplify the denominator of $\mathcal{F}_{n}$ by noting that
  \begin{align}
    % &\sum_{\substack{w\in\{\alpha,\beta\}^{\,n-1}\\
    %                        N_\beta(w)=p}} 1=\binom{n-1}{p},\nn\\
     & \sum_{\substack{w\in\{\alpha,\beta\}^{\,n-1}\\ N_\beta(w)=p}}\left[\,1+(2\sin(\pi x/n))^{2\Delta_L}g_{w}(n)+\cdots\,\right]=\binom{n-1}{p}+(2\sin(\pi x/n))^{2\Delta_L}\bar{g}_{p}(n)+\cdots\,\nn                                   \\
     % & \bar{g}_{p}(n) =\sum_{\substack{w\in\{\alpha,\beta\}^{\,n-1}\\ N_\beta(w)=p}}g_{w}(n)=C_{\alpha}\sum_{k=1}^{n-1}\frac{A_{k}}{(2\sin(\pi k/n))^{2\Delta_L}}+\sum_{k\neq k'=1}^{n-1}\frac{B_{kk'}}{(2\sin(\pi |k-k'|/n))^{2\Delta_L}},
    &\bar{g}_{p}(n) =\sum_{\substack{w\in\{\Psi,\Phi\}^{\,n-1}\\ N_\Phi(w)=p}}g_{w}(n)=C_{\Psi}\sum_{k=1}^{n-1}\frac{A_{k}}{(2\sin(\pi k/n))^{2\Delta_L}}+\frac{1}{2}\sum_{k\neq k'=1}^{n-1}\frac{B_{kk'}}{(2\sin(\pi |k-k'|/n))^{2\Delta_L}},
  \end{align}
  % \begin{align*}
  %   \bar{g}_{p}(n) =\sum_{\substack{w\in\{\Psi,\Phi\}^{\,n-1}\\ N_\Phi(w)=p}}g_{w}(n)=C_{\Psi}\sum_{k=1}^{n-1}\frac{A_{k}}{(2\sin(\pi k/n))^{2\Delta_L}}+\frac{1}{2}\sum_{k\neq k'=1}^{n-1}\frac{B_{kk'}}{(2\sin(\pi |k-k'|/n))^{2\Delta_L}},
  % \end{align*}
  where we have exchanged the order of the sum over words and the sum over
  replica indices and defined
  \begin{align}
     % & A_{k}(p,n)=\sum_{\substack{w\in\{\alpha,\beta\}^{\,n-1}\\ N_\beta(w)=p}}C_{w_k}=\binom{n-2}{p}C_{\alpha}+\binom{n-2}{p-1}C_{\beta}\equiv A(p,n)\nn    
     % \\
      & A_{k}(p,n)=\sum_{\substack{w\in\{\Psi,\Phi\}^{\,n-1}\\ N_\Phi(w)=p}}C_{w_k}=\binom{n-2}{p}C_{\Psi}+\binom{n-2}{p-1}C_{\Phi}\equiv A(p,n)\nn  \\
     % & B_{kk'}(p,n)=\sum_{\substack{w\in\{\alpha,\beta\}^{\,n-1}\\ N_\beta(w)=p}}C_{w_k}C_{w_{k'}}=\binom{n-3}{p}C_{\alpha}^{2}+2\binom{n-3}{p-1}C_{\alpha}C_{\beta}+\binom{n-3}{p-2}C_{\beta}^{2}\equiv B(p,n),  \nn\\
     & B_{kk'}(p,n)=\sum_{\substack{w\in\{\Psi,\Phi\}^{\,n-1}\\ N_\Phi(w)=p}}C_{w_k}C_{w_{k'}}=\binom{n-3}{p}C_{\Psi}^{2}+2\binom{n-3}{p-1}C_{\Psi}C_{\Phi}+\binom{n-3}{p-2}C_{\Phi}^{2}\equiv B(p,n),
    % &=\frac{(n-1-p)(n-2-p)}{(n-1)(n-2)}C_{\alpha}^2+\frac{2p(n-1-p)}{(n-1)(n-2)}C_{\alpha}C_{\beta}+\frac{p(p-1)}{(n-1)(n-2)}C_{\beta}^2\nn
  \end{align}
  where we performed the sums over words for fixed $k\neq k'$ and found that the
  result is independent of $k$ and $k'$. Therefore, we obtain
  % \begin{align}
  %   g_{p}(n) & =C_{\alpha}A(p,n)\sum_{k=1}^{n-1}\frac{1}{(2\sin(\pi k/n))^{2\Delta_L}}+B(p,n)\sum_{k\neq k'=1}^{n-1}\frac{1}{(2\sin(\pi |k-k'|/n))^{2\Delta_L}}\nn      \\
  %            & =(C_{\alpha}A(p,n)-B(p,n))\sum_{k=1}^{n-1}\frac{1}{(2\sin(\pi k/n))^{2\Delta_L}}+B(p,n)\sum_{k\neq k'=0}^{n-1}\frac{1}{(2\sin(\pi |k-k'|/n))^{2\Delta_L}}\nn \\
  %            & =\frac{2S_{2}(n)}{n}(C_{\alpha}A(p,n)-B(p,n))+B(p,n)S_{2}(n)                                                                                         % \\
  %   %& =A(p,n)\frac{2S_{2}(n)}{n}+B(p,n)S_{2}(n)
  % \end{align}
  \begin{align}
    g_{p}(n) & =C_{\Psi}A(p,n)\sum_{k=1}^{n-1}\frac{1}{(2\sin(\pi k/n))^{2\Delta_L}}+\frac{B(p,n)}{2}\sum_{k\neq k'=1}^{n-1}\frac{1}{(2\sin(\pi |k-k'|/n))^{2\Delta_L}}\nn      \\
             & =(C_{\Psi}A(p,n)-B(p,n))\sum_{k=1}^{n-1}\frac{1}{(2\sin(\pi k/n))^{2\Delta_L}}+\frac{B(p,n)}{2}\sum_{k\neq k'=0}^{n-1}\frac{1}{(2\sin(\pi |k-k'|/n))^{2\Delta_L}}\nn \\
             & =\frac{2S_{2}(n)}{n}(C_{\Psi}A(p,n)-B(p,n))+B(p,n)S_{2}(n)                                                                                         % \\
    %& =A(p,n)\frac{2S_{2}(n)}{n}+B(p,n)S_{2}(n)
  \end{align}
  and as a result we can write
  \begin{align}
     & \left[\,1+(2\sin(\pi x/n))^{2\Delta_L}g_{w}(n)+\cdots\,\right]=\binom{n-1}{p}+(2\sin(\pi x/n))^{2\Delta_L}g_{p}(n)+\cdots\,
  \end{align}

We note the following identities
  \begin{align}
     & Z(n)\equiv \sum_{p=0}^{\,n-1}y^{\,n-1-p}((1-y)D_{n})^{p}\binom{n-1}{p}= (y+(1-y) D_{n})^{n-1}                                 \\
     % & A_{\alpha}(n)\equiv \sum_{p=0}^{\,n-1}y^{\,n-1-p}((1-y)D_{n})^{p}\binom{n-2}{p}=y(y+(1-y) D_{n})^{n-2}                        \\
     & A_{\Psi}(n)\equiv \sum_{p=0}^{\,n-1}y^{\,n-1-p}((1-y)D_{n})^{p}\binom{n-2}{p}=y(y+(1-y) D_{n})^{n-2} \                       \\
     % & A_{\beta}(n)\equiv \sum_{p=0}^{\,n-1}y^{\,n-1-p}((1-y)D_{n})^{p}\binom{n-2}{p-1}=(1-y)D_{n}(y+(1-y) D_{n})^{n-2}              \\
     & A_{\Phi}(n)\equiv \sum_{p=0}^{\,n-1}y^{\,n-1-p}((1-y)D_{n})^{p}\binom{n-2}{p-1}=(1-y)D_{n}(y+(1-y) D_{n})^{n-2}              \\
     % & B_{\alpha\alpha}(n)\equiv\sum_{p=0}^{\,n-1}y^{\,n-1-p}((1-y)D_{n})^{p}\binom{n-3}{p}= y^{2}(y+(1-y)D_{n})^{n-3}               \\
     & B_{\Psi\Psi}(n)\equiv\sum_{p=0}^{\,n-1}y^{\,n-1-p}((1-y)D_{n})^{p}\binom{n-3}{p}= y^{2}(y+(1-y)D_{n})^{n-3}               \\
     % & B_{\alpha\beta}(n)\equiv\sum_{p=0}^{\,n-1}y^{\,n-1-p}((1-y)D_{n})^{p}\binom{n-3}{p-1}= y(1-y) D_{n}(y+(1-y)D_{n})^{n-3}       \\
     & B_{\Psi\Phi}(n)\equiv\sum_{p=0}^{\,n-1}y^{\,n-1-p}((1-y)D_{n})^{p}\binom{n-3}{p-1}= y(1-y) D_{n}(y+(1-y)D_{n})^{n-3}       \\
     % & B_{\beta\beta}(n)\equiv\sum_{p=0}^{\,n-1}y^{\,n-1-p}((1-y)D_{n})^{p}\binom{n-3}{p-2}= (1-y)^{2}D_{n}^{2}(y+(1-y)D_{n})^{n-3},\\
     & B_{\Phi\Phi}(n)\equiv\sum_{p=0}^{\,n-1}y^{\,n-1-p}((1-y)D_{n})^{p}\binom{n-3}{p-2}= (1-y)^{2}D_{n}^{2}(y+(1-y)D_{n})^{n-3}, 
  \end{align}
  For simplicity, we define
  \begin{align}
     & X(n)=(y+(1-y)D_{n}), \qquad X(1)=1\nn                 \\
     % & Y_{\alpha\beta}(n)=(y C_{\alpha}+(1-y)D_{n}C_{\beta}) \\
     & Y_{\Psi\Phi}(n)=(y C_{\Psi}+(1-y)D_{n}C_{\Phi}) 
  \end{align}
  so that
  % \begin{align}
  %   B(p,n) & = C_{\alpha}^{2}B_{\alpha\alpha}(n)+2C_{\alpha}C_{\beta}B_{\alpha\beta}(n)+C_{\beta}^{2}B_{\beta\beta}(n)\nn \\
  %          & =Y_{\alpha\beta}(n)^{2}X(n)^{n-3}\nn                                                                                         \\
  %   A(p,n) & = C_{\alpha}A_{\alpha}(n)+C_{\beta}A_{\beta}(n)=Y_{\alpha\beta}(n)X(n)^{n-2}
  % \end{align}
  \begin{align}
    B(p,n) & = C_{\Psi}^{2}B_{\Psi\Psi}(n)+2C_{\Psi}C_{\Phi}B_{\Psi\Phi}(n)+C_{\Phi}^{2}B_{\Phi\Phi}(n)\nn \\
           & =Y_{\Psi\Phi}(n)^{2}X(n)^{n-3}\nn                 \\
    A(p,n) & = C_{\Psi}A_{\Psi}(n)+C_{\Phi}A_{\Phi}(n)=Y_{\Psi\Phi}(n)X(n)^{n-2}
  \end{align}
  Then, in the denominator we have
  \begin{align}
     & X(n)^{n-1}\lb 1+\lb\frac{2\pi x}{n}\rb^{2\Delta_L}\Delta(n)+\cdots \rb
  \end{align}
  where
  % \begin{align}
  %   \Delta(n) & \equiv \frac{2S_{2}(n)}{n}\frac{(C_{\alpha}A(p,n)-B(p,n))}{X(n)^{n-1}}+\frac{B(p,n)}{X(n)^{n-1}}S_{2}(n)\nn                                                          \\
  %             & =\frac{2S_{2}(n)}{n}\lb \frac{C_{\alpha}Y_{\alpha\beta}(n)}{ X(n)}-\frac{Y_{\alpha\beta}(n)^{2}}{X(n)^{2}}\rb+S_{2}(n)\lb \frac{Y_{\alpha\beta}(n)^{2}}{X(n)^{2}}\rb \\
  %             & =S_{2}(n)\lb \frac{Y_{\alpha\beta}(n)}{X(n)}\left(\frac{2C_{\alpha}}{n}\right)+\frac{Y_{\alpha\beta}^{2}(n)}{X(n)^{2}}\lb 1-\frac{2}{n}\rb\rb                        %& g_{1}(n)=\frac{2(y C_{\alpha}+(1-y)D_{n}C_{\beta})}{n(y+(1-y) D_{n})}+\frac{(yC_{\alpha}+(1-y)D_{n}C_{\beta})^{2}}{(y+(1-y) D_{n})^{2}}
  % \end{align}
\begin{align}
    \Delta(n) & \equiv \frac{2S_{2}(n)}{n}\frac{(C_{\Psi}A(p,n)-B(p,n))}{X(n)^{n-1}}+\frac{B(p,n)}{X(n)^{n-1}}S_{2}(n)\nn                                                          \\
              & =\frac{2S_{2}(n)}{n}\lb \frac{C_{\Psi}Y_{\Psi\Phi}(n)}{ X(n)}-\frac{Y_{\Psi\Phi}(n)^{2}}{X(n)^{2}}\rb+S_{2}(n)\lb \frac{Y_{\Psi\Phi}(n)^{2}}{X(n)^{2}}\rb \\
              & =S_{2}(n)\lb \frac{Y_{\Psi\Phi}(n)}{X(n)}\left(\frac{2C_{\Psi}}{n}\right)+\frac{Y_{\Psi\Phi}^{2}(n)}{X(n)^{2}}\lb 1-\frac{2}{n}\rb\rb                       %& g_{1}(n)=\frac{2(y C_{\alpha}+(1-y)D_{n}C_{\beta})}{n(y+(1-y) D_{n})}+\frac{(yC_{\alpha}+(1-y)D_{n}C_{\beta})^{2}}{(y+(1-y) D_{n})^{2}}
\end{align}
  Putting this back into (\ref{eq:F_n_mixed}) we find
  % \begin{align}
  %   \log\mathcal{F}_{n}=(1-n)\log X(n)+\log\left[\frac{1+\lb \frac{2\pi x}{n}\rb^{2\Delta_L}S_{2}(n)\lb C_{\alpha}^{2}+\cdots\rb}{1+\lb \frac{2\pi x}{n}\rb^{2\Delta_L}\Delta(n)+\cdots}\right]
  % \end{align}
    \begin{align}
    \log\mathcal{F}_{n}=(1-n)\log X(n)+\log\left[\frac{1+\lb \frac{2\pi x}{n}\rb^{2\Delta_L}S_{2}(n)\lb C_{\Psi}^{2}+\cdots\rb}{1+\lb \frac{2\pi x}{n}\rb^{2\Delta_L}\Delta(n)+\cdots}\right]
  \end{align}
  We first note that taking the $n$-derivative of the first factor in $\mathcal{F}
  _{n}$ and sending $n\to 1$ vanishes
  \begin{align}
    \partial_{n}\left[(1-n)\log(y+(1-y)D_{n})\right]\Big|_{n=1}=0
  \end{align}
  and we are left with
  % \begin{align}
  %   S(\rho_{\alpha}\|y\rho_{\alpha}+(1-y)\rho_{\beta}) & =\partial_{n}\left[\lb \frac{2\pi x}{n}\rb^{2\Delta_L}(S_{2}(n) C_{\alpha}^{2}-\Delta(n))+\cdots\right]_{n=1}
  % \end{align}
  \begin{align}
    S(\rho_\Psi\|y\rho_{\Psi}+(1-y)\rho_{\Phi}) & =\partial_{n}\left[\lb \frac{2\pi x}{n}\rb^{2\Delta_L}(S_{2}(n) C_\Psi^{2}-\Delta(n))+\cdots\right]_{n=1}
  \end{align}
Note that this is consistent with our previous observation that $D_n(x)$ does not contribute to the relative entropy at order $x^2$. However, the analytic continuation we achieved shows that the higher order terms of $D_n(x)$ do not contribute either.

  Since $S_{2}(1)=0$ and $X(1)=1$, we have
  % \begin{align}
  %   S(\rho_{\alpha}\|y\rho_{\alpha}+(1-y)\rho_{\beta}) & = (2\pi x)^{2\Delta_L}S_{2}'(1)(C_{\alpha}^{2}-2C_{\alpha}Y_{\alpha\beta}(1)+Y_{\alpha\beta}(1)^{2})+\cdots\nn                \\
  %                                                          & =(\pi x)^{2\Delta_L}\frac{\sqrt{\pi}\Gamma(\Delta_{L}+1)}{4\Gamma(\Delta_{L}+3/2)}(1-y)^{2}(C_{\alpha}-C_{\beta})^{2}+\cdots
  % \end{align}

   \begin{align}
  S(\rho_{\Psi}\|y\rho_{\Psi}+(1-y)\rho_{\Phi}) & = (2\pi x)^{2\Delta_L}S_{2}'(1)(C_{\Psi}^{2}-2C_{\Psi}Y_{\Psi\Phi}(1)+Y_{\Psi\Phi}(1)^{2})+\cdots\nn                \\
   & =(\pi x)^{2\Delta_L}\frac{\sqrt{\pi}\Gamma(\Delta_{L}+1)}{4\Gamma(\Delta_{L}+3/2)}(1-y)^{2}(C_{\Psi}-C_{\Phi})^{2}+\cdots
  \end{align}
  
  where in the second line we have used $S_{2}(1)=0$. In the special case,
  $\alpha =\beta$ or $y\to 1$ the term vanishes as expected. In the case of vertex
  operators in free massless bosons the equations of motion are $\partial_z\partial_{\bar{z}}\phi(z,\bar{z})=0$, therefore the field splits into a chiral and anti-chiral mode $\phi(z,\bar{z})=\varphi(z)+\bar{\varphi}(\bar{z})$. This matches the result of \cite{Lashkari_2014,S_rosi_2016}. In the case $y\to 0$ we have
 \begin{align}
  S(\rho_\alpha\|\rho_\beta)=(C_\alpha-C_\beta)^2\frac{\pi^2 x^2}{3}+\cdots 
 \end{align}

  

  % \NL{HERE}
  % \begin{align}
  %    & \sum_{p=0}^{n-1}y^{n-1-p}((1-y)A)^{p}S(p)={}\nn                                                                                              \\
  %    & \frac{(n-1)(n-2)}{(n-1)(n-2)}C_{\alpha}^{2}y^{2}(y+(1-y)A)^{n-3}+\frac{2(n-1)(n-1)}{(n-1)(n-2)}C_{\alpha}C_{\beta}A y^{2}(y+(1-y)A)^{n-3}\nn \\
  %    & =(y+(1-y)A)^{n-3}(y C_{\alpha}+(1-y)A C_{\beta})^{2}
  % \end{align}
  % and the denominator of $\mathcal{F}_{n}$ becomes
  % \begin{align}
  %    & (y+(1-y)A)^{n-1}\left[1+(2\sin(\pi x/n))^{4h_L}h(n)+\cdots\right] \nn \\
  %    & h(n)=\lb\frac{y C_{\alpha}+(1-y)A_{n}C_{\beta}}{y+(1-y)A_{n}}\rb^{2}
  % \end{align}
  % Therefore, $\mathcal{F}_{n}$ in the small $x$ expansion becomes
  % \begin{align}
  %   \mathcal{F}_{n}= & \lb y+(1-y)A_{n}\rb^{1-n}\lb 1+(2\sin(\pi x/n))^{4h_L}S_{2}(n)(C_{\alpha}^{2}-h(n))+\cdots\rb\nn                                                      \\
  %   =                & \lb y+(1-y)A_{n}\rb^{1-n}\lb 1+F^{(1)}_{n}+\cdots\rb\nn                                                                                               \\
  %                    & F_{n}^{(1)}=(2\sin(\pi x/n))^{4h_L}S_{2}(n)\frac{A_{n}(C_{\alpha}-C_{\beta})(1-y)(A_{n}(C_{\alpha}+C_{\beta})(1-y)+2yC_{\alpha})}{(y+(1-y)A_{n})^{2}}
  % \end{align}

  % Taking the logarithm and $n$ derivative at $n\to 1$ of the first factor in $\mathcal{F}
  % _{n}$ vanishes
  % \begin{align}
  %   \partial_{n}\lb (1-n)\log(y+(1-y)A_{n})\rb_{n\to 1}=-\log(y+(1-y)A_{1})=0\ .
  % \end{align}
  % Therefore, the leading contribution to the relative entropy in the small $x$
  % limit comes from the second factor and since $S_{2}(1)=0$ we find
  % \begin{align}
  %   \partial_{n}\log\mathcal{F}_{n}\big|_{n\to 1}=\partial_{n}F_{n}^{(1)}\big|_{n\to 1}=(2\sin(\pi x))^{4h_L}S_{2}'(1)(1-y)(C_{\alpha}-C_{\beta})(C_{\alpha}+C_{\beta}+y(C_{\alpha}-C_{\beta}))\ .
  % \end{align}
  % Putting it all together, we find the leading small $x$ behavior of the relative
  % entropy
  % \begin{align}
  %   S(\rho_{\alpha}|y\rho_{\alpha}+(1-y)\rho_{\beta})=(2\sin(\pi x))^{4h_L}\frac{\sqrt{\pi}\Gamma(2h_{L}+1)}{4\Gamma(2h_{L}+3/2)}(1-y)(C_{\alpha}-C_{\beta})(C_{\alpha}+C_{\beta}+y(C_{\alpha}-C_{\beta}))+\cdots
  % \end{align}
  % For the special case $y=1/2$ and $\alpha=0$ we find
  % \begin{align}
  %   S(\rho_{0}|(\rho_{0}+\rho_{\beta})/2)=(2\sin(\pi x))^{4h_L}\frac{\sqrt{\pi}\Gamma(2h_{L}+1)}{8\Gamma(2h_{L}+3/2)}C_{\beta}^{2}+\cdots
  % \end{align}
  % In the special case of vertex operators in the massless free boson, the lowest
  % dimension primary is $\partial\phi$ and we find
  % \begin{align}
  %   S(\rho_{0}|(\rho_{0}+\rho_{\beta})/2)=(2\sin(\pi x))^{\beta^2}\frac{\sqrt{\pi}\Gamma(\beta^{2}+1)}{8\Gamma(\beta^{2}+3/2)}+\cdots
  % \end{align}



%\subsection{Limit $y\to 0$}

In the limit $y\to 0$, this is the replica trick already discussed in \cite{Lashkari_2014} that computes the relative entropy $S(\rho_\alpha||\rho_\beta)$. In this case, in the denominator $\tilde{O}_{\alpha|\beta;y}^{(n-1)}$ picks up only the $p=(n-1)$ term with only one word $w=\{\beta,\beta,\cdots, \beta\}$
\begin{align}
  \tilde{O}_\beta^{(n-1)} = \prod_{k=1}^{n-1}\frac{V_{\beta}(z_{k,n}^{+})V_{-\beta}(z_{k,n}^{-})}{\braket{V_{\beta}(z^+_{1,1})V_{-\beta}(z^-_{1,1})}}.
\end{align}
The denominator computes 
\begin{align}
\braket{V_\alpha|(H_0-H_\beta)|V_\alpha}-\tr((\rho_\alpha-\rho_0)H_0)
\end{align}
Similarly, repeating our calculation for the case where $\Delta_L>2$, in the lowest order we get the contribution from the two-point function of stress tensors. In full generality, wef obtain, 
% \begin{align}
% \boxed{ S(\rho_\Psi\|y\rho_\Psi + (1-y)\rho_\Phi)=\begin{cases} \frac{\sqrt{\pi}\Gamma(\Delta_L+1)}{4\Gamma(\Delta_L+3/2)}(C^L_{\Psi\Psi^*}-C^L_{\Phi\Phi^*})^2 (1-y)^2(\pi x)^{2\Delta_L}+\cdots&\text{if }\Delta_L\leq 2\\ \frac{8}{15c}\lb (h_\Psi - h_\Phi)^2 + (\bar h_\Psi - \bar h_\Phi)^2 \rb(1-y)^2(\pi x)^4 + \cdots &\text{if }\Delta_L>2 \end{cases}} \end{align}
\begin{align}\label{smallxMixture}
\boxed{
S(\rho_\Psi\|y\rho_\Psi+(1-y)\rho_\Phi)
=
\begin{cases}
\begin{aligned}
&\frac{\sqrt{\pi}\Gamma(\Delta_L+1)}
{4\Gamma(\Delta_L+\frac32)}
\sum_L
(C^L_{\Psi\Psi^*}-C^L_{\Phi\Phi^*})^2
(1-y)^2(\pi x)^{2\Delta_L}
+\cdots
\end{aligned}
&
\text{if }\Delta_L\le 2,
\\[2ex]
\begin{aligned}
&\frac{8}{15c}
\Big[
(h_\Psi-h_\Phi)^2 
% &\qquad\qquad
+(\bar h_\Psi-\bar h_\Phi)^2
\Big]
(1-y)^2(\pi x)^4
+\cdots
\end{aligned}
&
\text{if }\Delta_L>2 .
\end{cases}
}
\end{align}
where the sum over $L$ is over all operator of lowest dimension $\Delta_L$ in the OPE of $\Psi\Psi^*$ and $\Phi \Phi^*$.
An interesting feature of this equation is that it is symmetric under switching $\Psi$ and $\Phi$:
\begin{align}
S(\rho_\Psi\|y\rho_\Psi+(1-y)\rho_\Phi)=S(\rho_\Phi\|y\rho_\Psi+(1-y)\rho_\Phi)
\end{align}
and if we define $\bar{\rho}_y=y\rho_\Psi+(1-y)\rho_\Phi$ then
\begin{align}
    y S(\rho_\Psi\|\bar{\rho}_y)=(1-y) S(\rho_\Phi\|\bar{\rho}_y)\ .
\end{align}
\subsection{}
This can be generalized further by computing the relative entropy of a mixture with respect to another mixture, namely $S(\bar{\rho}_{y_1} \| \bar{\rho}_{y_2} )$, where $\bar{\rho}_{y_1}=y_1\rho_\Psi+(1-y_1)\rho_\Phi$ and $\bar{\rho}_{y_2}=y_2\rho_\Psi+(1-y_2)\rho_\Phi$. 
We want to compute it using replica trick, by computing the following,
\begin{align}
    S(\bar{\rho}_{y_1} \| \bar{\rho}_{y_2} ) = \p_n \log \left[\frac{\text{tr}(\bar{\rho}_{y_1}^n)}{\text{tr}(\bar{\rho}_{y_1}\bar{\rho}_{y_2}^{n-1})}\right]
\end{align}
For the numerator, we want to evaluate,
\begin{align}
    B_{kk'}(p,n) = \sum_{\substack{w\in\{\Psi,\Phi\}^{\,n}\\ N_\Phi(w)=p}} C_\Psi^2 \binom{n-2}{p} + 2C_\Psi C_\Phi \binom{n-2}{p-1} +C_\Phi^2\binom{n-2}{p-2}
\end{align}
The term other than the identity piece in the numerator evaluates to the following,
\begin{align}
    & \sum_{p=0}^{n} y_1^{n-p} \lb (1-y_1)D_n\rb^p\lb C_\Psi^2 \binom{n-2}{p} + 2C_\Psi C_\Phi \binom{n-2}{p-1} +C_\Phi^2\binom{n-2}{p-2} \rb \\
    &= \lb y_1+ (1-y_1 )D_n\rb^{n-2} \lb C_\Psi^2 y_1^2 + 2C_\Psi C_\Phi y_1(1-y_1)D_n + C_\Phi^2 ((1-y_1)D_n)^2  \rb \\
    &= X_1(n)^{n-2} \lb C_\Psi y_1 + C_\Phi(1-y_1)D_n \rb^2 \\
    &= X_1(n)^{n-2}Y_1(n)^{2}
\end{align}
where we defined
\begin{align}
    X_1(n) &= y_1 + (1-y_1)D_n \\
    Y_1(n) &=  C_\Psi y_1 + C_\Phi(1-y_1)D_n 
\end{align}
Finally the numerator is given by,
\begin{align}
    \text{tr}(\bar{\rho}_{y_1}^n) = X_1^n\lb1 + \lb\frac{2\pi x}{n}\rb^{2\Delta_L} S_2(n) \frac{Y_1(n)^2}{X_1(n)^2} + \cdots \rb
\end{align}
Now for the denominator, we want to evaluate the following, 
\begin{align}
    \text{tr}\lb \bar{\rho}_{y_1}\bar{\rho}_{y_2}^{n-1}  \rb
\end{align}
Again, we can define the following for the denominator,
\begin{align}
    X_2(n) &= y_2 + (1-y_2)D_n \\
    Y_2(n) &=  C_\Psi y_2 + C_\Phi(1-y_2)D_n 
\end{align}
The identity piece comes from precisely from $X_1 X_2^{n-1}$. The piece coming from the two point function evaluates to (following the same method we used earlier),
\begin{align}
    &X_2(n)^{n-2}\lb Y_1(n)Y_2 (n)\frac{2S_2(n)}{n} + \frac{1}{2}Y_2(n)^2 \lb S_2(n) - \frac{2S_2(n)}{n}\rb \rb \\
\end{align}
For simplicity of notation, we use the notation,
\begin{align}
    X_1(n) &\equiv X_1 ,\qquad Y_1(n) \equiv Y_1 ,\\
    X_2(n)&\equiv X_2, \qquad Y_2(n) \equiv Y_2.
\end{align}
The denominator is then given by,
\begin{align}
    \text{tr}\lb \bar{\rho}_{y_1}\bar{\rho}_{y_2}^{n-1}  \rb &= X_1 X_2^{n-1} \Bigg[ 1+  \lb\frac{2\pi x}{n}\rb^{2\Delta_L} \Bigg(\frac{Y_1Y_2}{X_1X_2}\lb\frac{2S_2(n)}{n}\rb  + \lb\frac{Y_2}{X_2}\rb^2 \lb S_2(n)-\frac{2S_2(n)}{n}\rb\Bigg)+\cdots\Bigg]
\end{align}
We finally have,
\begin{align}
    S(\bar{\rho}_{y_1} \| \bar{\rho}_{y_2} ) &= \p_n \log \Bigg[ \frac{X_1^n\lb1 + \lb\frac{2\pi x}{n}\rb^{2\Delta_L} S_2(n) \frac{Y_1^2}{X_1^2} + \cdots \rb}{X_1 X_2^{n-1} \Bigg[ 1+  \lb\frac{2\pi x}{n}\rb^{2\Delta_L} \Bigg(\frac{Y_1Y_2}{X_1X_2}\lb\frac{2S_2(n)}{n}\rb + \lb\frac{Y_2}{X_2}\rb^2 \lb S_2(n)-\frac{2S_2(n)}{n}\rb\Bigg)+\cdots\Bigg]}\Bigg] \\
    &= \p_n \log \lb \frac{X_1}{X_2}\rb^{n-1} \lb 1+ \lb\frac{2\pi x}{n}\rb^{2\Delta_L} S_2'(1) \lb Y_1^2 -2Y_1Y_2 +Y_2^2\rb\rb
\end{align}
where we used $X_1(1) = X_2(1)= 1$. Using the definitions of $Y_1,Y_2$, we finally obtain the relative entropy of mixtures of two excited states,
\begin{align}
    S(\bar{\rho}_{y_1} \| \bar{\rho}_{y_2} ) &= \frac{\sqrt{\pi} \Gamma(\Delta_L + 1)}{4\Gamma(\Delta_L + 3/2)}\sum_L (C_{\Psi \Psi^*}^L - C_{\Phi \Phi^*}^L)^2(y_1-y_2)^2 (\pi x)^{2\Delta_L} + \cdots \qquad \text{if }\Delta_L \leq 2.
\end{align}
Generalizing our results for $\Delta_L \geq 2$, we thus have,
\begin{align}\label{smallxMixture_wrt_mixture}
\boxed{
S(\bar{\rho}_{y_1} \| \bar{\rho}_{y_2} )
=
\begin{cases}
\begin{aligned}
&\frac{\sqrt{\pi}\Gamma(\Delta_L+1)}
{4\Gamma(\Delta_L+\frac32)}
\sum_L
(C^L_{\Psi\Psi^*}-C^L_{\Phi\Phi^*})^2
(y_1-y_2)^2(\pi x)^{2\Delta_L}
+\cdots
\end{aligned}
&
\text{if }\Delta_L\le 2,
\\[2ex]
\begin{aligned}
&\frac{8}{15c}
\Big[
(h_\Psi-h_\Phi)^2 
% &\qquad\qquad
+(\bar h_\Psi-\bar h_\Phi)^2
\Big]
(y_1-y_2)^2(\pi x)^4
+\cdots
\end{aligned}
&
\text{if }\Delta_L>2 .
\end{cases}
}
\end{align}

\NL{Understand why this is symmetric under $(y_1,y_2)\to (y_1+c,y_2+c)$. The fact that it is proportional to $(y_1-y_2)^2$ is expected from the Fisher information form of it, but in principle it could be proportional to $(y_1-y_2)^2 f(y_1)$.}

% There are two important functions above: 1) $(1-\pi x\cot(\pi x))$ has a simple pole at $x=1$. 2) $\log(2\sin(\pi x))$ has a branch point at $x=1$. 

% Near $x=1$ and $n$ a positive integer we have
% \begin{align}
%   D_n(1-x)\simeq \lb \frac{n\sin(\pi/n)}{\pi (1-x)}\rb^{-4p(h_\alpha-h_\beta)}\lb 1-\frac{4(h_\alpha-h_\beta)p}{n}\pi\cot(\pi/n)(1-x)+O((1-x)^2)\rb
% \end{align}
% However, this expansion does not help us much with the analytic continuation because the leading term still has the form of function of $x$ to power $p$. Therefore, we keep $D_n(1-x)$ as it is.



  % \subsection{Small $1-x$ limit}
\subsection{\texorpdfstring{Small $1-x$ limit}{Small 1-x limit}}

Here, we expand the correlator in (\ref{CorrFn_mixed}) in the limit of small $\ep=1-x$ to compute the relative entropy of mixtures in the limit $x\to 1^-$. In this limit, we have to expand in a different OPE channel (at $z_k^+$ and $z_{k+1}^-$) where operators from sheet $k$ are close to operators on sheet $k+1$. The correlators with $n=1$ (normalizations of density matrices) are symmetric under $x\to 1-x$. For the $n$-dependent correlators, at small $1-x$ the OPE expansion of the numerator remains the same as small $x$ with $x\to 1-x$. However, the denominator changes significantly because for $p>0$ there are terms that involve the OPE of operators $\Psi$ and $\Phi^*$ of weights $h_\Psi$ and $h_{\Phi}$ at $z_k^+$ and $z_{k+1}^-$. It can be shown that the phase terms also get cancelled exactly, one from the conformal factor and the other from $\zeta_n$ from each OPE. So, there is no dependency on spin, which showed up from these phases. 
% The leading term is given by
%   \begin{align}
%     V_{w_k}(z_{k}^{+})V_{-w_{k+1}}(z_{k+1}^{-})&=(2\sin(\pi (1-x)/n)^{-2(h_{\omega_k}+h_{\omega_{k+1}})} \times\nn\\
%     &\lb \delta_{w_kw_{k+1}}+ C_{w_kw_{k+1}}(2\sin(\pi (1-x)/n))^{\Delta_L}\mO_L(e^{2\pi i k/n})+\cdots\rb\nn
%   \end{align}
  The key difference compared to the previous case is that previously, all OPEs involved operators of the same weight, which starts with the identity term, whereas now, for $h_\Phi\neq h_\Psi$, there is no identity term in the OPE, and the fact that the lowest dimension operator is $\mO_{L'}$ comes with a new OPE coefficient that we denote $C_{\Psi\Phi^*}$. Since the overall factors $(2\sin(\pi (1-x)/n))$ come with powers that are additive in the operator dimensions, the new expression for $\mathcal{F}_n$  takes the form
% \begin{align}
%     &\label{eq:F_n_1_minus_x_mixed}\mathcal{F}_{n} \nn \\
%     &=\lb \,1+\lb 2\sin\lb\frac{\pi (1-x)}{n}\rb\rb^{2\Delta_L}f(n)+\cdots\,\rb\times 
%     \nn\\
%     &\quad\Bigg(\sum_{p=0}^{\,n-1}y^{\,n-1-p}((1-y)D_n(1-x))^{p}\sum_{\substack{w\in\{\alpha,\beta\}^{\,n-1}\\ N_\beta(w)=p}}\Big[\,\prod_{k=0}^{n-1}\delta_{w_kw_  {k+1}}\nn\\
%     &\qquad+\lb 2\sin\lb\frac{\pi (1-x)}{n}\rb\rb^{2\Delta_L}g_{w}(n)+\cdots\,\Big]\Bigg)^{-1}
%     \nn\\
%     &=\lb\,1+\lb\frac{2\pi (1-x)}{n}\rb^{2\Delta_L}f(n)+\cdots\,\rb\times
%     \nn\\
%     &\quad\Bigg(\sum_{p=0}^{\,n-1}y^{\,n-1-p}((1-y) D_n(1-x))^{p}\sum_{\substack{w\in\{\alpha,\beta\}^{\,n-1}\\ N_\beta(w)=p}}\Big[\,\prod_{k=0}^{n-1}\delta_{w_kw_  {k+1}}\nn\\
%     &\qquad+\lb\frac{2\pi (1-x)}{n}\rb^{2\Delta_L}g_{w}(n)+\cdots\,\Big]\Bigg)^{-1}
%     %D_n(x)&=\lb\frac{n\sin(\pi x/n)}{\sin(\pi x)}\rb^{4(h_\alpha-h_\beta)}
%     \end{align}
\begin{align}\label{eq:F_n_1_minus_x_mixed}
    &\mathcal{F}_{n}
    =\lb \,1+\lb 2\sin\lb\frac{\pi \ep}{n}\rb\rb^{2\Delta_L}f(n)+\cdots\,\rb\times 
    \nn\\
    &\quad\Bigg(\sum_{p=0}^{\,n-1}y^{\,n-1-p}((1-y)D_n(1-x))^{p}\sum_{\substack{w\in\{\Psi,\Phi\}^{\,n-1}\\ N_\Phi(w)=p}}\Big[\,\prod_{k=0}^{n-1}\delta_{w_kw_  {k+1}}+\lb 2\sin\lb\frac{\pi \ep}{n}\rb\rb^{2\Delta_L}g_{w}(n)+\cdots\,\Big]\Bigg)^{-1}
    \nn\\
    &=\lb\,1+\lb\frac{2\pi \ep}{n}\rb^{2\Delta_L}f(n)+\cdots\,\rb\times
    \nn\\
    &\quad\Bigg(\sum_{p=0}^{\,n-1}y^{\,n-1-p}((1-y) D_n\ep)^{p}\sum_{\substack{w\in\{\Psi,\Phi\}^{\,n-1}\\ N_\Phi(w)=p}}\Big[\,\prod_{k=0}^{n-1}\delta_{w_kw_  {k+1}}+\lb\frac{2\pi \ep}{n}\rb^{2\Delta_L}g_{w}(n)+\cdots\,\Big]\Bigg)^{-1}
    %D_n(x)&=\lb\frac{n\sin(\pi x/n)}{\sin(\pi x)}\rb^{4(h_\alpha-h_\beta)}
    \end{align}
  where the leading correction $f(n)$ remains the same, and hence we end up with the same $S_{2}(n)$ and the constraint with the spin of the lowest primary $s_{L'}$ being a multiple of four. Whereas $g_{w}(n)$ changes to
  \begin{align}\label{gnw_1_minus_x_mixed}
    % g_{w}(n)&=\sum_{k\neq k'=0}^{n-1}\frac{C_{w_kw_{k+1}}C_{w_{k'}w_{k'+1}}}{(2\sin(\pi |k-k'|/n))^{2\Delta_L}}\prod_{\substack{i=1\\i\neq k,i\neq k'}}^n  \delta_{w_i w_{i+1}} \\ 
    g_{w}(n)&=\sum_{k\neq k'=0}^{n-1}\frac{C_{w_kw_{k+1}}C_{w_{k'}w_{k'+1}}}{(2\sin(\pi |k-k'|/n))^{2\Delta_{\alpha-\beta}}}\prod_{\substack{i=1\\i\neq k,i\neq k'}}^n  \delta_{w_i w_{i+1}} \nn
  \end{align}  
  where we have defined $w_0=\Psi=w_n$. The product $\prod_{k=0}^{n-1}\delta_{\omega_k\omega_{k+1}}$ is non-zero only if $w_k=\Psi$) for all $k$, which means $p=0$ and 
  % $w=(\alpha,\dots,\alpha)$ 
  $w=(\Psi,\dots,\Psi)$. Therefore, the leading term in the denominator is
  % \begin{align}
  %   &\sum_{p=0}^{n-1}y^{n-1-p}((1-y)D_n(1-x))^p\sum_{\substack{w\in\{\alpha,\beta\}^{n-1}\\N_\beta(w)=p}}\prod_{\substack{i=1}}^n\delta_{w_i w_{i+1}}\nn\\
  %   &=\sum_{p=0}^{n-1}y^{n-1-p}((1-y) D_n(1-x))^p\delta_{p,0}=y^{n-1}.
  % \end{align}
  \begin{align}
    &\sum_{p=0}^{n-1}y^{n-1-p}((1-y)D_n\ep)^p\sum_{\substack{w\in\{\Psi,\Phi\}^{\,n-1}\\ N_\Phi(w)=p}}\prod_{k=0}^{n-1}\delta_{w_kw_  {k+1}}\nn\\
    &=\sum_{p=0}^{n-1}y^{n-1-p}((1-y) D_n\ep)^p\delta_{p,0}=y^{n-1}.
  \end{align}
It is clear from (\ref{gnw_1_minus_x_mixed}) that since there is no identity element in the OPE of 
% $\alpha$ and $\beta$
$\Psi$ and $\Phi^*$, the only words and $k\neq k'$ configurations that can give non-vanishing contributions with two pairs OPEs $\Psi\Phi^*$ must have $p$ consecutive $\Phi$ with $k+1$ the 
% location of the first $\beta$ and $k'-1$ the location of the last $\beta$ 
% \JM{(
location of last $\Phi$ should be k', then only, the total number of $\Phi$ should match with $k' - (k+1) +1 = k'-k =  p$
% )
% }
. All other terms are higher order because they involve four or more pairs of $\Psi\Phi^*$ OPEs.
Therefore, for fixed $k$ and $k'$ only one word contributes and we find
% \begin{align}
%   \sum_{\substack{w\in \{\alpha,\beta\}^{n-1}\\N_\beta(w)=p}}C_{w_k w_{k+1}}C_{w_{k'}w_{k'+1}}\prod_{\substack{i=1\\i\neq k,i\neq k'}}^n  \delta_{w_i w_{i+1}} =C_\alpha^2\delta_{p,0}+C_{\alpha\beta}^2\delta_{p |k-k'|}
% \end{align}
\begin{align}
  \sum_{\substack{w\in \{\Psi,\Phi\}^{n-1}\\N_\Phi}(w)=p}C_{w_k w_{k+1}}C_{w_{k'}w_{k'+1}}\prod_{\substack{i=1\\i\neq k,i\neq k'}}^n  \delta_{w_i w_{i+1}} =C_\Psi^2\delta_{p,0}+C_{\Psi\Phi}^2\delta_{p |k-k'|}
\end{align}
The $p=0$ in the denominator becomes
\begin{align}
y^{n-1} (1+(2\pi (1-x)/n)^{2\Delta_L}f(n))
\end{align}
whereas for non-zero $p$ we get
\begin{align}
  % \sum_{k\neq k'}\frac{\delta_{p|k-k'|}}{(2\sin(\pi|k-k'|/n))^{2\Delta_L}}=\frac{2(n-p)}{(2\sin(\pi p/n))^{2\Delta_L}}\\
  \frac{1}{2}\sum_{k\neq k'}\frac{\delta_{p|k-k'|}}{(2\sin(\pi|k-k'|/n))^{2\Delta_{L'}}}=\frac{1}{2}\frac{2(n-p)}{(2\sin(\pi p/n))^{2\Delta_{L'}}}
\end{align}
% \JM{(For the above equation, one should take 1/2 of it, since double counting).}
and as a result we pick up the contribution
\begin{align}
  % y^{n-1-p}((1-y) D_n(1-x))^p \frac{2(n-p) C_{\alpha\beta}^2}{(2\sin(\pi p/n))^{2\Delta_L}}(2\pi (1-x)/n)^{2\Delta_L} \\
  y^{n-1-p}((1-y) D_n(1-x))^p \frac{(n-p) C_{\Psi\Phi}^2}{(2\sin(\pi p/n))^{2\Delta_{L'}}}(2\pi (1-x)/n)^{2\Delta_{L'}}
\end{align}
Note that the $p=0$ term in the denominator corresponds to the same correlation function as the one appearing in the numerator with an extra factor of $y^{n-1}$, therefore we only need to keep track of $p>0$ terms. Plugging these back in $\mathcal{F}_n$ and expanding in small $1-x$ we find
% \begin{align}
% \mathcal{F}_n
%     &=\frac{\,1+(2\pi (1-x)/n)^{2\Delta_L}f(n)+\cdots\,}{\displaystyle y^{n-1}\lb 1+(2\pi (1-x)/n)^{2\Delta_L}\lb f(n)+2C_{\alpha\beta}^2\mu(n,x)\rb\rb}\nn\\
% &=y^{1-n}\lb 1-\lb\frac{2\pi (1-x)}{n}\rb^{2\Delta_L} 2 C_{\alpha\beta}^2 \:\mu(n,x)+\cdots \rb\nn\\
% \mu(n,x)&=\sum_{p=1}^{n-1}\frac{(n-p)((y^{-1}-1)D_n(1-x))^p}{(2\sin(\pi p/n))^{2\Delta_L}}
% \end{align}
\begin{align}
\mathcal{F}_n
    &=\frac{\,1+(2\pi (1-x)/n)^{2\Delta_L}f(n)+\cdots\,}{\displaystyle y^{n-1}\lb 1+(2\pi (1-x)/n)^{2\Delta_L} f(n)+\sum_{L'}\lb 2\pi (1-x)/n\rb ^{2\Delta_{L'}} C_{\Psi\Phi}^2\mu(n,x,y;\Delta_{L'})+\cdots\rb}\nn\\
&=y^{1-n}\lb 1-\sum_{L'}\lb\frac{2\pi (1-x)}{n}\rb^{2\Delta_{L'}}  C_{\Psi\Phi}^2 \:\mu(n,x,y;\Delta_{L'})+\cdots \rb\nn\\
\mu(n,x,y;\Delta_{L'})&=\sum_{p=1}^{n-1}\frac{(n-p)((y^{-1}-1)D_n(1-x))^p}{(2\sin(\pi p/n))^{2\Delta_{L'}}} \label{eq: mu (n,x,DeltaL')}
\end{align}
where $D_n$ is defined in (\ref{Dn_mixed}). 


Note, here we took into account the correction due to the two point functions of the higher order primaries arising from the mixing OPE of $\Psi \Phi^*$ and $\Psi^* \Phi$. They are sub-leading corrections, however, we should keep the contribution from every primary from the mixing OPE which will have the same dimension that of the sub-leading correction of the numerator. For example, in the coherent states, in order to produce the exact result, we should take two operators in the OPE of $V_\alpha$ and $V_{-\beta}$, which are $V_{\alpha - \beta}$ and $C_{\alpha\beta} V_{\alpha-\beta}J$. This will reproduce our previous result of coherent states in eq.~[\ref{eq: coherent states rel entr}] for $y\to 0$ limit. 

The reader should also note that, in our previous calculation of relative entropy $S(\rho_\Psi \| \rho_\Phi)$, we were essentially computing the $y \to 0$ limit of this relative entropy of mixture calculation. The leading divergent piece came from the fact that we had just $p=n-2$ $\Phi \Phi^*$ OPE insertions, and the mixing OPEs of $\Psi \Phi^*$ and $\Phi \Psi^*$. That gave us a term in the denominator of $\mathcal{F}_n$,
\begin{align}
    \frac{1}{(2\sin (\pi/n))^{2\Delta_{L'}}}
\end{align}
which is clearly divergent when we compute the relative entropy. However, in the mixture, we have a summation over $p$ in eq.~[\ref{eq: mu (n,x,DeltaL')}], where $p$ is the number of pairs of $\Phi \Phi^*$ insertions. We expect this summation to regulate our final answer, as we saw previously in the case where 
\begin{align}
    S_2(n,\Delta) = \sum_{k=1}^{n-1}\frac{1}{(2\sin(\pi k/n))^{2\Delta}}
\end{align}
in the limit $n\to 1$ vanishes, however, the individual terms diverge. 

So, for a general mixture, we obtain,
\begin{align}
    S(\rho_\Psi \| y\rho_\Psi + (1-y)\rho_\Phi) &= -\log y - C_{\Psi\Phi}^2 \lb 2\pi (1-x) \rb^{2\Delta_{L'}} \mu'(1,x,y;\Delta_{L'}) + \cdots
\end{align}
For simplicity, assume $y=1/2$ so that we can expand in small $\ep=1-x$ as
% \begin{align}
%   \mu(n,\ep)=\sum_{p=1}^{n-1}\frac{(n-p)(1+p\pi^2\ep^2(n^2-1)/(6n^2)+O(\ep^4))}{(2\sin(\pi p/n))^{2\Delta_L}}=\mu_0(n)+\mu_1(n)\ep^2+\cdots.
% \end{align}
\begin{align}
  \mu(n,\ep,\Delta_{L'})=\sum_{p=1}^{n-1}\frac{(n-p)(1+2p\pi^2\ep^2(n^2-1)(h_\Psi - h_\Phi)/(3n^2)+O(\ep^4))}{(2\sin(\pi p/n))^{2\Delta_{L'}}}=\mu_0(n,\Delta_{L'})+\mu_1(n,\Delta_{L'})\ep^2+\cdots.
\end{align}
We note that since the denominator of the summand is symmetric under $p\to n-p$ we have
\begin{align}
  \mu_0(n,\Delta_{L})&=\sum_{p=1}^{n-1}\frac{n-p}{(2\sin(\pi p/n))^{2\Delta_L}}=\frac{1}{2}\sum_{p=1}^{n-1}\frac{p+(n-p)}{(2\sin(\pi p/n))^{2\Delta_L}}\nn\\
  &=\frac{n}{2}\sum_{p=1}^{n-1}\frac{1}{(2\sin(\pi p/n))^{2\Delta_L}}=S_2(n,\Delta_{L})\ .
\end{align}
We find
% \begin{align}
%   S(\rho_\alpha\|(\rho_\alpha+\rho_\beta)/2))&=\log 2-(2\pi \ep)^{2\Delta_L} 2 C_{\alpha\beta}^2 (\mu'_0(1)+\mu'_1(1)\ep^2+\cdots)\nn\\
%   &=\log 2-(\pi \ep)^{2\Delta_L} 2 C_{\alpha\beta}^2 \frac{\sqrt{\pi}\Gamma(\Delta_L+1)}{4\Gamma(\Delta_L+3/2)} +\cdots
% \end{align}
\begin{align}
 S(\rho_\Psi\|(\rho_\Psi+\rho_\Phi)/2))&=\log 2-\sum_{L'}(2\pi \ep)^{2\Delta_{L'}}  C_{\Psi \Phi}^2(\mu'_0(1)+\mu'_1(1)\ep^2+\cdots)
\end{align} 
Therefore,
\begin{align}
    \boxed{S(\rho_\Psi\|(\rho_\Psi+\rho_\Phi)/2))=\log 2-\sum_{L'}(\pi \ep)^{2\Delta_{L'}} C_{\Psi \Phi}^2 \frac{\sqrt{\pi}\Gamma(\Delta_{L'}+1)}{4\Gamma(\Delta_{L'}+3/2)} +\cdots}
\end{align}
For the next order term, we find
\begin{align}
    \mu_1(n) = \sum_{p=1}^{n-1}2\pi^2 \lb n-1/n\rb  (h_\Psi - h_\Phi) \frac{p}{(2\sin(\pi p/n))^{2\Delta_{L'}}} - \sum_{p=1}^{n-1}2\pi^2 (1-1/n^2) (h_\Psi- h_\Phi) \frac{p^2}{(2\sin(\pi p/n))^{2\Delta_{L'}}} 
\end{align}
The first sum can be carried out again using the symmetry $\sin(x) = \sin (\pi - x)$, and that would give $S_2(n)$. Here the derivative at $n=1$ from that term would contribute $S_2'(1)$.

The analytic continuation of the second series $\mu_1'(n)$ is given in the appendix,
\begin{align}
    \mu_1'(n)\big|_{n=1} &= 2^{-2\Delta_{L'}+1}\pi^2(1-1/n^2)(h_\Psi - h_\Phi) \Bigg(\frac{\sqrt{\pi}\Gamma(\Delta_{L'}+1)}{4\Gamma(\Delta_{L'}+3/2)} \nn  \\
    &\qquad\qquad
    % \qquad\qquad \qquad
    % \frac{2^{2\Delta_{L'}+1}\sin(\pi \Delta_{L'})\Gamma(\Delta_{L'})\Gamma(-2\Delta_{L'})}{\pi \Gamma(-\Delta_{L'})}
    -\frac{ B(\Delta_{L'},-2\Delta_{L'})\cos(\pi \Delta_{L'})}{4}\left[ -\frac{12}{\Delta_{L'}} + \frac{2\Delta_{L'} }{2\Delta_{L'} + 1} \lb \pi^2 -6\psi_1(\Delta_{L'}) + \frac{6}{\Delta_{L'}^2} \rb  \right]
    \Bigg) \\
    & = 0.
\end{align}


As a special case, we consider the following,
\begin{align}
    S(\rho_0 \| (\rho_0 + \rho_\Phi)/2) &= \p_n \log \frac{\,2^{n-1}}{\displaystyle \lb 1+\sum_{L'}\lb 2\pi (1-x)/n\rb ^{2\Delta_{L}} C_{\Psi\Phi}^2\mu_0(n,x,1/2;\Delta_{L})+\cdots\rb}\\
    S(\rho_\Phi \| (\rho_0 + \rho_\Phi)/2)&= \p_n \log \frac{\,2^{n-1}\lb 1+(2\pi (1-x)/n)^{2\Delta_L}f(n)+\cdots\rb}{\displaystyle \lb 1+(2\pi (1-x)/n)^{2\Delta_L}f(n)+\sum_{L'}\lb 2\pi (1-x)/n\rb ^{2\Delta_{L}} C_{\Psi\Phi}^2\mu_0(n,x,1/2;\Delta_{L})+\cdots\rb}
\end{align}
Notice that, in both the cases . 2e kept only $\mu_0$ piece in the denominator, which is an approximation we made. Also, in the second line, the numerator gets cancelled exactly from the $p=n-1$ contribution from the denominator, and we will be left with $2^{n-1}$ and $\mu_0(n,x,1/2;\Delta_{L})$. Hence, we have, 
\begin{align}
     S(\rho_0 \| (\rho_0 + \rho_\Phi)/2) = S(\rho_\Phi \| (\rho_0 + \rho_\Phi)/2)
\end{align}

----

We compute the contribution of three point function $\braket{\mO_{L'} \mO_{L'}\mO_L}$ for $p>0$, which we expect to be suppressed.

For $p=0$, the contribution from both numerator and denominator,
\begin{align}
    f_{3pt}(n) = \frac{1}{6}\sum_{k_1\neq k_2\neq k_3 =1}^{n-1} \frac{C_{\Psi}^3 }{\lb \sin(\pi (k_3 - k_2)/n) \sin (\pi (k_2 - k_1)/n) \sin(\pi (k_3-k_1)/n)\rb^{\Delta_L}}
\end{align}
which is suppressed by a factor of $\zeta_n^{3\Delta_L}$ coming from the OPEs of $\Psi \Psi^*$.
Note that, for $p=0$ case, the three point function contribution from the numerator and the denominator of $\mathcal{F}_n$ are identical, so they cancel each other.

% For $p=1$, there are $n-p-1= n-2$ copies of 
% \begin{align}
%     \braket{\mO_L(z_k) \mO_{L'}(z_0)\mO_{L'}(z_{n-1}) } &= C_{\Psi\Phi}^2C_\Phi\frac{e^{-i3\pi (2s_{L'} -s_L)/2}}{\lb2\sin(\pi/n)\rb^{2\Delta_{L'}-\Delta_L}}\times\nn \\ 
%     & \qquad \qquad \qquad \qquad 
%     \sum_{k=1}^{n-2}\Bigg( \frac{e^{-i\pi s_L(k/n+3/2)}}{\lb2\sin(\pi(k+1)/n)\rb^{\Delta_L}} \frac{e^{-i\pi s_L((k+1)/n+1/2)}}{\lb 2\sin(\pi k/n)\rb^{\Delta_L}}\Bigg) \nn \\
%     &= C_{\Psi\Phi}^2C_\Phi e^{-3\pi is_{L'} } e^{-i\pi s_L/2}\lb 2\sin (\pi /n)\rb^{\Delta_L-2\Delta_{L'}}\times \nn \\
%     &\qquad \qquad \qquad \qquad \sum_{k=1}^{n-2}\frac{e^{-i\pi s_L (2k+1)/n}}{\lb(  4\sin(\pi (k+1)/n)\sin(\pi k /n)\rb)^{\Delta_L}} 
% \end{align}
The numerator of $\mathcal{F}_n$ including the corrections from the three point function would be,
\begin{align}
    1 + \zeta_n^{2\Delta_L}f(n) + \zeta_n^{3\Delta_L}f_{3pt}(n) + \cdots
\end{align}
The denominator of $\mathcal{F}_n$ including the corrections from the three point functions would be,
\begin{align}
y^{n-1}\Bigg( 1+ \zeta_n^{2\Delta_L} f(n)  + + \zeta_n^{3\Delta_L}f_{3pt}(n) + \sum_{L'} \zeta_n^{2\Delta_{L'} } C_{\Psi \Phi}^2 \mu(n,x,y;\Delta_{L'}) \nn\\
+ \sum_{L,L'}\zeta_n^{2\Delta_{L'} + \Delta_L}C^2_{\Psi\Phi} \mu_G(n,x,y;\Delta_{L'},\Delta_L)+\cdots\Bigg)
\end{align}
where we define
\begin{align}
    q &= y^{-1}-1 \\
    % \lambda_p &= D_n(\ep)^p \lb \frac{\sin(\pi (1-x)/n)}{\sin (\pi p /n)}\rb^{2\Delta_{L'}}\\
    G_p(n) &= \sum_{m=1}^{p-1} \frac{1}{\lb \sin(\pi m/n) \sin(\pi (p-m)/n)\rb^{\Delta_L}}\\
    \mu_G (n,x,y;\Delta_{L'},\Delta_L) &= \sum_{p=1}^{n-1}   \frac{q^p (n-p)D_n(\ep)^{p}}{\lb \sin(\pi p/n)\rb^{2\Delta_{L'}-\Delta_L}} \lb C_\Phi G_p(n) + C_\Psi G_{n-p}(n)\rb
\end{align}

----

$S(\rho_0 \| y\rho_0 + (1-y)\rho_\Phi)$. 

The denominator starts from 1, coming from $p=0$. For $p=1$, we just have a sum over two-point functions. 
\begin{align}
    y^{n-1}\lb 1  + \frac{(n-1) (y^{-1}-1)D_n(\ep)\zeta_n^{2\Delta_\Phi}}{(2\sin(\pi (1-\ep)/n))^{2\Delta_\Phi}}+ \cdots\rb
\end{align}
For $p=2$, we can have either have a string of $\Phi \Phi^*$ of length $p=2$, or a pair of $\Phi \Phi^*$ which do not share common site. The contribution from $p=2$ is,
\begin{align}
    \lb(y^{-1}-1)D_n(\ep)\rb^2\lb  (n-2)\lb\frac{\sin(\pi \ep/n)}{\sin(2\pi (1-\ep)/n)}\rb^{2\Delta_\Phi}  + \mO\lb\sin(\pi \ep/n)^{4\Delta_\Phi}\rb\rb
\end{align}
where the higher order terms come from the non-stringy contribution, which is clearly subleading term. 
Ignoring the non-stringy contribution, we have,
\begin{align}
    \mathcal{F}_n = \frac{1}{y^{n-1}\lb 1+ \zeta_n^{2\Delta_\Phi} \sum_{p=1}^{n-1} (n-p)\frac{\lb(y^{-1}-1)(D_n(\ep)\rb^p}{\lb\sin(p\pi (1-\ep)/n)\rb^{2\Delta_\Phi}} + \cdots \rb}
\end{align}
If we want to consider $y\to 1$ limit and keep the term up to the order of $(1-y)^2$ to relate it to quantum Fisher information, we thus have the following in the limit $y \to 1$,
\begin{align}
    S(\rho_0 \|( \rho_0 + (1-y)(\rho_\Phi - \rho_0))) &= -\log y - \p_n \Bigg(  (n-1)  (y^{-1}-1)D_n(\ep) \lb\frac{\zeta_n}{\sin(\pi(1-\ep)/n)}\rb^{2\Delta_\Phi}\nn \\
    &\qquad+ (n-2)\lb(y^{-1}-1)(D_n(\ep)\rb^2\lb\frac{\zeta_n}{\sin(2\pi(1-\ep)/n)}\rb^{2\Delta_\Phi} +\cdots
    \Bigg)
\end{align}
In the limit $y \to 1$, we have
\begin{align}
S(\rho_0 \|( \rho_0 + (1-y)(\rho_\Phi - \rho_0))) 
    &= (-1)^{2\Delta_\Phi+1}2^{-2\Delta_\Phi}(1-y)^2 +  (1-y)^2 \p_n\lb D_n(\ep)^2  \lb\frac{\zeta_n}{\sin(2\pi(1-\ep)/n)}\rb^{2\Delta_\Phi}\rb 
\end{align}
where the first term is a result of the derivative hitting $(n-2)$, which gives $(-1)^{2\Delta_\Phi}$ from $\sin(2\pi (1-\ep)/n)^{2\Delta_\Phi}$. We then have,
\begin{align}
S(\rho_0 \|( \rho_0 + (1-y)(\rho_\Phi - \rho_0))) 
    &=  (-1)^{2\Delta_\Phi+1}2^{-2\Delta_\Phi}(1-y)^2 \nn \\ &\qquad\qquad+ (1-y)^2 \p_n\Bigg( \lb\frac{n\sin(\pi \ep/n)}{\sin(\pi \ep)}\rb^{-4\Delta_\Phi}\times 
\lb\frac{\zeta_n}{\sin(2\pi(1-\ep)/n)}\rb^{2\Delta_\Phi}\Bigg)  \\ 
    &= (-1)^{2\Delta_\Phi+1}2^{-2\Delta_\Phi}(1-y)^2 \nn\\
    & \qquad\qquad+ (1-y)^2 (\sin(\pi \ep))^{4\Delta_\Phi} \times 
    \p_n\lb\frac{1}{n^2 \sin(\pi \ep/n) \sin(2\pi (1-\ep)/n)}\rb^{2\Delta_\Phi} \\
    &=(-1)^{2\Delta_\Phi+1}2^{-2\Delta_\Phi}(1-y)^2 + 2\Delta_\Phi (1-y)^2 (\pi \ep)^{4\Delta_\Phi}\Bigg( -2  + \frac{\pi \ep}{n^2} \cot \lb\frac{\pi \ep}{n}\rb   \nn\\
    &\qquad\qquad +\frac{2\pi (1-\ep)}{n^2} \cot \lb\frac{2\pi (1-\ep)}{n}\rb \Bigg)\Big|_{n=1}\\
    &= (-1)^{2\Delta_\Phi+1}2^{-2\Delta_\Phi}(1-y)^2 + 2\Delta_\Phi (1-y)^2 (\pi \ep)^{4\Delta_\Phi}\Bigg( -2 + \pi \ep \lb \frac{1}{\pi \ep} - \frac{\pi \ep}{3} + \cdots\rb \nn\\
    &\qquad \qquad - 2\pi (1-\ep)\lb \frac{1}{2\pi \ep} - \frac{2\pi \ep}{3} + \cdots \rb\Bigg)\\
    &= (-1)^{2\Delta_\Phi+1}2^{-2\Delta_\Phi}(1-y)^2 + 2\Delta_\Phi (1-y)^2 (\pi \ep)^{4\Delta_\Phi} \lb -\frac{1}{\ep} + \frac{4\pi^2 \ep}{3} - \frac{5\pi^2 \ep^2}{3} + \cdots \rb
\end{align}
% \JM{Check the contribution of $\braket{\Phi \mO \Phi^*}$.
% % It seems, that the contribution has also divergence, but it depends on $\Delta_L,\Delta_\Phi$ which whether it is suppressed or even more divergent.
% }

% The denominator of $\mathcal{F}_n$ takes the following form when we include the contribution from the three point function,
% \begin{align}
%     y^{n-1}\Bigg( 1+ (n-2) \lb (y^{-1}-1)D_n(\ep)\rb^2 \lb\frac{\zeta_n}{\sin(2\pi(1-\ep)/n}\rb^{2\Delta_\Phi}\left[1+ C_\Phi\frac{\lb\zeta_n\sin(2\pi(1-\ep)/n)\rb^{\Delta_L}}{\lb \sin (\pi (1-\ep)/n)\rb^{2\Delta_L}}+\cdots\right]\Bigg)
% \end{align}
% where we did not include the $p=1$ case since it has a $(n-1)$ factor sitting infront of it.
% \begin{align}
%     &S(\rho_0 \|( \rho_0 + (1-y)(\rho_\Phi - \rho_0)))\big|_{y\to 1} =\nn\\
%     &\Bigg((-1)^{2\Delta_\Phi +1}2^{-2\Delta_\Phi}(1-y)^2 + 2\Delta_\Phi (1-y)^2  (\pi \ep)^{4\Delta_\Phi}\lb -\frac{1}{\ep} + \frac{4\pi^2 \ep}{3} - \frac{5\pi^2 \ep^2}{3} + \cdots \rb \Bigg) \lb 1 + C_\Phi(-2)^{\Delta_L} (\pi \ep)^{2\Delta_L -2\Delta_\Phi}\rb\nn\\
%     &+ C_\Phi(1-y)^2 (-1)^{2\Delta_\Phi + \Delta_L + 1 }2^{\Delta_L-2\Delta_\Phi}(\pi \ep)^{4\Delta_\Phi} \Bigg( (\Delta_L - 2\Delta_\Phi)\frac{1}{\ep} + (\Delta_\Phi - 2\Delta_L) \frac{2\pi^2 \ep}{3} \nn \\
%     & 
%     -(\Delta_\Phi - 6\Delta_L )\frac{\pi^2 \ep^2}{6} -2( \Delta_\Phi + \Delta_L)\Bigg)
% \end{align}

If we consider the following,
\begin{align}
   & S
(\rho_\Psi \|( \rho_\Psi + (1-y)(\rho_0 - \rho_\Psi))) \\&\qquad\qquad= \p_n \log \left[\frac{1+\zeta_n^{2\Delta_L}f(n) + \cdots}{y^{n-1}\lb 1+\zeta_n^{2\Delta_L}f(n) + (n-2)\lb (y^{-1}-1)D_n(\ep)\rb^2 \lb\frac{\zeta_n}{\sin(2\pi(1+\ep)/n}\rb^{2\Delta_\Phi} +\cdots\rb}\right]
\end{align}
Here we did not include the contribution from $p=1$ since it does not contribute to the relative entropy, because of the $n-1$ factor sitting in front of it. Also notice that the numerator gets cancelled exactly from the $p=0$ term in the denominator. This gives us,
\begin{align}
    S
(\rho_\Psi \|( \rho_\Psi + (1-y)(\rho_0 - \rho_\Psi))) &=-2^{-2\Delta_\Phi}(1-y)^2 + 2\Delta_\Phi (1-y)^2 (\pi \ep)^{4\Delta_\Phi}\Bigg( -2  + \frac{\pi \ep}{n^2} \cot \lb\frac{\pi \ep}{n}\rb   \nn\\
    &\qquad\qquad +\frac{2\pi (1+\ep)}{n} \cot \lb\frac{2\pi (1+\ep)}{n}\rb \Bigg)\Big|_{n=1}\\
    &= -2^{-2\Delta_\Phi}(1-y)^2 + 2\Delta_\Phi (1-y)^2 (\pi \ep)^{4\Delta_\Phi}\Bigg( -2 + \pi \ep \lb \frac{1}{\pi \ep} - \frac{\pi \ep}{3} + \cdots\rb \nn\\
    &\qquad \qquad + 2\pi (1+\ep)\lb \frac{1}{2\pi \ep} - \frac{2\pi \ep}{3} + \cdots \rb\Bigg)\\
    &= -2^{-2\Delta_\Phi}(1-y)^2 + 2\Delta_\Phi (1-y)^2 (\pi \ep)^{4\Delta_\Phi}\lb \frac{1}{\ep} - \frac{4\pi^2 \ep}{3} - \frac{5\pi^2 \ep^2}{3} + \cdots \rb
\end{align}
Here, we get positive divergence, in contrast to the previous case. 

----

\JM{If we consider the general OPE with the identity piece as given in eq.~[\ref{eq: general OPE with identity}], the denominator will be given by}
\begin{align}
    &\delta_{\Psi \Phi} \sum_{p=0}^{n-1}\Big[(y^{n-1-p}\lb (1-y)D_n(\ep) \rb^p\binom{n-1}{p} \lb 1 + \lb 2\pi (1-x)/n\rb ^{2\Delta_{L}} f(n) + \cdots\rb  \Big]\nn \\
    &+ (1- \delta_{\Psi \Phi}) \Big( \lb 2\pi (1-x)/n\rb ^{2\Delta_{L'}} \sum_{p=1}^{n-1}\Big[  \frac{y^{n-1-p}\lb (1-y)D_n(\ep)\rb^p(n-p)}{\lb 2 \sin(\pi p/n)\rb ^{2\Delta_{L'}}}\Big]\Big)
\end{align}
\JM{Since the first term has $\delta_{\Psi \Phi}$, it will act on $D_n(\ep)$ to make it equal to 1, and we are left with $(y + (1-y))^{n-1} =1$. Hence, the first term now exactly cancels the numerator. We are left with,}
\begin{align*}
\mathcal{F}_n
    &=\frac{1}{\JM{(1- \delta_{\Psi \Phi}) \Big( \lb 2\pi (1-x)/n\rb ^{2\Delta_{L'}} \sum_{p=1}^{n-1}\Big[  \frac{y^{n-1-p}\lb (1-y)D_n(\ep)\rb^p(n-p)}{\lb 2 \sin(\pi p/n)\rb ^{2\Delta_{L'}}}\Big]\Big)}}\nn\\ 
    &= \frac{1}{\JM{(1- \delta_{\Psi \Phi}) \Big( \lb 2\pi (1-x)/n\rb ^{2\Delta_{L'}} y^{n-1}\mu(n,x)\Big)}}\nn \\
% &=\JM{y^{1-n}\lb 1-\lb\frac{2\pi (1-x)}{n}\rb^{2\Delta_{L'}}  C_{\Psi\Phi}^2 \:\mu(n,x)+\cdots \rb}\nn\\
\mu(n,x)&=\sum_{p=1}^{n-1}\frac{(n-p)((y^{-1}-1)D_n(1-x))^p}{(2\sin(\pi p/n))^{2\JM{\Delta_{L'}}}}
\end{align*}
\JM{Then we have,}
\begin{align}
    \JM{\p_n \log \mathcal{F}_n =-\p_n \log \Bigg[ (1-\delta_{\Psi \Phi})\Bigg( \lb 2\pi (1-x)/n\rb ^{2\Delta_{L'}} \sum_{p=1}^{n-1}\Bigg[  \frac{y^{n-1-p}\lb (1-y)D_n(\ep)\rb^p(n-p)}{\lb 2 \sin(\pi p/n)\rb ^{2\Delta_{L'}}}\Bigg]\Bigg)  \Bigg]}
\end{align}
\JM{We can see that, if we take the limit $y\to 0$ limit in this expression, the only term that survives in $\mu(n,x)$ is when $p = n-1$, which gives the leading divergent piece. 
% This matches with our previous answer in eq.~[\ref{eq: small 1-x limit}]. 
Similarly, taking $y\to 1$ limit, we get $\log \mathcal{F}_n$ vanishes, which shows that the relative entropy between two identical states vanishes, as expected.}

\JM{For $y = 1/2$, we have}
\begin{align}
    \JM{S(\rho_{\Psi}\| (\rho_\Psi + \rho_\Phi)/2)} &= \JM{-\p_n \log \Bigg [ \zeta^{2\Delta_{L'}}_n 2^{-(n-1)}\sum_{p=1}^{n-1}\frac{(n-p)\Big((1  + 2p\pi^2 \ep^2 (n^2 - 1)(h_\Psi - h_\Phi)/(3n^2) + \mO(\ep^4)\Big)}{(2\sin(\pi p/n))^{2\Delta_{L'}}} \Bigg]\Bigg|_{n=1}} \nn \\
    &=\JM{\log 2 - \p_n \log \Big[  \zeta_n^{2\Delta_{L'} }\big( \mu_0(n) + \ep^2 \mu_1(n)\big)\Big]}
\end{align}
\JM{where}
\begin{align}
  \mu_0(n)&=\JM{\sum_{p=1}^{n-1}\frac{n-p}{(2\sin(\pi p/n))^{2\Delta_{L'}}}=\frac{1}{2}\sum_{p=1}^{n-1}\frac{p+(n-p)}{(2\sin(\pi p/n))^{2\Delta_{L'}}}}\nn\\
  &=\JM{\frac{n}{2}\sum_{p=1}^{n-1}\frac{1}{(2\sin(\pi p/n))^{2\Delta_{L'}}}=S_2(n) }\\
  \mu_1(n) &= \JM{\frac{(n^2-1)\lb2\pi^2 (h_\Psi - h_\Phi)\rb}{3n^2}\sum_{p=1}^{n-1} \frac{(n-p)p}{(2\sin(\pi p/n))^{2\Delta_{L'}}}}\nn \\
  &= \JM{\frac{(n^2-1)\lb2\pi^2 (h_\Psi - h_\Phi)\rb}{3n^2} \lb nS_2(n) - \sum_{p=1}^{n-1}\frac{p^2}{(2\sin(\pi p/n))^{2\Delta_{L'}}}\rb} 
\end{align}

\subsection{Some examples}
We consider a mixture of coherent states. We will evaluate the relative entropy $S(\rho_\alpha \|y\rho_\alpha + (1-y)\rho_\beta)$, where the $\rho_{\alpha},\rho_\beta$ are reduced density matrices of some excited states created by Vertex operators $V_\alpha, V_\beta$. As discussed in the previous section, the numerator of $\mathcal{F}_n$ evaluates to $1+ \zeta_n^{2\Delta_L}f(n)$, however for the denominator, we get the divergent piece from the two point function of the vertex operators $V_{\alpha- \beta}$ and $V_{\beta- \alpha}$. This gets multiplied with the two point function of the currents $J$, with the appropriate OPE coefficient that sits infront of $J$. The denominator takes the following form,
\begin{align}
    &\sum_{p=1}^{n-1} y^{n-1-p} (1-y)^p D_n(\ep)^p (n-p) \lb 2\sin(p\pi(1-\ep)/n)\rb^{-2(\alpha - \beta)^2} \left[  1+ \zeta_n^2 g_p(n) + \cdots  \right]\\
    &g(p) =  \frac{1}{2}\sum_{k\neq k' =0}^{n-1} c_k c_{k'} \lb 2\sin( \pi (k-k')/n )\rb^{-2}\\
    &c_k = \begin{cases}
        \beta , \qquad  k = 1,2,\cdots ,p \\
        \alpha , \qquad k = 0 ,p+1, p+2, \cdots, n-1
    \end{cases}
\end{align}
We can introduce an indicator function,
\begin{align}
    \chi_k^{(p)} = \begin{cases}
        1 \qquad k \in \{ 1, 2 ,\cdots , p\}\\
        0 \qquad k\in \{0,p+1,p+2, \cdots , n-1 \} 
    \end{cases}
\end{align}
to write 
\begin{align}
    c_k = \alpha + (\beta - \alpha)\chi_k^{(p)}
\end{align}
Then $g_p(n)$ simplifies to,
\begin{align}
    g_p(n) &= \frac{1}{2}\sum_{k\neq k' = 0}^{n-1} c_k c_{k'} \lb 2\sin( \pi (k-k')/n )\rb^{-2} \\
    &= \alpha^2 S_2(n) + \frac{\alpha (\beta - \alpha) }{2}\sum_{k\neq k'}^{n-1}\frac{\chi_k + \chi_{k'}}{\lb2\sin (\pi (1-\ep)/n)\rb^2} + \frac{1}{2}(\beta - \alpha)^2 \sum_{k\neq k'=0}^{n-1} \frac{\chi_k \chi_{k'}}{\lb2\sin (\pi (1-\ep)/n)\rb^2}
\end{align}
The second sum can be simplified to,
\begin{align}
     \frac{\alpha (\beta - \alpha) }{2}\sum_{k\neq k'}^{n-1}\frac{\chi_k + \chi_{k'}}{\lb2\sin (\pi (k-k')/n)\rb^2} &= \alpha (\beta - \alpha) \sum_{k\neq k' = 0}^{n-1} \frac{\chi_k}{\lb2\sin (\pi (k-k')/n)\rb^2} \\
     &=\alpha (\beta - \alpha) \sum_{k=1}^p \sum_{k'\neq k} \frac{1}{\lb2\sin (\pi (k-k')/n)\rb^2}\\
     &= \alpha (\beta - \alpha)p \frac{2S_2(n)}{n}
\end{align}
where in the last step, the sum over $k'$ is $\frac{2S_2(n)}{n}$. Consider the third sum,
\begin{align}
    \frac{1}{2}(\beta - \alpha)^2 \sum_{k\neq k'=0}^{n-1} \frac{\chi_k \chi_{k'}}{\lb2\sin (\pi (1-\ep)/n)\rb^2} &= \frac{1}{2}(\beta - \alpha)^2 \sum_{k\neq k'=1}^{p} \frac{\chi_k \chi_{k'}}{\lb2\sin (\pi (1-\ep)/n)\rb^2} \\
    &= (\beta-\alpha)^2 \sum_{m=1}^{p-1} \frac{p-m}{\lb2\sin(\pi m /n)\rb^2}
\end{align}
Now we can write the full $\mathcal{F}_n$,
\begin{align}
    \mathcal{F}_n &= \frac
    {1+\zeta_n^{2}f(n)+\cdots}{y^{n-1}\lb 1+\zeta_n^2f(n) + \sum_{p=1}^{n-1} (y^{-1}-1)^pD_n(\ep)^p (n-p) \frac{\zeta_n^{2(\alpha-\beta)^2}}{\lb 2\sin \lb p\pi (1-\ep)/n\rb\rb^2}\lb 1+\zeta_n^2 g_p(n) + \cdots \rb\rb } \\
    g(p) & = \lb  \alpha^2 + \frac{2\alpha (\beta - \alpha)p}{n}\rb S_2(n)  +(\beta - \alpha)^2 \\
    H_p(n) &= \sum_{m=1}^{p-1} \frac{(p-m)}{\lb 2\sin (\pi m /n)\rb^2}.
\end{align}
To simplify the notation, we define,
\begin{align}
    &q = y^{-1}-1, \qquad \lambda_p = D_n(\ep)^p \lb \frac{\zeta_n}{\sin(p \pi (1-\ep)/n)}\rb^{(\alpha - \beta)^2} \\
    &\mu_0(n,y) = \sum_{p=1}^{n-1} q^p (n-p)\lambda_p \\
    & \mu_1(n,y) = \sum_{p=1}^{n-1}q^p p(n-p) \lambda_p \\
    &\mu_H(n,y) = \sum_{p=1}^{n-1} q^p (n-p) \lambda_p H_p(n)
\end{align}
In terms of these sums, we get,
\begin{align}
    \mathcal{F}_n &= \frac{1 + \zeta_n^2 f(n)+\cdots}{y^{n-1}\lb  1+ \zeta_n^2 f(n) + \mu_0(n,y) + \zeta_n^2\lb  \alpha^2 S_2(n) \mu_0(n,y)+ \frac{2\alpha (\beta - \alpha)S_2(n)}{n} \mu_1(n,y) + (\beta - \alpha)^2 \mu_H(n,y)\rb \rb} \\
    \p_n \log \mathcal{F}_n &= -\log y - \mu_0'(1,y) - (\pi \ep)^2 (\beta - \alpha)^2 \p_n \lb \sum_{p=1}^{n-1} q^p (n-p) \lambda_p H_p \rb \Bigg|_{n=1}
\end{align}
where in the last step, we used the fact $S_2(n)=0$, when the derivative hits $\mu_0,\mu_1$ and is evaluated at $n=1$. 

The limit $y\to 1$ is trivial, since $q \to 0$, and hence $\mu_0 = \mu_1 = \mu_H = 0$. $\mathcal{F}_n = 1$, and hence the relative entropy vanishes as expected. 

For the other limit, $y \to 0$ limit, we pull the $y^{n-1}$ back into the paranthesis in the denominator. Then in the sums $\mu_0,\mu_1,\mu_H$, only $p = n-1$ contributes. Then we have the simplification,
\begin{align}
    \lambda_{n-1}  &= D_n(\ep)^{n-1} \\
    \mu_0(n,y) &= \lambda_{n-1} \\
    \mu_1 & = (n-1)\lambda_{n-1}\\
    \mu_H &= \lambda_{n-1} \sum_{p=1}^{n-2} \frac{n-1-m}{(2\sin(\pi m /n))^2} = \lambda_{n-1} \sum_{m=1}^{n-1}\frac{n-1-m}{(2\sin(\pi m/n))^2} \\
    &= \lambda_{n-1} \frac{n-2}{n}S_2(n) 
\end{align}
In this limit, we have
\begin{align}
    \mathcal{F}_n &= \frac{1 + \zeta_n^2 f(n)+\cdots}{  \lambda_{n-1} \lb 1 + \zeta_n^2 \lb \alpha^2 S_2(n) + \frac{2\alpha (\beta - \alpha)(n-1)}{n}S_2(n) + (\beta-\alpha)^2  \lb\frac{n-2}{n}\rb S_2(n)\rb\rb } 
\end{align}
Taking the log and derivative, the $\lambda_{n-1}$ piece gives the divergent piece, the $\alpha^2S_2(n)$ in the denominator cancels with the numerator, while the second term proportional to $\frac{n-1}{n}$ does not contribute to relative entropy when we take the limit $n\to 1$. So we are left with,
\begin{align}
    S(\rho_\alpha \| \rho_\beta) = (\alpha - \beta)^2 \lb \frac{1}{\ep} - \frac{\ep}{3}\rb + (\alpha - \beta)^2 \frac{(\pi \ep)^2}{3}+\cdots
\end{align}
Similarly, for the non-chiral vertex operators, we will get,
\begin{align}
    S(\rho_\alpha \| \rho_\beta) = 2(\alpha - \beta)^2 \lb \frac{1}{\ep} - \frac{\pi^2\ep}{3}\rb + 2(\alpha - \beta)^2 \frac{(\pi \ep)^2}{3}+\cdots
\end{align}

\subsection{Some limits}

\paragraph{Limit of $y\to 0$:}

It is instructive to think about the limit $y\to 0$ and $x\to 1^-$. To simplify, further, we take $\rho_\alpha$ to be the vacuum density matrix. In this case, the numerator of $\mathcal{F}_n$ is just one, and the denominator is contains only the $p=n-1$ term which is an $2(n-1)$ point function of $V_\beta$:
\begin{align}
  D_n(1-x)^{n-1}\frac{2C_{\alpha\beta}^2}{(2\sin(\pi(n-1)/n))^{2\Delta_L}}\lb\frac{2\pi (1-x)}{n}\rb^{2\Delta_L}+\cdots 
\end{align}

To find the relative entropy we need to compute
\begin{align}
  \partial\log\mathcal{F}_n|_{n\to 1}&=\partial\left[(1-n)\log y-2C_{0\beta}^2\lb\frac{2\pi (1-x)}{n} \rb^{2\Delta_L} \mu(n,x)+\cdots \right]_{n\to 1}\nn\\
  &=-\log y-2C_{\alpha\beta}^2(2\pi (1-x))^{2\Delta_L}(-2\Delta \: \mu(1,x)+\partial_n\mu(1,x)|_{n\to 1})+\cdots
\end{align}

In Appendix \ref{app:analytic-continuation}, we use a method similar to that of Tonni et al based on Carlson's theorem that analytically continues $\mu(n)$ to complex $n$, and show that 
\begin{align}
  \partial_n\mu(n)|_{n=1}=
\end{align}


\paragraph{Mixtures with vacuum:}

We focus on the case where $\alpha=0$. In this case, $\mathcal{F}_n$ simplifies to
\begin{align}
  1/\mathcal{F}_n&=\braket{\tilde{O}_{0|\beta}^{(n-1)}}\nn\\
\tilde{O}_{0|\beta}^{(n-1)}&=\sum_{p=0}^{n-1}y^{n-1-p}(1-y)^p
\end{align}
%  \NL{HERE}

  % In this case, we have
  % \begin{eqnarray}
  %   &&S(\rho_0\|y\rho_0+(1-y)\rho_\beta)=\partial_n\left[ \frac{(y\braket{1}_{\mathcal{C}_1}+(1-y)\braket{V_\beta(u^+) V_{-\beta}(u^-)}_{\mathcal{C}_1})^{n-1}}{\braket{\mO^{(n-1)}_\beta}_{\mathcal{C}_n}}\right]_{n=1}\nn\\
  %   &&\mO_\beta^{(n-1)}=\sum_{p=0}^{n-1}y^{n-1-p}(1-y)^p\sum_{\substack{w\in\{0,\beta\}^{n-1}\\ N_\beta(w)=p}}
  % V_{w_1}(u_1^+)V_{-w_1}(u_1^-)\cdots V_{w_{n-1}}(u_{n-1}^+)V_{-w_{n-1}}(u_{n-1}^-)\nn\\
  % \end{eqnarray}

% \section{Examples of $x\to 0^+$ and $x\to 1^-$ limits}
\section{\texorpdfstring{Examples of $x\to 0^+$ and $x\to 1^-$ limits}
{Examples of x->0+ and x->1- limits}}

\subsection{Non-chiral compact free boson 1D}
$\ket{00}$ denotes the ground state, while $\ket{ss}$ denotes the primary state that corresponds to the vertex operator $V(z,\bar z) = \exp \lb \frac{i}{\sqrt{2}}\phi(z) + \frac{i}{\sqrt{2}}\bar \phi(\bar z) \rb$. We also consider another primary $\ket{3s,s}$ that corresponds to the vertex operator $V(z,\bar z) = \exp \lb i\frac{3}{\sqrt{2}}\phi(z) + \frac{i}{\sqrt{2}}\bar \phi(\bar z) \rb$.
Consider the mixture state with $p\equiv y=1/2, \bar \rho_A = \frac{1}{2}\lb \rho_A^{00} + \rho_A^\chi \rb$. 

For $x \to 0^+$ limit, we have $\Delta_L =1$. For the state $\ket{ss}$ we have,  $\alpha = C_{\alpha (-\alpha)}^L =  \bar \alpha = \bar C _{\alpha(- \alpha)}^L = 1 $. And similarly, for $\ket{3s,s}$, we have  $\alpha = C_{\alpha (-\alpha)}^L = 3/\sqrt{2}; \qquad  \bar \alpha = \bar C _{\alpha(- \alpha)}^L = 1/\sqrt{2} $. Hence,we have
\begin{align}
    S(\rho_A^{00}\| \bar \rho_A^{(\chi)}) = \begin{cases}
        \frac{\pi^2}{12} x^2 & \chi = \ket{ss} \\
        \frac{5\pi^2 }{12}x^2 & \chi = \ket{3s,s}
    \end{cases}
\end{align}

For $x \to 1^-$ limit, we have to consider the $\Delta_{L'}$ for the mixing OPE. For $\ket{ss}$, $\Delta_{L'} = \frac{1}{2}\lb \lb 1/\sqrt{2} -0\rb^2 + \lb 1/\sqrt{2} -0\rb^2 \rb= 1/2$. The mixing OPE coefficient $C_{\alpha ,I}^\alpha = \bar C_{\bar \alpha, I}^{\bar \alpha} = 1$. Hence, 
\begin{align}
    \delta_\beta^{(\chi)} &= \log 2 - \frac{1}{2}\lb S(\rho_B^{00} \| \bar \rho^{(\chi)}_B  ) + S(\rho^\chi_B \| \bar \rho ^{(\chi)}_B)   \rb  \\&= \log 2 - \frac{1}{2}\lb \log 2 - (\pi \ep) \frac{\sqrt{\pi} \Gamma(3/2)}{4\Gamma(2)}\rb - \frac{1}{2}\lb  \log 2 - (\pi \ep) \frac{\sqrt{\pi} \Gamma(3/2)}{4\Gamma(2)}\rb \\
    &= \frac{\pi^2\ep}{8} \\
    &= 1.234 \ep. 
\end{align}


Similarly, for $\ket{3s,s}$,
$\Delta_{L'} = \frac{1}{2}\lb\left(3/\sqrt{2}-0\right)^2
+ \left(1/\sqrt{2}-0\right)^2\rb = 5/2$.
And the mixing OPE coefficients are
$C_{\alpha,I}^{\alpha}=1$ and
$\bar{C}_{\bar{\alpha},I}^{\bar{\alpha}}=1$.
Hence,
\begin{align}
    \delta_\beta^{(\chi)}  &= \log 2 - \frac{1}{2}\lb \log 2 - (\pi \ep)^{5}\frac{\sqrt{\pi}\Gamma(7/2)}{4\Gamma(4)} \rb - \frac{1}{2}\lb  \log 2 - (\pi \ep)^{5}\frac{\sqrt{\pi}\Gamma(7/2)}{4\Gamma(4)}\rb  \\
    &=(\pi \ep)^{5}\frac{\sqrt{\pi}\Gamma(7/2)}{4\Gamma(4)} \\
    &= 75.108 \ep^{5} 
\end{align}
In the above two cases, notice, $S(\rho_B^{00} \| \bar \rho^{(\chi)}_B  ) = S(\rho^\chi_B \| \bar \rho ^{(\chi)}_B) $. 

\subsection{Ising CFT}
We have two non-identity primaries, namely, $\sigma \equiv
 \lb \frac{1}{16}, \frac{1}{16}\rb$ and $\varepsilon\equiv \lb \frac{1}{2},\frac{1}{2}\rb$. 
We have the following OPEs,
\begin{align}
    \sigma(z_1)\times \sigma(z_2) &= \frac{1}{|z_{12}|^{1/4}}\lb 1+ \frac{1}{2} |z_{12}| \varepsilon(z_{2}) + \cdots \rb \\
    \varepsilon(z_1) \times \varepsilon(z_2) &= \frac{1}{|z_{12}|^2} \lb 1+ 2|z_{12}|^2 T(z_{12}) + 2|\bar{z}_{12}|^2\bar{T}(\bar{z}_{12})  + \cdots\rb \\
    \varepsilon(z_1) \times \sigma(z_2) &= \frac{1}{2|z_{12}|}\lb \sigma(z_2) + 4|z_{12}| \p \sigma(z_2) + 4|\bar{z}_{12}|\bar{\p}\sigma  (z_2) + \cdots \rb \\
    \sigma(z_1) \times \varepsilon(z_2) &= \frac{1}{2|z_{12}|}\lb \sigma(z_2) -3|z_{12}| \p \sigma(z_2) -3|\bar{z}_{12}|\bar{\p}\sigma  (z_2) + \cdots \rb 
\end{align}
So, for small interval, we have the following results,
\begin{align}
    S(\rho_I \| y \rho_I + (1-y)\rho_\sigma) &= \frac{1}{12} (1-y)^2(\pi x)^2 \\
    S(\rho_\sigma \| y \rho_I + (1-y)\rho_\sigma) &= \frac{1}{12} y^2(\pi x)^2 \\ 
    S(\rho_I \| y \rho_I + (1-y)\rho_\varepsilon) &= \frac{8}{15} (1-y)^2 (\pi x)^4 \\
    S(\rho_\varepsilon \| y \rho_I + (1-y)\rho_\varepsilon) &= \frac{8}{15} y^2 (\pi x)^4 \\
    S(\rho_\sigma \| y \rho_\sigma + (1-y)\rho_\varepsilon) &= \frac{1}{12} (1-y)^2(\pi x)^2 \\
    S(\rho_\varepsilon \| y \rho_\sigma + (1-y)\rho_\varepsilon) &= \frac{1}{12} y^2(\pi x)^2
\end{align}
where we used the two point function of stress tensor in the third and fourth line, which comes from the OPEs of two $\varepsilon$.

Similarly, for the large interval limit, we have the following,
\begin{align}
    S(\rho_I \| y \rho_I + (1-y)\rho_\sigma) &= -\log y - (2\pi (1-x))^{1/4} \mu'(1,x,y;1/8)\\
    S(\rho_I \| y \rho_I + (1-y)\rho_\varepsilon) &= -\log y - (2\pi (1-x))^{2} \mu'(1,x,y;1)\\ 
    S(\rho_\sigma \| y \rho_\sigma + (1-y)\rho_\varepsilon) &= -\log y -  \frac{1}{4}(2\pi (1-x))^{1/4} \mu'(1,x,y;1/8)\\
\end{align}
where the $1/4$ in the third line comes from the square of the OPE coefficient $C_{\sigma\varepsilon} = 1/2$. 

\textbf{Numerics matching: } 

 For $x \to 1^-$ limit, we consider the mixture of $\ket I -\ket \sigma$, and  $\ket I -\ket \varepsilon$. For $\ket{I}-\ket{\sigma}$, $\Delta_{L'} = 1/8$ , while for $\ket{I}-\ket{\varepsilon}$, $\Delta_{L'} = 1$. The mixing OPE coefficients are trivial and equal to 1 in these cases, as the second operator is Identity. Hence we have,
\begin{align}
    \delta_\beta|_{\ket I-\ket{\sigma}} = (\pi \ep)^{1/4} \frac{\sqrt{\pi} \Gamma(9/8)}{4\Gamma(13/8)} = 0.619 \ep^{1/4}. 
\end{align}
Similarly, we can compute it for the other mixture with $\ket{\varepsilon}$. Finally,
 \begin{align}
     \delta_\beta^{(\chi)} = \begin{cases}
         0.619  \ep^{1/4} & \ket{I}-\ket{\sigma} \\
         3.289 \ep^2 & \ket{I}-\ket{\varepsilon}
     \end{cases}
 \end{align}

\subsection{Free boson}




\subsection{Minimal models}

% \section{Beyond small intervals}

% \subsection{Massless free boson}

%   In the free boson theory, the equation of motion is $\partial\bar{\partial}\phi(z,\bar{z})=0$ which means that the solutions to the equations of motion can be written as
%   \begin{align}
%     \phi(z,\bar{z})=\varphi(z)+\bar{\varphi}(\bar{z}).
%   \end{align}
%   The algebra of observables factors as the tensor product of a chiral and anti-chiral $U(1)$ current algebras with two–point functions 
% \[
% \langle \varphi(z)\varphi(0)\rangle=-\ln z, \qquad
% \langle \bar\varphi(\bar z)\bar\varphi(0)\rangle=-\ln \bar z, \qquad
% \langle \varphi(z)\bar\varphi(0)\rangle=0 .
% \]
% From the Schwinger-Dyson equations we know that up to contact terms $\partial\bar\partial\phi=0$ holds as an operator equation. This means that we can replace vertex operators in $\phi$ and $\partial_\mu \phi$ with their on-shell values up to contact terms. In free field theory, normal ordering keeps track of these contact terms. A general vertex operator of conformal weight $(\alpha ^2/2,\bar{\alpha}^2/2)$ is
% \begin{align}
% V_{\alpha,\bar{\alpha}}(z,\bar{z})&=:e^{i(\alpha\varphi(z)+\bar{\alpha}\bar{\varphi}(\bar{z}))}:=V_\alpha(z)V_{\bar\alpha}(\bar z)\nn\\
% V_\alpha(z)&=:e^{i \alpha\varphi(z)}:,\qquad V_{\bar\alpha}(\bar{z})=:e^{i \bar\alpha\bar\varphi(\bar z)}:
% \end{align}
% For now, we are assuming that the field is non-compact, i.e., $\phi:\mathbb{R}^{1,1}\to\mathbb{R}$, therefore $\alpha$ and $\bar{\alpha}$ are arbitrary real charges. We postpone a discussion of compact boson, the quantization of charge and winding until subsection \ref{sec:compact-boson}. 

% For linear functional of free fields we have
% \begin{align}
% :e^A::e^B:=e^{\langle A,B\rangle}\: :e^{A+B}:
% \end{align}
% Therefore, for the product of vertex operators we have the Wick contractions
% \[
% V_{\alpha,\bar\alpha}(z,\bar z)V_{\beta,\bar\beta}(0,0)
% =
% z^{\alpha\beta}\bar z^{\bar\alpha\bar\beta}V_{\gamma,\bar{\gamma}},\qquad \gamma=\alpha+\beta,\qquad \bar{\gamma}=\bar{\alpha}+\bar{\beta},
% % \nn\\
% % V_{}
% % :\exp\!\big(i\alpha\varphi(z)+i\beta\varphi(0)
% % +i\bar\alpha\bar\varphi(\bar z)+i\bar\beta\bar\varphi(0)\big): .
% \]

% \subsection{Vertex operators}

% The modular Hamiltonian of the vacuum on a cylinder of circumference $2\pi$ in the interval $(-x,x)$ is given by
% \begin{align}
% H_0&=\frac{(2\pi)^2}{2}\int_{-x}^x\frac{\cos(\pi y)-\cos(\pi x)}{\sin(\pi x)}T_{00}(x).
% \end{align}
% In a state with uniform energy density such as an exact energy eigenstate with dimension $h_\alpha$ on a cylinder of circumference $2\pi$ we have $\Delta T_{00}(y)=h/\pi$, and as a result
% \begin{align}
% \Delta H_\sigma\equiv \text{tr}((\rho_\alpha-\rho_0)H_0)=4h(1-\pi x\cot(\pi x))
% \end{align}
% For vertex operator with $h_\alpha=\bar{h}_\alpha$, a direct calculation shows that $S(\rho_\alpha)=S(\rho_0)$ which implies that 
% \begin{align}
% S(\rho_\alpha\|\rho_0)&=2\alpha^2(1-\pi x \cot(\pi x))=\Delta H_0\ .
% \end{align}
% As a result, we find $S(\rho_\alpha)=S(\rho_0)$ for all $x$ and $\alpha$. This suggests that $\rho_\alpha=V\rho_0V^\dagger$ for some unitary $V$. In fact, since $H_0$ generates a geometric flow in the Euclidean plane, it is indeed true that 
% \begin{align}
% V_{\alpha,\bar{\alpha}}(z,\bar{z})\sigma_0^{1/2}=\sigma_0^{1/2}V_{\alpha,\bar{\alpha}}(z',\bar{z}')
% \end{align}
% Then, since $\rho_\alpha=\sigma_0^{(1/2-x)}V_\alpha \sigma_0^{2x}V_{-\alpha}\sigma_0^{(1/2-x)}$, we have
% \begin{align}
%   \rho_\alpha=V_\alpha \sigma_0 V_\alpha^\dagger
% \end{align}
% and since $V_\alpha$ is unitary, we have
% \begin{align}
%   H_\alpha=V_\alpha \rho_0 V_\alpha^\dagger.
% \end{align}
% As a result,
% \begin{align}
% \tr(\rho_\alpha\log\rho_\beta)&=\tr(\rho_{\alpha-\beta}\log \rho_0)\nn\\
% S(\rho_\alpha\|\rho_\beta)&=S(\rho_{\alpha-\beta}\|\rho_0)=S(\rho_0\|\rho_{\alpha-\beta}).
% \end{align}
% In particular, this implies further that for a probability distribution $y_0+\sum_i y_i=0$ we have
% \begin{align}
% S(\rho_\alpha\|y_0\rho_\alpha+\sum_i y_i\rho_{\beta_i})&=S(\rho_0\|y_0\rho_0+\sum_i y_i\rho_{\beta_i-\alpha})
% \end{align}

% For an arbitrary energy eigenket we have
% \begin{align}
% S(\rho_h\|\rho_0)=S(\rho_h)-S(\rho_0)+4h(1-\pi x\cot(\pi x))
% \end{align}
% For instance, in the case of the eigenket generated by the action of $i\partial\varphi$ in the chiral free boson theory, we have
% \begin{align}
% &S(\rho_{i\partial\phi}\|\rho_0)=2(1-\pi x \cot(\pi x))+2\lb \log(2\sin(\pi x))+\psi_0(\csc(\pi x)/2)+\sin(\pi x)\rb
% \end{align}
% which implies the second term 
% \begin{align}
% 2\lb \log(2\sin(\pi x))+\psi_0(\csc(\pi x)/2)+\sin(\pi x)\rb
% \end{align}
% has the interpretation of $\Delta S$ as can be checked by the symmetry $x\to 1-x$. For these states we also have
% \begin{align}
% S(\rho_{i\partial \varphi}\|\rho_\beta)=S(\rho_{i\partial \varphi}\|\rho_0)+S(\rho_0\|\rho_\beta)
% \end{align}
% which requires an extra condition 
% \begin{align}
%   \tr((\rho_{i\partial \varphi}-\rho_0)(H_\beta-H_0))=0.
% \end{align}
% To derive this relation, we note that the current operator can be defined either as a spatial derivative of a vertex operator at zero charge, or the normal-ordered product of two vertex operators with unit and opposite charges
% \begin{align}
%   (i\partial \varphi)&=\frac{1}{\alpha}\partial_z V_\alpha|_{\alpha=0}=\partial_\alpha\partial_z V_\alpha|_{\alpha=0}\nn\\
%   (i\partial \varphi)&=\lim_{z\to 0} :V_\alpha(z)V_{-\alpha}(0):
% \end{align}
% Note that vertex operators of unit charge $\psi=e^{i\varphi}$ can be interpreted as a complex fermion:
% \begin{align}
% &\psi(z) \psi(w)=\frac{1}{(z-w)}:e^{i(\varphi(z)+\varphi(w))}=-\psi(w)\psi(z)\nn\\
% &(i\partial\varphi)(z)=J(z)=:\psi^\dagger\psi(z):
% \end{align}




% \NL{HERE}

% Expand the chiral/anti-chiral modes around the center of coordinates:
% \[
% \varphi(z)=\varphi(0)+z\,\partial\varphi(0)+\frac{z^2}{2}\partial^2\varphi(0)+O(z^3),
% \]

% \[
% \bar\varphi(\bar z)=
% \bar\varphi(0)+\bar z\,\bar\partial\bar\varphi(0)
% +\frac{\bar z^2}{2}\bar\partial^2\bar\varphi(0)+O(\bar z^3).
% \]
% and substituting into the exponent gives
% \begin{align}
% i\alpha\varphi(z)+i\beta\varphi(0)
% +i\bar\alpha\bar\varphi(\bar z)+i\bar\beta\bar\varphi(0)&=
% i\gamma\varphi(0)+i\bar\gamma\bar\varphi(0)
% +i\alpha z\,\partial\varphi(0)
% +i\bar\alpha\bar z\,\bar\partial\bar\varphi(0)\nn\\
% &\quad
% +\frac{i\alpha}{2}z^2\partial^2\varphi(0)
% +\frac{i\bar\alpha}{2}\bar z^2\bar\partial^2\bar\phi(0)
% +\cdots .
% \end{align}
% where we have defined the fused charges $
% \gamma=\alpha+\beta$ and $\bar\gamma=\bar\alpha+\bar\beta$.
% Therefore
% \[
% \begin{aligned}
% V_{\alpha,\bar\alpha}(z,\bar z)V_{\beta,\bar\beta}(0,0)
% &=
% z^{\alpha\beta}\bar z^{\bar\alpha\bar\beta}
% :\exp\!\Big(i\gamma\phi(0)+i\bar\gamma\bar\phi(0)\Big) \\
% &\qquad \times
% \exp\!\Big(
% i\alpha z\,\partial\phi(0)
% +i\bar\alpha \bar z\,\bar\partial\bar\phi(0) \\
% &\qquad\quad
% +\frac{i\alpha}{2}z^2\partial^2\phi(0)
% +\frac{i\bar\alpha}{2}\bar z^2\bar\partial^2\bar\phi(0)
% +\cdots
% \Big):\nn\\
% &=z^{\alpha\beta}\bar z^{\bar\alpha\bar\beta}
% :V_{\gamma,\bar\gamma}(0,0)e^Y:\nn\\
% Y&=i\alpha z\partial\phi(0)
% +i\bar\alpha\bar z\bar\partial\bar\phi(0)
% +\frac{i\alpha}{2}z^2\partial^2\phi(0)
% +\frac{i\bar\alpha}{2}\bar z^2\bar\partial^2\bar\phi(0)
% +\cdots .
% \end{aligned}
% \]
% Expanding $e^Y$ up to second order in derivatives we find
% \begin{align}
% e^Y
% &=
% 1
% +i\alpha z\,\partial\phi
% +i\bar\alpha \bar z\,\bar\partial\bar\phi
% +\frac{i\alpha}{2}z^2\partial^2\phi
% +\frac{i\bar\alpha}{2}\bar z^2\bar\partial^2\bar\phi\nn\\
% &\quad-\frac{\alpha^2}{2}z^2(\partial\phi)^2
% -\frac{\bar\alpha^2}{2}\bar z^2(\bar\partial\bar\phi)^2
% -\alpha\bar\alpha z\bar z\,\partial\phi\,\bar\partial\bar\phi
% +O(\text{level }3).
% \end{align}
% Combining everything, and defining $J=i \partial_z \varphi(z)$ and $\bar J=i \bar\partial_z \bar\varphi(z)$ we obtain
% \begin{align}
% V_{\alpha,\bar\alpha}(z,\bar z)V_{\beta,\bar\beta}(0,0)
% &=
% z^{\alpha\beta}\bar z^{\bar\alpha\bar\beta}
% \Big[
% V_{\gamma,\bar\gamma}(0,0)+:(\alpha zJ
% +\bar{\alpha}\bar{z} \bar{J})\,V_{\gamma,\bar\gamma}(0,0):\nn\\
% &\frac{1}{2}:\lb z^2(\alpha J'+\alpha^2 J^2)+\bar{z}^2(\bar{\alpha}\bar{J}'+\bar{\alpha}^2\bar{J}^2)\rb V_{\gamma,\bar\gamma}(0,0):\nn\\
% &+\alpha\bar{\alpha} z\bar{z}:J\bar{J} V_{\gamma,\bar\gamma}(0,0):
% +\cdots
% \Big].
% \end{align}
% In the case $\gamma=\alpha+\beta=0$ and $\bar{\gamma}=\bar{\alpha}+\bar{\beta}=0$, we have
% \begin{align}
% V_\alpha(z,\bar{z})V_{-\alpha}(0,0)&=(z\bar{z})^{-\alpha^2}\Big[ 1+\alpha :(z J+\bar{z}\bar{J})(0):\nn\\
% &\quad+\frac{\alpha}{2}:\lb z^2J+\bar{z}^2\bar{J}\rb(0):+\frac{\alpha^2}{2}:(z J+\bar{z}\bar{J})^2(0):+\cdots\Big]
% \end{align}
% We saw that the only operators that contribute for us have spin $h-\bar{h}$ is a multiple of $4$. Therefore, to this order the only operator that contributes to our correlator is 
% \begin{align}
% V_\alpha(z,\bar{z})V_{-\alpha}(0,0)=(z\bar{z})^{-\alpha^2}\lb 1+\alpha^2 \bar{z}\bar{z} (J\bar{J})(0,0)+\cdots\rb
% \end{align} 
% which has $h_p=\bar{h}_p=1$ and $C_\alpha=\alpha$. Therefore, $\Delta_L=2$ and
% \begin{align}
% S(\rho_\alpha\|y\rho_\alpha+(1-y)\rho_\beta)&=(\pi x)^4\frac{\sqrt{\pi}\Gamma(3)}{4\Gamma(7/2)}(1-y)^2(\alpha-\beta)^2+\cdots \nn\\
% &=\frac{4(\pi x)^4}{15}(1-y)^2(\alpha-\beta)^2+\cdots
% \end{align}
% At $y=1/2$ and $\alpha=0$ we find
% \begin{align}
%   S(\rho_0\|(\rho_0+\rho_\beta)/2)=\frac{(\pi x)^4\beta^2}{15}+\cdots
% \end{align}
% Note that this is $x^4$ order, not $x^2$ order. This is a consequence of the fact that the lowest primaries $J$ and $\bar{J}$ did not contribute because they are spin one. If instead of $V_{\alpha,\alpha}$ we considered chiral primaries, then $J$ would contribute at $x^2$ order \NL{Repeat the calculation in that case.}


% \subsection{Compact boson}\label{sec:compact-boson}

% For compact boson $\phi\sim \phi+2\pi R$ we have the momentum $m/R$ and winding modes $w R$ for $m,w\in\mathbb{Z}$, and left and right charges are
% \begin{align}
%   p_L=\frac{m}{R}+\frac{w}{2\pi R}, \qquad p_R=\frac{m}{R}-\frac{w}{2\pi R}.
% \end{align}
% The vertex operators are 
% \begin{align}
%   V_{m,w}(z,\bar{z})&=:e^{i(p_L\varphi(z)+p_R\bar\varphi(\bar{z}))}:,\qquad (h,\bar{h})=(p_L^2/2,p_R^2/2),\nn\\
%   \Delta&=h+\bar{h}=\frac{m^2}{R^2}+w^2 R^2,\qquad s=h-\bar{h}=m w.
% \end{align}
% The OPE is
% \begin{align}
%   V_{m,w}(z,\bar{z})V_{m',w'}(0,0)=z^{p_L p'_L}\bar{z}^{p_R p'_R}V_{m+m',w+w'}(0,0)\lb 1+\cdots\rb
% \end{align}
% We can use the Coulomb gas formalism to describe the operator content and OPE of all minimal models in terms of the OPE of vertex operators combined with screening charges. \NL{Do Ising model as an example}

% \subsection{Global symmetry}

% Symmetry resolved relative entropy

% In an exact energy eigenstate, the reduced density matrix is supposed to commute with the charge operator $Q$. There are superselection sectors, and the reduced state takes the form
% \begin{align}
%   &\rho=\oplus_q p_q \rho_q, \qquad \rho_q=\Pi_qr\rho_q \Pi_q,
% \end{align}
% where $\Pi_q$ is the projector onto the sector with charge $q$. 
% In the case of $U(1)$ current algebra, we have
% \begin{align}
%   &\Pi_q=\frac{1}{2\pi}\int d\alpha e^{-i\alpha Q}e^{i \pi q}
% \end{align}
% As explained in \cite{Capizzi} this project is given by the insertion of two vertex operators at the end points of the region. 



%   \subsection*{Operator Insertions on the Strip}

%   We perform the calculation on a strip of width $2\pi$, obtained via the
%   standard conformal mapping from the $n$-sheeted cylinder. The replica sheets are
%   indexed by $k = 0, \dots, n-1$.

%   The operator insertion points $t_{k}^{\pm}$ on the strip are:
%   \begin{equation}
%     t_{k}^{\pm}= \frac{2\pi k}{n}\pm \frac{\pi x}{n}
%   \end{equation}
%   where $x = l/R$ is the dimensionless interval size. For the correlator $\langle
%   \mathcal{O}_{\alpha}^{(q)}\mathcal{O}_{\beta}^{(p)}\rangle$ with $q=n-p$, the
%   operators are explicitly distributed as follows: On sheets $k=0, \dots, q-1$ ($\alpha$-sector)
%   we insert $V_{\alpha}$ at $t_{k}^{+}$ and $V_{-\alpha}$ at $t_{k}^{-}$, and on
%   sheets $k=q, \dots, n-1$ ($\beta$-sector) we insert $V_{\beta}$ at $t_{k}^{+}$
%   and $V_{-\beta}$ at $t_{k}^{-}$.

%   \subsection*{Special case}

%   For simplicity, focus on the special case $y=1/2$ so that we need to compute
%   \begin{eqnarray}
%     2^{-(n-1)}\frac{\tr(\rho_{0}(\rho_{0}+\rho_{\beta})^{n-1})}{\tr(\rho_{0}^{n})}&&=2^{-(n-1)}\Big(
%     1+(n-1)\frac{\tr(\rho_{0}^{n-1}\rho_{\beta})}{\tr(\rho_{0}^{n})}\nn\\
%     &&\quad+(n-1)\sum_{k=1}^{n-2}\frac{\tr(\rho_{0}^{k}\rho_{\beta}\rho_{0}^{n-1-k}\rho_{\beta})}{\tr(\rho_{0}^{n})}
%     \nn\\
%     &&\quad+(n-1)\sum_{k_1=1}^{n-3}\sum_{k_2=0}^{n-3-k_1}\frac{\tr(\rho_{0}^{k_1}\rho_{\beta}\rho_{0}^{k_2}\rho_{\beta}\rho_{0}^{n-3-k_1-k_2}\rho_{\beta})}{\tr(\rho_{0}^{n})}+\cdots\Big)
%   \end{eqnarray}
%   In this case, the first term $p=1$ gives
%   \begin{eqnarray}
%     &&\frac{\tr(\rho_{0}^{n-1}\rho_{\beta})}{\tr(\rho_{0}^{n})}=f(x)^{-\beta^2}\nn\\
%     &&f(x)\equiv (2\sin(\pi x/n)),
%   \end{eqnarray}
%   whereas the $p=2$ becomes
%   \begin{eqnarray}
%     &&\sum_{k=1}^{n-2}\frac{\tr(\rho_{0}^{k}\rho_{\beta}\rho_{0}^{n-1-k}\rho_{\beta})}{\tr(\rho_{0}^{n})}=f_1(x)^{-2\beta^2}
%     g_2(x;\beta)\nn\\
%     &&g_2(x;\beta) = \sum_{k=1}^{n-2} \left[
%     \frac{\sin^{2}\left(\frac{\pi k}{n}\right)}{\sin\left(\frac{\pi(k+x)}{n}\right)
%     \sin\left(\frac{\pi(k-x)}{n}\right)}
%     \right]^{\beta^2}\ .
%   \end{eqnarray}
%   For the $p=3$ term, we find
%   \begin{eqnarray}
%     \sum_{k_1=1}^{n-3}\sum_{k_2=0}^{n-3-k_1}\frac{\tr(\rho_{0}^{k_1}\rho_{\beta}\rho_{0}^{k_2}\rho_{\beta}\rho_{0}^{n-3-k_1-k_2}\rho_{\beta})}{\tr(\rho_{0}^{n})}=\cdots
%   \end{eqnarray}

% %  \section{General $x$}

  

%   \subsection*{Calculation of the correlator}

%   The correlation function for primary vertex operators $V_{Q}= :e^{i Q \phi}:$
%   is given by the product of distances raised to the power of the product of
%   their charges. The chordal distance on the strip is $d(t_{i}, t_{j}) = 2 \sin(\frac{t_{i}-
%   t_{j}}{2})$.
%   \begin{equation}
%     \langle \prod_{i}V_{Q_i}(t_{i}) \rangle \sim \prod_{i<j}\left[ 2 \sin\left( \frac{t_{i}-
%     t_{j}}{2}\right) \right]^{Q_i Q_j}
%   \end{equation}
%   \begin{itemize}
%     \item \textbf{Like charges} ($Q_{i}Q_{j}> 0$): The operators repel (regular
%       zero), leading to a positive exponent $+Q_{i}Q_{j}$.

%     \item \textbf{Opposite charges} ($Q_{i}Q_{j}< 0$): The operators attract (singular
%       pole), leading to a negative exponent $-|Q_{i}Q_{j}|$.
%   \end{itemize}

%   We have the following contributions:
%   \paragraph{1. Nearest Neighbor Interactions (Same sheet, same sector)}
%   These are the contractions between $t_{k}^{+}$ and $t_{k}^{-}$ on the same sheet.
%   \begin{itemize}
%     \item $\alpha$-sector ($q$ pairs): $(\alpha, -\alpha)$ separated by
%       $\Delta t = 2\pi x/n$.
%       \[
%         \implies \left( 2 \sin \frac{\pi x}{n}\right)^{-q \alpha^2}
%       \]

%     \item $\beta$-sector ($p$ pairs): $(\beta, -\beta)$ separated by
%       $\Delta t = 2\pi x/n$.
%       \[
%         \implies \left( 2 \sin \frac{\pi x}{n}\right)^{-p \beta^2}
%       \]
%   \end{itemize}
%   Putting the two together, we have
%   \begin{eqnarray}
%     \label{Nfunc} \mathcal{N}(q,p;x)=\left( 2 \sin \frac{\pi x}{n} \right)^{-(q\alpha^2+p\beta^2)}
%   \end{eqnarray}

%   \paragraph{2. Intra-sector Interactions (Different sheets, same sector)}
%   For the $\alpha$-sector, we sum over sheet separations $k = 1, \dots, q-1$. The
%   number of pairs at distance $k$ is $(q-k)$.
%   \begin{itemize}
%     \item Like charges $(\alpha, \alpha)$ and $(-\alpha, -\alpha)$:
%       \[
%         \text{Dist}= \frac{2\pi k}{n}\implies \text{Exp}= +\alpha^{2}\times 2
%       \]

%     \item Opposite charges $(\alpha, -\alpha)$: Two diagonals
%       $\frac{2\pi k}{n}\pm \frac{2\pi x}{n}$.
%       \[
%         \text{Dist}= \frac{2\pi (k\pm x)}{n}\implies \text{Exp}= -\alpha^{2}\text{
%         (each)}
%       \]
%   \end{itemize}
%   Collecting these terms for the $\alpha$ sector gives $[g(q;x)]^{\alpha^2}$:
%   \begin{eqnarray}
%     g(q;x) &= \prod_{k=1}^{q-1} \left(
%     \frac{\sin^{2}(\frac{\pi k}{n})}{\sin(\frac{\pi(k+x)}{n}) \sin(\frac{\pi(k-x)}{n})}
%     \right)^{q-k} \ .
%   \end{eqnarray}
%   The $\beta$ sector is identical with $q \to p$ and $\alpha \to \beta$, and we
%   have the overall contribution
%   \begin{eqnarray}
%     g(q;x)^{\alpha^2}g(p;x)^{\beta^2}\ .
%   \end{eqnarray}

%   \paragraph{3. Cross-sector Interactions}
%   Interactions between the $q$ $\alpha$-sheets and $p$ $\beta$-sheets yield
%   terms involving the product $\alpha \beta$. These collect into $[h(p,q;x)]^{\alpha\beta}$:
%   \begin{align}
%     \label{hfunct}h(q,p;x) & = \prod_{l=1}^{p}\prod_{k=1}^{q}\frac{\sin^{2}(\frac{\pi(k+l-1)}{n})}{\sin(\frac{\pi(k+l-1+x)}{n}) \sin(\frac{\pi(k+l-1-x)}{n})}
%   \end{align}

%   Putting all the pieces together, we find that the full correlator is
%   \begin{equation}
%     \label{eq:mixed_correlator}\langle \mathcal{O}_{\alpha}^{(q)}\mathcal{O}_{\beta}
%     ^{(p)}\rangle = \mathcal{N}(n-p,p;x) \left[ g(n-p;x) \right]^{\alpha^2}\left[
%     g(p;x) \right]^{\beta^2}\left[ h(n-p,p;x) \right]^{\alpha\beta}\ .
%   \end{equation}

%   We need to continue in $n$ analytically:
%   \begin{eqnarray}
%     \label{Fn} F(n)\equiv \text{tr}(\rho(\rho+\sigma)^{n-1}) &&= \sum_{p=0}^{n-1}
%     \binom{n-1}{p} \langle \mathcal{O}_{\alpha}^{(q)} \mathcal{O}_{\beta}^{(p)} \rangle
%   \end{eqnarray}

%   \subsection{Analytic continuation in $n$}

%   We consider the relative entropy $S(\rho_{\alpha}|| \sigma)$ where
%   $\sigma = y \rho_{\alpha}+ (1-y)\rho_{\beta}$ for excited coherent states in a
%   1+1D free boson CFT. The replica trick requires evaluating:
%   \begin{equation}
%     \text{tr}(\rho_{\alpha}\sigma^{n-1}) = \sum_{p=0}^{n-1}\binom{n-1}{p}y^{n-1-p}
%     (1-y)^{p}C_{n-p, p}
%   \end{equation}
%   where
%   $C_{n-p, p}= \langle \mathcal{O}_{\alpha}^{(n-p)}\mathcal{O}_{\beta}^{(p)}\rangle$
%   is the mixed correlator on the $n$-sheeted replica surface.

%   \paragraph{Step 1 Log-Sum expansion of the correlator:}
%   The correlator is a product of geometric functions $g$ and $h$. We focus
%   on the intra-sector term $\log g(p;x)$. For integer $p$, it is defined as:
%   \begin{equation}
%     \log g(p;x) = \sum_{k=1}^{p-1}(p-k) \left[ 2 \log d\left(\frac{k}{n}\right) -
%     \log d\left(\frac{k+x}{n}\right) - \log d\left(\frac{k-x}{n}\right) \right]
%   \end{equation}
%   where $d(y) = 2\sin(\pi y)$ is the distance between operator insertions. We
%   utilize the Fourier expansion for $\log \sin(\pi y)$:
%   \begin{equation}
%     \log \sin(\pi y) = -\log 2 - \sum_{m=1}^{\infty}\frac{\cos(2\pi m y)}{m}
%   \end{equation}
%   Substituting this into the potential $\mathcal{V}(k,x)$ inside the sum, the
%   constants cancel, yielding:
%   \begin{equation}
%     \mathcal{V}(k,x) = -2 \sum_{m=1}^{\infty}\frac{1}{m}\left[ 1 - \cos\left(\frac{2\pi
%     m x}{n}\right) \right] \cos\left(\frac{2\pi m k}{n}\right)
%   \end{equation}
%   The expression above has the advantage that it has separated the contribution
%   of $k$ from those of $x$, and we can now perform the sum over $k$.

%   \paragraph{Step 2 Sum over $k$:}
%   Following Section 6 of Calabrese, Cardy, and Tonni (2011), we perform the sum over
%   $k$ for each mode $m$:
%   \begin{equation}
%     \log g(p;x) = -2 \sum_{m=1}^{\infty}\frac{1 - \cos(2\pi m x/n)}{m}\sum_{k=1}^{p-1}
%     (p-k) \cos\left(\frac{2\pi m k}{n}\right)
%   \end{equation}
%   The inner sum $S_{m}(p) = \sum_{k=1}^{p-1}(p-k) \cos(k\theta_{m})$ is
%   evaluated using geometric series. For integer $p$ and $\theta_{m}=2\pi m/n$:
%   \begin{equation}
%     S_{m}(p) = \text{Re}\left[ \frac{p e^{i\theta_m}}{1-e^{i\theta_m}}- \frac{e^{i\theta_m}(1-e^{ip\theta_m})}{(1-e^{i\theta_m})^{2}}
%     \right]=\frac{-p}{2}+\frac{(1-\cos(p\theta_{m}))}{4\sin^{2}(\theta_{m}/2)},
%   \end{equation}
%   Therefore, we find
%   \begin{eqnarray}
%     \log g(p;x)=\sum_{m=1}^\infty \frac{1-\cos(\frac{2\pi m x}{n})}{m}\lb p-\frac{\sin^{2}(\pi
%     p m/n)}{\sin^{2}(\pi m/n)}\rb
%   \end{eqnarray}

%   Define the mode kernels
%   \begin{equation}
%     L(n,x):=\sum_{m=1}^{\infty}\frac{1-\cos(2\pi m x/n)}{m}, \qquad K_{m}(n,x):=
%     \frac{1-\cos(2\pi m x/n)}{2m\,\sin^{2}(\pi m/n)}. \label{eq:L-Km}
%   \end{equation}
%   to write
%   \begin{equation}
%     \log g(p;x)= p\,L(n,x)\;-\;\sum_{m=1}^{\infty}K_{m}(n,x)\,\bigl(1-\cos(p\theta
%     _{m})\bigr). \label{eq:logg-final}
%   \end{equation}
%   All $p$-independent pieces can be absorbed into $\Omega_{n}^{\rm (rest)}(x)$
%   later.

%   A technical note to keep in mind is that to make every rearrangement fully
%   justified, one may define all mode sums with a common Abel factor $e^{-\eps m}$
%   in the summand and take $\eps\to 0^{+}$ at the end. This guarantees absolute
%   and uniform convergence in intermediate steps.

%   Next, we move to the function $h(q,p;x)$ in (\ref{hfunct}) which includes a double
%   product over $(k,l)$ with $1\le k\le q$, $1\le l\le p$. Taking logs and using the
%   same log-sum representation as in we did for the $g$ function yields
%   \begin{equation}
%     \log h(q,p;x)=\sum_{r=1}^{n-1}w_{r}(p,q)\,V(r,x), \label{eq:logh-weighted}
%   \end{equation}
%   where $w_{r}(p,q)$ counts the number of pairs $(k,l)$ such that $k+l-1=r$.

%   It is convenient to evaluate the weighted cosine sum via generating functions.
%   Let $\theta$ be arbitrary and note
%   \begin{equation}
%     \sum_{r=1}^{n-1}w_{r}(p,q)e^{\ii r\theta}=\sum_{k=1}^{q}\sum_{l=1}^{p}e^{\ii(k+l-1)\theta}
%     =e^{-\ii\theta}\Big(\sum_{k=1}^{q}e^{\ii k\theta}\Big)\Big(\sum_{l=1}^{p}e^{\ii
%     l\theta}\Big). \label{eq:w-gen}
%   \end{equation}
%   Now specialize to the case $q=n-p$ and $\theta=\theta_{m}=2\pi m/n$. Using
%   \begin{equation}
%   \sum_{k=1}^{n-p}e^{\ii k\theta_m}=-e^{-\ii p\theta_m/2}\,\frac{\sin(p\theta_{m}/2)}{\sin(\theta_{m}/2)}
%   \quad\text{and}\quad
%   \sum_{l=1}^{p}e^{\ii l\theta_m}=e^{\ii(p+1)\theta_m/2}\,\frac{\sin(p\theta_{m}/2)}{\sin(\theta_{m}/2)},
%   \end{equation}
%   the right-hand side of \eqref{eq:w-gen} becomes real and equals
%   \begin{equation}
%     \sum_{r=1}^{n-1}w_{r}(p,n-p)\cos(r\theta_{m}) = -\left(\frac{\sin(p\theta_{m}/2)}{\sin(\theta_{m}/2)}
%     \right)^{2}= -\frac{1-\cos(p\theta_{m})}{2\sin^{2}(\theta_{m}/2)}. \label{eq:wcos-eval}
%   \end{equation}
%   Substituting \eqref{eq:wcos-eval} into \eqref{eq:logh-weighted} gives the exact
%   mode form
%   \begin{equation}
%     \log h(n-p,p;x)= 2\sum_{m=1}^{\infty}K_{m}(n,x)\,\bigl(1-\cos(p\theta_{m})\bigr
%     ), \label{eq:logh-final}
%   \end{equation}
%   with the \emph{same} $K_{m}(n,x)$ as in \eqref{eq:L-Km}.

%   We are now ready to put together all the pieces back into $C_{n-p,p}$. Using
%   \eqref{eq:logg-final} and \eqref{eq:logh-final}, we find
%   \begin{align}
%     \alpha^{2}\log g(q;x)+\beta^{2}\log g(p;x)+\alpha\beta\log h(q,p;x) & =(\alpha^{2}q+\beta^{2}p)\,L(n,x) \nonumber                                                             \\
%                                                                         & \quad-(\alpha-\beta)^{2}\sum_{m=1}^{\infty}K_{m}(n,x)\,\bigl(1-\cos(p\theta_{m})\bigr) \label{eq:combo}
%   \end{align}
%   where we have used a similar log-sum expansion for the function $\mathcal{N}$.
%   Absorbing all $p$-independent constants (including the $-(\alpha-\beta)^{2}\sum
%   _{m}K_{m}$ term) into $\Omega_{n}(x)$, we obtain
%   \begin{equation}
%     C_{n-p,p}(x) =\Omega_{n}(x)\,\Lambda_{\alpha}^{\,n-p}\,\Lambda_{\beta}^{\,p}\,
%     \exp\!\left[\sum_{m=1}^{\infty}a_{m}(n,x)\cos(p\theta_{m})\right], \label{eq:C-expcos}
%   \end{equation}
%   where
%   \begin{equation}
%     a_{m}(n,x)=(\alpha-\beta)^{2}\,K_{m}(n,x) =(\alpha-\beta)^{2}\,\frac{1-\cos(2\pi
%     m x/n)}{2m\,\sin^{2}(\pi m/n)}\ . \label{eq:am}
%   \end{equation}

%   \subsection{Fourier/Bessel expansion and definition of the $R$-modes}
%   Let $\theta:=\frac{2\pi p}{n}$ so that $\cos(p\theta_{m})=\cos(m\theta)$. Using
%   $e^{a\cos(m\theta)}=\sum_{r\in\mathbb{Z}}I_{r}(a)\,e^{\ii r m\theta}$ (modified
%   Bessel functions), we expand
%   \begin{equation}
%     \exp\!\left[\sum_{m\ge 1}a_{m}\cos(m\theta)\right] =\sum_{R\in\mathbb{Z}}W_{R}
%     (n,x;\alpha,\beta)\,e^{\ii R\theta}, \label{eq:WR-Fourier}
%   \end{equation}
%   where the coefficients are the (infinite) Bessel convolution
%   \begin{equation}
%     W_{R}(n,x;\alpha,\beta)= \sum_{\{r_m\}\,:\,\sum_{m\ge1} m r_m=R}\;\prod_{m\ge1}
%     I_{r_m}\!\big(a_{m}(n,x)\big), \qquad W_{-R}=W_{R}. \label{eq:WR-def}
%   \end{equation}
%   This allows us to rewrite
%   \begin{equation}
%     C_{n-p,p}(x) =\Omega_{n}\,\Lambda_{\alpha}^{\,n-p}\Lambda_{\beta}^{\,p}\; \sum
%     _{R\in\mathbb{Z}}W_{R}(n)\,e^{\ii \frac{2\pi R}{n}p}. \label{eq:C-R}
%   \end{equation}

%   Note that with an Abel regulator in the defining mode sums, the function of
%   $\theta$ in \eqref{eq:WR-Fourier} is real-analytic in a complex strip, which
%   implies exponential decay of $W_{R}$ and in particular
%   $\sum_{R}|W_{R}|<\infty$.

%   The expression in (\ref{eq:C-R}) has the advantage that if we can swap the sum
%   over $R$ with the sum over $p$, we can perform the $p$-sum explicitly. That is
%   our next step. {\bf I believe swapping the two sums can be justified using Abel regulators but I am not entirely sure!}

%   Plugging these back in \ref{Fn} and swapping the order of summation, we obtain
%   \begin{align}
%     F(n) & =\Omega_{n}\,\Lambda_{\alpha}\sum_{R\in\mathbb{Z}}W_{R}(n) \sum_{p=0}^{n-1}\binom{n-1}{p}\, (y\Lambda_{\alpha})^{n-1-p}\, \bigl((1-y)\Lambda_{\beta}e^{\ii2\pi R/n}\bigr)^{p}\nonumber \\
%          & = \Omega_{n}\,\Lambda_{\alpha}\sum_{R\in\mathbb{Z}}W_{R}(n)\, \Bigl(y\Lambda_{\alpha}+(1-y)\Lambda_{\beta}e^{\ii2\pi R/n}\Bigr)^{n-1}\ . \label{eq:F-R}
%   \end{align}

%   Next, we would like to argue that only the $R=0$ term in the expression above contrinutes
%   to the $n\to 1$ limit. To establish this, define $\eps=n-1$ and
%   \begin{equation}
%     A_{R}(n):=y\Lambda_{\alpha}+(1-y)\Lambda_{\beta}e^{\ii2\pi R/n},\qquad A_{0}:
%     =y\Lambda_{\alpha}+(1-y)\Lambda_{\beta}>0.
%   \end{equation}
%   Then \eqref{eq:F-R} reads
%   \begin{equation}
%     F(n)=\Omega_{n}\Lambda_{\alpha}\sum_{R\in\mathbb{Z}}W_{R}(n)\,A_{R}(n)^{\eps}
%     . \label{eq:F-eps}
%   \end{equation}

%   With an Abel regulator $W_{R}$ decays exponentially and has finite moments.
%   Expand the phase for $n=1+\eps$:
%   \begin{equation}
%     e^{\ii2\pi R/n}=e^{\ii2\pi R/(1+\eps)}=\exp\!\bigl(-\ii2\pi R\eps+O(R\eps^{2}
%     )\bigr),
%   \end{equation}
%   so $A_{R}(n)=A_{0}+\delta_{R}$ with $\delta_{R}=O(R\eps)$ for fixed $R$ and $\delta
%   _{-R}=\overline{\delta_R}$. Using
%   \begin{equation}
%     A_{R}^{\eps}=\exp\!\bigl(\eps\log(A_{0}+\delta_{R})\bigr)=A_{0}^{\eps}\left[1
%     +\eps\frac{\delta_{R}}{A_{0}}+O(\eps|\delta_{R}|^{2})\right],
%   \end{equation}
%   and the evenness $W_{-R}=W_{R}$, the potentially dangerous linear term $\sum_{R}
%   W_{R}\,\delta_{R}$ cancels at $O(\eps)$ (its leading part is odd/purely
%   imaginary), leaving only $O(\eps^{2})$ corrections after summing over $R$. Consequently,
%   after absorbing $\sum_{R}W_{R}$ into $\Omega_{n}$ (a $p$-independent normalization),
%   \begin{equation}
%     F(1+\eps)=\Omega_{1+\eps}\Lambda_{\alpha}\,A_{0}^{\eps}\bigl(1+O(\eps^{2})\bigr
%     ). \label{eq:F-asymp}
%   \end{equation}
%   Thus the $R\neq 0$ modes do not affect $\partial_{n}\log F(n)$ at $n=1$.

%   The replica prescription for the relative entropy is
%   \begin{equation}
%     S(\rho_{\alpha}\|\sigma_{y})= \left.\partial_{n}\Bigl(\log \tr(\rho_{\alpha}^{n}
%     )-\log \tr(\rho_{\alpha}\sigma_{y}^{n-1})\Bigr)\right|_{n=1}. \label{eq:S-replica}
%   \end{equation}
%   In the same framework one has $\tr(\rho_{\alpha}^{n})=\widetilde\Omega_{n}\,\Lambda
%   _{\alpha}^{\,n}$ so that $\partial_{n}\log \tr(\rho_{\alpha}^{n})|_{n=1}=\log\Lambda
%   _{\alpha}$ (all $\widetilde \Omega_{n}$ terms cancel against $\Omega_{n}$ in
%   the difference). Using \eqref{eq:F-asymp},
%   \begin{equation}
%     \left.\partial_{n}\log \tr(\rho_{\alpha}\sigma_{y}^{n-1})\right|_{n=1}=\log A
%     _{0}=\log\!\bigl(y\Lambda_{\alpha}+(1-y)\Lambda_{\beta}\bigr).
%   \end{equation}
%   Therefore
%   \begin{equation}
%     S(\rho_{\alpha}\|\sigma_{y})=\log\Lambda_{\alpha}-\log\!\bigl(y\Lambda_{\alpha}
%     +(1-y)\Lambda_{\beta}\bigr) = -\log\!\left[y+(1-y)\frac{\Lambda_{\beta}}{\Lambda_{\alpha}}
%     \right]. \label{eq:S-Lambda}
%   \end{equation}

%   For coherent-state excitations in the free boson CFT one may identify
%   \begin{equation}
%     \frac{\Lambda_{\beta}}{\Lambda_{\alpha}}=e^{-S(\alpha\|\beta)}, \qquad S(\alpha
%     \|\beta)=(\alpha-\beta)^{2}K(x), \qquad K(x)=1-\pi x\cot(\pi x). \label{eq:Salpha}
%   \end{equation}
%   Substituting \eqref{eq:Salpha} into \eqref{eq:S-Lambda} gives the final answer
%   \begin{equation}
%     \boxed{ S\!\left(\rho_\alpha\,\Big\|\,y\rho_\alpha+(1-y)\rho_\beta\right) = -\log\!\left( y+(1-y)\exp\!\big[-(\alpha-\beta)^2(1-\pi x\cot(\pi x))\big] \right). }
%     \label{eq:S-final}
%   \end{equation}

% \subsection{Current algebra: Symmetry resolved relative entropy}



%   \section{Consistency tests}
%   \label{sec:tests}

%   Let $s:=(\alpha-\beta)^{2}K(x)\ge 0$ and $a:=e^{-s}\in(0,1]$ so that
%   \begin{eqnarray}
%     S(y)\equiv S\!\left(\rho_\alpha\,\Big\|\,y\rho_\alpha+(1-y)\rho_\beta\right)=-\log\lb
%     y+(1-y)a\rb=-\log(a+(1-y)a)
%   \end{eqnarray}
%   We are gonna test the following limits:
%   \begin{itemize}
%     \item \textbf{Identical states:}
%       $\alpha=\beta\Rightarrow s=0\Rightarrow a=1\Rightarrow S(y)=0$ as expected.

%     \item \textbf{No mixing:} $S(1)\to 0$ as expected and
%       $S(0)=S(\alpha\|\beta)$ as expected. Therefore, we can write
%       \begin{eqnarray}
%         e^{-S(y)}=y e^{-S(1)}+(1-y)e^{-S(0)}\ .
%       \end{eqnarray}

%     \item \textbf{Much heavier state:}
%       $s\to\infty\Rightarrow a\to 0\Rightarrow S_{y}\to -\log y$.
%   \end{itemize}

%   Next, we test the following properties:

%   \begin{itemize}
%     \item \textbf{Positivity:} Since $y\le a+y(1-a)\le 1$, we have $0\le S(y)\le
%       -\log y$ as required fo r relative entropy.

%     \item \textbf{Convexity in $y$:} Taking derivatives with respect to $y$
%       \begin{equation}
%         S'(y)= -\frac{1-a}{a+y(1-a)}\le 0, \qquad S''(y)=\frac{(1-a)^{2}}{(a+y(1-a))^{2}}
%         \ge 0.
%       \end{equation}
%       Thus $S(y)$ is decreasing and convex in the mixing parameter $y$ (strictly
%       convex if $s>0$), consistent with the general theorem that $S(\rho\|\sigma)$
%       is convex in $\sigma$.

%     \item \textbf{Monotonicity and concavity in $s$:} Treating $S$ as a function
%       of $s$:
%       \begin{equation}
%         \frac{\partial S}{\partial s}=\frac{(1-y)e^{-s}}{y+(1-y)e^{-s}}\ge 0, \qquad
%         \frac{\partial^{2}S}{\partial s^{2}}=-\frac{y(1-y)e^{-s}}{(y+(1-y)e^{-s})^{2}}
%         \le 0.
%       \end{equation}
%       So $S$ increases with $s$ but is concave in $s$, consistent with
%       saturation at $-\log y$. A direct derivative gives the derivative of relative
%       entropy in region size
%       \begin{equation}
%         K'(x)=\frac{\pi}{\sin^{2}(\pi x)}\left(\pi x-\frac{1}{2}\sin(2\pi x)\right
%         )>0 \qquad \text{for }x\in(0,1),
%       \end{equation}
%       since $\sin u<u$ for $u>0$. Thus $K(x)$ (and therefore $S$) is increasing
%       in $x$ for $\alpha\neq\beta$, as expected from data processing.
%   \end{itemize}

%   Finally, we note that in the small interval limit $x\ll 1$ we have
%   \begin{equation}
%     K(x)=1-\pi x\cot(\pi x)=\frac{\pi^{2}}{3}x^{2}+\frac{\pi^{4}}{45}x^{4}+O(x^{6}
%     ).
%   \end{equation}
%   Hence
%   \begin{equation}
%     S(y)= (1-y)\,(\alpha-\beta)^{2}\frac{\pi^{2}}{3}x^{2}+O(x^{4}).
%   \end{equation}
%   The relative entropy vanishes as $x^{2}$ for small intervals, as expected.

%   \subsection*{Alternative analytic continuation}
%   \label{sec:small-n1} An alternative analytic continuation follows from observing
%   that since we are interested in small $n-1$, and the sum over
%   $p=1,\cdots , n-1$, we linearize our functions in $p$ and discard higher order
%   terms:
%   \begin{equation}
%     S_{m}(p)= -\frac{p}{2}+\frac{1-\cos(p\theta_{m})}{4\sin^{2}(\theta_{m}/2)}= -
%     \frac{p}{2}+O(p^{2}). \label{eq:Sm-smallp}
%   \end{equation}
%   Expanding the $g$ function in $n=1+\epsilon$ results in
%   \begin{equation}
%     \log g(p;x)=p\,L(n,x)+O(p^{2})=p\,L(n,x)+O(\eps^{2}),
%   \end{equation}
%   uniformly over the relevant range $p\in\{0,\dots,n-1\}$. This reproduces the linear-in-$p$
%   structure that one would obtain by keeping only the ``smooth'' term $S_{m}(p)\approx
%   -p/2$ and effectively discarding the oscillatory $\cos(p\theta_{m})$ piece at this
%   order. While this results in a different analytic continuation than what we
%   computed above, the $n\to 1$ limit gives the same result, and the same final
%   expression \eqref{eq:S-final} for relative entropy.

  %------------------------------------------------------------
  % Appendix-style writeup: pole-matching analytic continuation
  %------------------------------------------------------------

  \section*{Pole-matching analytic continuation of $g(p;x)$ and $h(n-p,p;x)$}

  Here, we reformulate the analytic continuation of the building blocks $g(p;x)$
  and $h(n-p,p;x)$ using the method used in the appendix of 
  % \texttt{arXiv:1508.03506}
  \cite{Lashkari_2016}.
  The idea is to this: To find the analytic cotinuation of a meromorphic function
  $f_{1}(z)$, you find another entire function with the same poles and zeros
  $f_{2}(z)$ such that $f_{1}(z)/f_{z}(2)$ is bounded. Then, by Liouville's
  theorem it must be a constant $C$. This means that $C f_{2}(z)$ is the analytic
  continuation of $f_{1}(z)$.

  \subsection*{The basic building block and the definitions}

  Introduce
  \begin{equation}
    f_{n,m}(x) :=\frac{\sin^{2}(\pi m/n)}{\sin\!\bigl(\pi(m+x)/n\bigr)\,\sin\!\bigl(\pi(m-x)/n\bigr)}
    \,, \qquad m\in\mathbb{Z}.
  \end{equation}
  For integer $n\ge 2$ and $p\in\{1,\dots,n-1\}$, the functions appearing in the
  correlator are
  \begin{align}
    g(p;x)   & =\prod_{m=1}^{p-1}f_{n,m}(x)^{\,p-m}, \label{eq:g-def-poleapp}                                                                  \\[2pt]
    h(q,p;x) & =\prod_{l=1}^{p}\prod_{k=1}^{q}f_{n,k+l-1}(x) =\prod_{r=1}^{n-1}f_{n,r}(x)^{\,w_r(p,q)}, \qquad q=n-p, \label{eq:h-def-poleapp}
  \end{align}
  where $w_{r}(p,q)$ counts the number of pairs $(k,l)$ such that $k+l-1=r$.

  \subsection*{Poles of $g(p;x)$ as a function of $x$}

  From \eqref{eq:g-def-poleapp}, all $x$-dependence sits in the denominator
  sines. For fixed $m$, the factor $\sin(\pi(m\pm x)/n)$ vanishes iff
  $m\pm x=\ell n$ for some $\ell\in\mathbb{Z}$, i.e.
  \begin{equation}
    x=\pm m+\ell n.
  \end{equation}
  Hence the pole set of $g(p;x)$ is
  \begin{equation}
    \boxed{ \mathrm{Pole}(g)=\Bigl\{\,x=\pm m+\ell n\ \Big|\ m=1,\dots,p-1,\ \ell\in\mathbb{Z}\Bigr\}. }
  \end{equation}
  Moreover, each vanishing sine gives a simple pole, and the factor $f_{n,m}(x)$
  is raised to power $(p-m)$. Therefore the \emph{pole order} of $g(p;x)$ at
  $x=\pm m+\ell n$ equals $(p-m)$:
  \begin{equation}
    \boxed{ \mathrm{ord}_{x=\pm m+\ell n}\, g(p;x)=p-m. }
  \end{equation}
  Finally, $g(p;x)$ has no zeros as a function of $x$, since the numerator $\sin^{2}
  (\pi m/n)$ is independent of $x$.

  \subsection*{A multiple-sine candidate and its divisor}

  We now introduce a candidate $g_{\rm MS}(p;x)$ built from multiple sine
  functions. Let $S_{2}(z|1,n)$ and $S_{3}(z|1,1,n)$ denote the double and triple
  sine functions (in a standard normalization), characterized by the shift
  relations
  \begin{equation}
    \frac{S_{2}(z+1\,|\,1,n)}{S_{2}(z\,|\,1,n)}=2\sin\!\Bigl(\frac{\pi z}{n}\Bigr
    ), \qquad \frac{S_{3}(z+1\,|\,1,1,n)}{S_{3}(z\,|\,1,1,n)}=\frac{1}{S_{2}(z\,|\,1,n)}
    . \label{eq:S2S3-shifts}
  \end{equation}
  Define the triangular-product package
  \begin{equation}
    T_{p}(z):= \frac{S_{3}(z+2\,|\,1,1,n)^{\,p}}{S_{3}(z+p+1\,|\,1,1,n)\,S_{3}(z+1\,|\,1,1,n)^{\,p-1}}
    , \label{eq:T-def}
  \end{equation}
  and the candidate analytic continuation
  \begin{equation}
    g_{\rm MS}(p;x):=\frac{T_{p}(0)^{2}}{T_{p}(x)\,T_{p}(-x)}. \label{eq:gMS-def}
  \end{equation}

  \paragraph{Claim (telescoping to the finite sine product).}
  By iterating the shift relations \eqref{eq:S2S3-shifts}, one checks that
  $T_{p}(z)$ is proportional to the triangular sine product
  \begin{equation}
    T_{p}(z)\ \propto\ \prod_{m=1}^{p-1}\sin\!\Bigl(\frac{\pi(m+z)}{n}\Bigr)^{p-m}
    , \label{eq:T-telescope}
  \end{equation}
  where the proportionality factor is independent of $z$ (it may depend on $n,p$).
  Consequently, \eqref{eq:gMS-def} implies that for integer $n$,
  $g_{\rm MS}(p;x)$ reproduces the same pole lattice and the same pole orders as
  $g(p;x)$.

  \paragraph{Zeros/poles of $g_{\rm MS}(p;x)$.}
  From \eqref{eq:T-telescope}, $T_{p}(z)$ has zeros whenever $z=\ell n-m$ for $m=
  1,\dots,p-1$, with zero order $(p-m)$. Therefore $g_{\rm MS}(p;x)$ has poles whenever
  $T_{p}(x)=0$ or $T_{p}(-x)=0$, i.e.
  \begin{equation}
    x=\ell n-m\quad \text{or}\quad x=-(\ell n-m)=m-\ell n \quad\Longleftrightarrow
    \quad x=\pm m+\ell n,
  \end{equation}
  and the pole order at each such point equals $(p-m)$. In particular,
  \begin{equation}
    \boxed{ \mathrm{Pole}(g_{\rm MS})=\mathrm{Pole}(g),\qquad \mathrm{ord}_{x=\pm m+\ell n}\,g_{\rm MS}(p;x)=p-m. }
  \end{equation}
  Like $g$, the function $g_{\rm MS}$ has no zeros as a function of $x$.

  \subsection*{The ratio has no poles and is bounded}

  Define the ratio
  \begin{equation}
    R(x):=\frac{g_{\rm MS}(p;x)}{g(p;x)}.
  \end{equation}
  Since $g_{\rm MS}$ and $g$ have the same poles with the same orders and no zeros,
  all singularities cancel:
  \begin{equation}
    \boxed{R(x)\ \text{is entire in }x.}
  \end{equation}
  Moreover, both $g$ and $g_{\rm MS}$ are even and $n$-periodic in $x$, hence so
  is $R$:
  \begin{equation}
    R(-x)=R(x),\qquad R(x+n)=R(x). \label{eq:R-periodic}
  \end{equation}

  To argue boundedness, it suffices (by periodicity) to bound $R$ on a single
  vertical strip $\{x=u+\ii v: u\in[0,n]\}$. On such a strip, the original
  $g(p;x)$ obeys an exponential-type estimate as $|v|\to\infty$: for each fixed $m$,
  \begin{equation}
    \Bigl|\sin\!\Bigl(\frac{\pi(m\pm(u+\ii v))}{n}\Bigr)\Bigr| \sim \frac{1}{2}e^{\pi
    |v|/n}, \qquad |v|\to\infty,
  \end{equation}
  so
  \begin{equation}
    |g(p;u+\ii v)| \sim C\,\exp\!\Big[-\frac{2\pi|v|}{n}\sum_{m=1}^{p-1}(p-m)\Big] = C\,\exp\!\Bigl[-\frac{\pi p(p-1)}{n}|v|\Bigr], \qquad |v|\to\infty, \label{eq:g-asymp}
  \end{equation}
  with $C$ independent of $v$ (uniformly for $u$ in compact sets away from poles).

  On the other hand, the multiple sine functions $S_{2},S_{3}$ have well-known large-imaginary
  asymptotics of the form
  \begin{equation}
    \log S_{r}(u+\ii v\,|\,\boldsymbol{\omega}) = \pm \frac{\pi}{r!\,\omega_{1}\cdots
    \omega_{r}}\,|v|^{r}\ +\ \text{lower powers in }|v| \ +\ O(1), \qquad |v|\to
    \infty,
  \end{equation}
  so that in ratios such as \eqref{eq:T-def} the leading polynomial growth
  cancels, and $T_{p}(z)$ has exponential type matching the finite sine product in
  \eqref{eq:T-telescope}. Consequently $g_{\rm MS}(p;u+\ii v)$ has the same exponential
  type as $g(p;u+\ii v)$, and hence their ratio $R(u+\ii v)$ remains bounded as
  $|v|\to\infty$ uniformly in $u\in[0,n]$. Together with continuity on the compact
  region $\{u\in[0,n],\,|v|\le V\}$, we conclude:
  \begin{equation}
    \boxed{R(x)\ \text{is bounded on }\mathbb{C}.}
  \end{equation}

  Since $R$ is entire and bounded, Liouville's theorem implies $R$ is constant.
  Finally, $g(p;0)=1$ by definition and $g_{\rm MS}(p;0)=1$ by \eqref{eq:gMS-def},
  so $R(0)=1$ and thus
  \begin{equation}
    \boxed{g_{\rm MS}(p;x)=g(p;x)\quad \text{for all }x\ \text{(at integer }n\text{).}}
  \end{equation}
  This is precisely the ``same poles + bounded ratio'' strategy.

  \subsection*{Corollary: $h(n-p,p;x)$ from a full-product identity}

  Using \eqref{eq:h-def-poleapp} with $q=n-p$ and the explicit multiplicities
  \begin{equation}
    w_{r}(p,n-p)=
    \begin{cases}
      r,   & 1\le r\le p,  \\
      p,   & p<r\le n-p,   \\
      n-r, & n-p<r\le n-1,
    \end{cases}
  \end{equation}
  one may rewrite $h(n-p,p;x)$ as
  \begin{equation}
    h(n-p,p;x)=\Bigl(\prod_{r=1}^{n-1}f_{n,r}(x)\Bigr)^{p}\,g(p;x)^{-2}. \label{eq:h-via-fullprod}
  \end{equation}
  The appendix identity of 
  % \texttt{arXiv:1508.03506} 
  \cite{Lashkari_2016}
  proves
  \begin{equation}
    \prod_{r=1}^{n-1}f_{n,r}(x)=\left(\frac{n\sin(\pi x/n)}{\sin(\pi x)}\right)^{2}
    ,
  \end{equation}
  again by the same poles/bounded ratio argument. Substituting into \eqref{eq:h-via-fullprod}
  yields
  \begin{equation}
    \boxed{ h(n-p,p;x) = \left(\frac{n\sin(\pi x/n)}{\sin(\pi x)}\right)^{2p}\,g(p;x)^{-2}, }
  \end{equation}
  so the analytic continuation of $h$ follows once that of $g$ is fixed by the
  pole-matching method.

  \section{Information-theoretic interpretation}

  \appendix
\section{Analytic continuation of the replica sum with a phase}
\label{app:analytic-continuation}

In this appendix we derive the analytic continuation in the replica index
of the sum
\begin{align}
\mu(n)&=\sum_{p=1}^{n-1}\frac{(n-p)e^{-2\pi i \kappa_n p}}{\sin^{2\Delta}\!\left(\frac{\pi p}{n}\right)}\nn\\
e^{-2\pi i \kappa_n}&:=r_n(x):=(y^{-1}-1)D_n(1-x)=(y^{-1}-1)\lb\frac{n\sin(\pi (1-x)/n)}{\sin(\pi (1-x))} \rb^{4(h_\alpha-h_\beta)},
\end{align}
and compute its derivative at $n=1$. First note that for all positive integer $n$ the function $D_n(1-x)\geq 1$, which means that for $h_\beta\geq h_\alpha$ and $0\leq y\leq 1/2$ the function $r_n(x)$ is real, and bounded above by one:
\begin{align}
  \forall\: 0\leq y\leq 1/2, h_\alpha\leq h_\beta: \qquad 0\leq r_n(x)\leq 1,   
\end{align}
We are going to focus on this case. Redefining $p\to n-p$ we have
\begin{align}
\mu(n)=r_n(x)^{n}\sum_{p=1}^{n-1}\frac{p e^{2\pi i \kappa_n p}}{\sin^{2\Delta}\left(\frac{\pi p}{n}\right)}.
\end{align}


Our derivation of the analytic continuation $\mu(1)$ and $\mu'(1)$ in this limit follows closely the
method used by Cardy, Tonni and collaborators, based on Fourier analysis
and Carlson's theorem.

\subsection{Fourier coefficients of the sine kernel}

Define the periodic kernel
\begin{equation}
f(x)=\sin^{-2\Delta}(\pi x).
\end{equation}
Its Fourier coefficients are
\begin{equation}
f_\Delta(m)=\int_0^1 dx\,\frac{e^{-2\pi i m x}}{\sin^{2\Delta}(\pi x)},
\qquad m\in\mathbb Z .
\end{equation}
Using the standard beta-function integral one finds
\begin{equation}
f_\Delta(m)=
\frac{2^{2\Delta}\Gamma(1-2\Delta)e^{-i\pi m}}
{\Gamma(1-\Delta+m)\Gamma(1-\Delta-m)} .
\end{equation}
For analytic continuation it is convenient to introduce a holomorphic
interpolation of the positive Fourier coefficients to complex $m$,
\begin{equation}\label{eq: Fourier coeff of csc^2Delta}
g_\Delta(m)=c_\Delta\,\frac{\Gamma(m+\Delta)}{\Gamma(m+1-\Delta)},
\qquad
c_\Delta=\frac{2^{2\Delta}\Gamma(1-2\Delta)\sin(\pi\Delta)}{\pi}.
\end{equation}
This function reproduces the discrete data
\(
g_\Delta(m)=f_\Delta(m)
\)
for integer $m\ge0$ and satisfies the growth conditions required for the
application of Carlson's theorem.

\subsection{Analytic continuation of the replica sum}

Following the Poisson-resummation strategy used in
\cite{calabrese2011entanglement}, the analytic continuation of the replica sum can be
written as
\begin{equation}
S_n(k_0,\Delta)=
\JM{\frac{1}{2}}\sum_{m\ge1}\big[n\,g(nm-\kappa_n)-g(m-\kappa_n)\big]
+
\sum_{m\ge0}\big[n\,g(-nm-\kappa_n)-g(-m-\kappa_n)\big],
\end{equation}
where
\begin{equation}
\kappa_n=-\frac{n k_0}{2\pi}.
\end{equation}

This expression reproduces the original discrete sum for integer $n$ and
provides its analytic continuation.

% \subsection{Derivative at $n=1$}
\subsection{\texorpdfstring{Derivative at $n=1$}{Derivative at n=1}}

To compute the derivative at $n=1$ we differentiate term by term.  In the
strip $0<\Re\Delta<1/2$ the series converge absolutely, allowing the
interchange of differentiation and summation.  One obtains
\begin{equation}
\partial_n S_n\big|_{n=1}=
\sum_{m=1}^\infty\big[g(m-\kappa)+m g'(m-\kappa)\big]
+
\sum_{m=0}^\infty\big[g(m+\kappa)-m g'(m+\kappa)\big],
\end{equation}
\begin{align*}
\partial_n S_n\big|_{n=1}&=\frac{1}{2}
\sum_{m=1}^\infty\big[g(m-\kappa)+m g'(m-\kappa)\big]
+
\sum_{m=0}^\infty\JM{\big[g(-m-\kappa)-m g'(-m-\kappa)\big],}  \nn \\
\JM{\p_n S_n \big|_{n=1}} &= \frac{1}{2}
\sum_{m=1}^\infty\big[g(m-\kappa)+m g'(m-\kappa)\big] + \sum_{m=0}^\infty\JM{\big[g(m-\kappa)+m g'(m-\kappa)\big].}
\end{align*}
where $\kappa=\kappa_{n=1}=-k_0/(2\pi)$.

\subsection{Integral representation}

To evaluate these sums we use the beta-function identity
\begin{equation}
\frac{\Gamma(z+\Delta)}{\Gamma(z+1-\Delta)}
=
\frac{1}{\Gamma(1-2\Delta)}
\int_0^1 dt\, t^{z+\Delta-1}(1-t)^{-2\Delta},
\end{equation}
valid for $0<\Re\Delta<1/2$.

This gives
\begin{equation}
g(z)=A_\Delta\int_0^1 dt\, t^{z+\Delta-1}(1-t)^{-2\Delta},
\qquad
A_\Delta=\frac{2^{2\Delta}\sin(\pi\Delta)}{\pi}.
\end{equation}

Differentiating,
\begin{equation}
g'(z)=A_\Delta\int_0^1 dt\, t^{z+\Delta-1}(1-t)^{-2\Delta}\ln t .
\end{equation}

Substituting these expressions into the derivative formula and
interchanging the order of summation and integration we obtain
\begin{equation}
\partial_n S_n\big|_{n=1}=
A_\Delta\int_0^1 dt (1-t)^{-2\Delta}
\left[
\sum_{m=1}^\infty t^{m-\kappa+\Delta-1}(1+m\ln t)
+
\sum_{m=0}^\infty t^{m+\kappa+\Delta-1}(1-m\ln t)
\right].
\end{equation}
\begin{equation}
\partial_n S_n\big|_{n=1}=
\JM{\frac{A_\Delta}{2}}\int_0^1 dt (1-t)^{-2\Delta}
\left[
\sum_{m=1}^\infty t^{m-\kappa+\Delta-1}(1+m\ln t)
+
\sum_{m=0}^\infty t^{m\JM{-}\kappa+\Delta-1}(1\JM{+}m\ln t)
\right].
\end{equation}

\subsection{Performing the sums}

Using the geometric-series identities
\begin{align}
\sum_{m=1}^\infty t^m &=\frac{t}{1-t},\\
\sum_{m=1}^\infty m t^m &=\frac{t}{(1-t)^2},\\
\sum_{m=0}^\infty t^m &=\frac{1}{1-t},
\end{align}
the sums can be evaluated explicitly.  This yields
\begin{equation}
\partial_n S_n\big|_{n=1}=
A_\Delta\int_0^1 dt
\left[
\frac{t^{\Delta-\kappa}+t^{\Delta+\kappa-1}}{(1-t)^{2\Delta+1}}
+
\frac{(t^{\Delta-\kappa}-t^{\Delta+\kappa})\ln t}{(1-t)^{2\Delta+2}}
\right].
\end{equation}
\begin{equation}
\partial_n S_n\big|_{n=1}=
\JM{\frac{A_\Delta}{2}}\int_0^1 dt
\left[
\frac{t^{\Delta-\kappa}+t^{\Delta\JM{-}\kappa-1}}{(1-t)^{2\Delta+1}}
+
\frac{\JM{2t^{\Delta-\kappa}}\ln t}{(1-t)^{2\Delta+2}}
\right].
\end{equation}

\subsection{Final evaluation}

The remaining integrals can be evaluated using
\begin{align}
\int_0^1 dt\, t^{a-1}(1-t)^{b-1}&=B(a,b),\\
\int_0^1 dt\, t^{a-1}(1-t)^{b-1}\ln t
&=B(a,b)\,[\psi(a)-\psi(a+b)].
\end{align}

After straightforward algebra we obtain
\begin{equation}
\begin{aligned}
\partial_n S_n(k_0,\Delta)\big|_{n=1}
&=
\frac{2^{2\Delta}\sin(\pi\Delta)}{\pi}
\Big[
B(1+\Delta-\kappa,-2\Delta)
+
B(\Delta+\kappa,-2\Delta)
\\
&\quad
+
B(1+\Delta-\kappa,-1-2\Delta)
(\psi(1+\Delta-\kappa)-\psi(-\Delta-\kappa))
\\
&\quad
-
B(1+\Delta+\kappa,-1-2\Delta)
(\psi(1+\Delta+\kappa)-\psi(\kappa-\Delta))
\Big],
\end{aligned}
\end{equation}
\begin{equation}
\begin{aligned}
\partial_n S_n(k_0,\Delta)\big|_{n=1}
&=
\JM{\frac{2^{2\Delta-1}\sin(\pi\Delta)}{\pi}}
\Big[
B(1+\Delta-\kappa,-2\Delta)
+
B(\Delta\JM{-}\kappa,-2\Delta)
\\
&\quad
+ \JM{2}
B(1+\Delta-\kappa,-1-2\Delta)
(\psi(1+\Delta-\kappa)-\psi(-\Delta-\kappa))\Big],
% \\
% &\quad
% -
% B(1+\Delta+\kappa,-1-2\Delta)
% (\psi(1+\Delta+\kappa)-\psi(\kappa-\Delta))
% \Big],
\end{aligned}
\end{equation}
with
\begin{equation}
\kappa=-\frac{k_0}{2\pi}.
\end{equation}

This expression gives the analytically continued derivative of the replica
sum at $n=1$.

Now if we set $\kappa = 0$, we get the desired result. This can be shown by using the identities of Gamma and digamma functions.
\begin{align}
    \mu_0'(1) =  \p_nS_n(k_0 = 0,\Delta)\big|_{n=1}& = \frac{2^{2\Delta -1}\sin (\pi \Delta)}{\pi} \Big[ \frac{\Gamma(\Delta +1) \Gamma(-2\Delta)}{\Gamma(-\Delta +1)} + \frac{\Gamma(\Delta) \Gamma(-2\Delta)}{\Gamma(-\Delta)} \nn \\
    &\qquad \qquad \qquad+ 2B(\Delta + 1, -2\Delta -1) (\psi(1+ \Delta) - \psi(-\Delta))\Big]  
    \nn \\
    & = \frac{2^{2\Delta -1}\sin (\pi \Delta)}{\pi}  \Big[ \frac{\Delta \Gamma(\Delta)\Gamma(-2\Delta)}{-\Delta \Gamma(-\Delta) } + \frac{\Gamma(\Delta)\Gamma(-2\Delta)}{\Gamma(-\Delta)}\nn \\
    &\qquad \qquad \qquad+ \frac{2\Gamma(\Delta + 1)\Gamma(-2\Delta - 1)}{\Gamma(-\Delta)}(-\pi \cot (\pi \Delta))
    \Big]\nn \\
    &= 2^{2\Delta}\cos (\pi \Delta) \Big[  \frac{\Gamma(\Delta + 1)\Gamma(-2\Delta -1)}{\frac{\pi }{\sin(\pi \Delta)\Gamma(1+\Delta)}}\Big]\nn \\
    &=  \frac{2^{2\Delta - 1}} {\pi}\sin(2\pi \Delta) (\Gamma(1+\Delta))^2 \Gamma(-2\Delta - 1) \nn \\
    &= 2^{2\Delta - 1}  \frac{(\Gamma(1+\Delta))^2}{\Gamma(2+2\Delta)} \nn \\
    &= \frac{2^{2\Delta -1}(\Gamma(1+\Delta))^2}{\Big(\frac{\Gamma(1+\Delta)\Gamma(\Delta + 3/2)}{\sqrt{\pi} 2^{1-2(1+\Delta)}}\Big)} \nn \\
    &= \frac{\sqrt{\pi }\Gamma(1+\Delta)}{4 \Gamma(\Delta + 3/2)}.
\end{align}

where we used the following identities,
\begin{align*}
    \psi(1+\Delta) - \psi(-\Delta) &= -\pi \cot(\pi \Delta)\nn \\
    \Gamma(1+z)\Gamma(-z) &= -\frac{\pi}{\sin(\pi z)} \nn \\ 
    \Gamma(z)\Gamma(z+1/2)& = 2^{1-2z}\sqrt{\pi}\Gamma(2z)
\end{align*}

% \JM{The result above should match with $\mu_0'(1)= \frac{2^{-2\Delta_L}\sqrt{\pi}\Gamma(\Delta_L+1)}{4\Gamma(\Delta_L+3/2)}$, however from the result above, $\partial_n S_n(k_0 =0,\Delta)\big|_{n=1} = 0$. This can be shown by taking $\kappa =0$, the third and fourth term cancel, and we are left with
% \begin{align}
%     B(1+\Delta - ,-2\Delta) + B(\Delta + ,-2\Delta) &= \frac{\Gamma(1+\Delta)\Gamma(-2\Delta)}{\Gamma(1-\Delta)} + \frac{\Gamma(\Delta)\Gamma(-2\Delta)}{\Gamma(-\Delta)} \nn \\
%     &= \frac{\Delta \Gamma(\Delta) \Gamma(-2\Delta)}{-\Delta \Gamma(-\Delta)} + \frac{\Gamma(\Delta)\Gamma(-2\Delta)}{\Gamma(-\Delta)} \nn \\
%     &= 0.
% \end{align}
% }

% \section{Analytic continuation of $\mu_1(n)$}
\section{\texorpdfstring{Analytic continuation of $\mu_1(n)$}{Analytic continuation of mu1(n)}}
We want to compute the analytic continuation of the second summation of series in $\mu_1(n)$, which is the second moment (it has $p^2$ in the numerator, instead of $p$ which was the case for $\mu_0$ and the first first summation of $\mu_1$). We derive the analytic continuation of 
\begin{align}
    \mu_1(n) = r_n(x) \sum_{p=1}^{n-1} \frac{p^2 e^{2 \pi i \kappa_n p}}{\sin ^{2\Delta} \lb \frac{\pi p}{n}\rb} .
\end{align}
  The analyic continuation of the first moment was computed, and is given in eq. \textbf{A.8}. We can compute the second moment, by taking a derivative with respect to $\kappa_n$, and finally take derivative with $n$ before taking the limit $n \to 1$. 

  The derivative of eq. A.8 with $\kappa_n$ gives,
  \begin{align}
      \p_{\kappa_n}S_n(k_0,\Delta) = \sum_{m\geq 1} \left[ -n g'(n(m-\kappa)) + g'(m-n\kappa) \right] + \sum_{m\geq 0} \left[ -ng'(-n(m+k))  + g'(-m-nk)\right]. 
  \end{align}
Now we can take the derivative with n, and set $n\to 1$,
\begin{align}
    \p_n\p_{\kappa_n}S_n\big|_{n= 1} = \sum_{m\geq 1}\left[ -g'(m-\kappa) -mg''(m-\kappa) \right] + \sum_{m\geq 0}\left[ -g'(m+\kappa)  + mg''(m+\kappa)\right]
\end{align}
\begin{align*}
    \p_n\p_{\kappa_n}S_n\big|_{n= 1} = \sum_{m\geq 1}\left[ -g'(m-\kappa) -mg''(m-\kappa) \right] + \sum_{m\geq 0}\left[ -g'(\JM{-m-\kappa})  + mg''(\JM{-m-\kappa})\right]
\end{align*}
Now, we can use the integral representation of $g(z),\text{ } g'(z) \text{ and } g''(z)$, where 
\begin{align}
    g''(z) = A_\Delta \int_0^1 dt \, t^{z + \Delta -1} \, (1-t)^{-2\Delta} \, (\ln t)^2
\end{align}
This gives,
\begin{align}
    \p_n\p_{\kappa_n}S_n\big|_{n= 1} = -A_{\Delta} \int_0^1 dt\, \left[\frac{(t^{\Delta - \kappa} + t^{\Delta + \kappa - 1})(\ln t)}{(1-t)^{2\Delta +1}} + \frac{(t^{\Delta - \kappa} - t^{\Delta + \kappa})(\ln t)^2}{(1-t)^{2\Delta +2}}\right] 
\end{align}
\begin{align*}
    \p_n\p_{\kappa_n}S_n\big|_{n= 1} = -A_{\Delta} \int_0^1 dt\, \left[\frac{(t^{\Delta - \kappa} + t^{\Delta \JM{-}\kappa - 1})(\ln t)}{(1-t)^{2\Delta +1}} + \frac{\JM{2t^{\Delta - \kappa} (\ln t)^2}}{(1-t)^{2\Delta +2}}\right] 
\end{align*}
Using the identities from A.19, A.20, and 
\begin{align}
    \int_0^1 dt \, t^{a-1} (1-t)^{b-1}(\ln t)^2 = B(a,b)\left[( \psi (a) - \psi(a+b))^2 + (\psi_1(a) - \psi(a+b))\right]
\end{align}
We obtain,
\begin{align}
    \p_n \p_{\kappa_n}S_n(k_0,\Delta)\big|_{n=1} =& -A_{\Delta}\Big[ B(1+ \Delta - \kappa, -2\Delta)(\psi(1+ \Delta - \kappa) - \psi(1-\Delta - \kappa)) \nn \\
    &+ B(\Delta + \kappa, - 2\Delta)(\psi(\Delta + \kappa) - \psi(\kappa - \Delta))  \nn \\
    &+B(1+ \Delta - \kappa, -2\Delta - 1)\big( \lb \psi(1+\Delta - \kappa) - \psi(-\kappa-\Delta)\rb^2 \nn \\
    & \qquad \qquad \qquad \qquad \qquad +\lb \psi_1(1+\Delta - \kappa) - \psi_1(-\kappa - \Delta)\rb \big)\nn \\
    &- B(1+\Delta + \kappa, -2\Delta - 1) \big( \lb \psi(1+ \Delta + \kappa) - \psi(\kappa - \Delta) \rb^2 \nn \\
    &\qquad \qquad \qquad \qquad \qquad + \lb \psi_1(1+ \Delta + \kappa ) - \psi_1(\kappa - \Delta) \rb 
    \big)
    \Big] 
\end{align}
\begin{align*}
    \p_n \p_{\kappa_n}S_n(k_0,\Delta)\big|_{n=1} =& -A_{\Delta}\Big[ B(1+ \Delta - \kappa, -2\Delta)(\psi(1+ \Delta - \kappa) - \psi(1-\Delta - \kappa)) \nn \\
    &+ B(\Delta \JM{-} \kappa, - 2\Delta)(\psi(\Delta \JM{-} \kappa) - \psi(\JM{-}\kappa - \Delta))  \nn \\
    &+\JM{2}B(1+ \Delta - \kappa, -2\Delta - 1)\big( \lb \psi(1+\Delta - \kappa) - \psi(-\kappa-\Delta)\rb^2 \nn \\
    & \qquad \qquad \qquad \qquad \qquad +\lb \psi_1(1+\Delta - \kappa) - \psi_1(-\kappa - \Delta)\rb \big)\Big]\nn \\
    % &- B(1+\Delta + \kappa, -2\Delta - 1) \big( \lb \psi(1+ \Delta + \kappa) - \psi(\kappa - \Delta) \rb^2 \nn \\
    % &\qquad \qquad \qquad \qquad \qquad + \lb \psi_1(1+ \Delta + \kappa ) - \psi_1(\kappa - \Delta) \rb 
    % \big)
    % \Big] 
\end{align*}
We can set $\kappa = 0 $ in the above result, and again the third and fourth term gets cancelled. We are left with,
\begin{align}
    \p_n \p_{\kappa_n}S_n(k_0 = 0,\Delta)\big|_{n=1} &= A_\Delta B(\Delta , -2\Delta)\big( \psi(1+\Delta) - \psi(1-\Delta)  - \psi(\Delta ) - \psi(-\Delta)\big) 
\end{align}
\JM{Using the following identities,}
\begin{align*}
    \JM{B(\Delta +1, -2\Delta)} &= \JM{-B(\Delta, -2\Delta)} \nn\\
    \JM{ B(\Delta + 1, -2\Delta -1)} &= \JM{-\frac{\Delta}{2\Delta +1} B(\Delta,-2\Delta)} \nn\\
    \JM{\psi(1+\Delta) - \psi(1-\Delta)} &= \JM{-\pi \cot(\pi \Delta)  + \frac{1}{\Delta}} \nn \\
    \JM{\psi(\Delta) - \psi(-\Delta)} &= \JM{-\pi \cot(\pi \Delta) -\frac{1}{\Delta}} \nn \\
    \JM{\psi(1+\Delta) - \psi(-\Delta)} &= \JM{-\pi \cot(\pi \Delta)}\nn \\
    \JM{\psi_1(1+\Delta) + \psi_1(-\Delta) } &= \JM{\pi^2 \csc ^2(\pi\Delta)}
\end{align*}
\JM{we get,}
\begin{align*}
    \p_n \p_{\kappa_n}S_n(k_0 = 0,\Delta)\big|_{n=1} &= A_\Delta \JM{B(\Delta , -2\Delta)}\Big[ \JM{\frac{2}{\Delta} + \frac{2\Delta}{\Delta + 1}\lb \pi^2 \cot^2(\pi \Delta)  - \psi_1(1+\Delta) + \psi(-\Delta))\rb}\Big]
\end{align*}
We use the following identity,
\begin{align}
    \psi(z+1)= \psi(z) + \frac{1}{z}
\end{align}
We obtain,
\begin{align}
    \p_n \p_{\kappa_n}S_n(k_0 =0,\Delta)\big|_{n=1} &= \frac{2A_\Delta B(\Delta,-2\Delta)}{\pi} = \frac{2^{2\Delta + 1} \sin(\pi \Delta) \Gamma(\Delta) \Gamma(-2\Delta)}{\pi \Gamma(-\Delta)}
\end{align}
\begin{align*}
    \p_n \p_{\kappa_n}S_n(k_0 =0,\Delta)\big|_{n=1} &= \JM{\frac{2^{2\Delta}\sin^2(\pi \Delta)}{\pi^2}(\Gamma(\Delta))^2 \Gamma(1-2\Delta) \Big[ \frac{1}{\Delta} + \frac{\Delta}{2\Delta + 1} \lb \pi^2 \lb1+2\cot^2(\pi\Delta)\rb-2\psi_1(1+\Delta)\rb\Big]}
    % \frac{2A_\Delta B(\Delta,-2\Delta)}{\pi} = \frac{2^{2\Delta + 1} \sin(\pi \Delta) \Gamma(\Delta) \Gamma(-2\Delta)}{\pi \Gamma(-\Delta)}
\end{align*}

% \section{Analytic continuation of $\mu_1(n)$ using Cardy, Tonni and collaborators}
\section{\texorpdfstring{Analytic continuation of $\mu_1(n)$ using Cardy, Tonni and collaborators}
{Analytic continuation of mu1(n) using Cardy, Tonni and collaborators}}
Cardy, Tonni and collaborators analyticaaly continued the sum $S_2(n)$ and its derivative at $n=1$ using Carlson's theorem. They analytically continued 
\begin{align*}
    S_2(n) = \frac{1}{2}\sum_{j=1}^{n-1} \frac{1}{\big[\sin^2\lb \frac{2\pi j}{2n} \rb\big]^\alpha} \equiv \frac{1}{2}\sum_{k = -\infty}^{\infty}[ng(nk)-g(k)]
\end{align*}
where $g(k)$ is the analytic continuation of the Fourier modes $f(k)$ of $\lb\sin^2(u/2)\rb^{-\alpha}$. However, for $\mu_1(n)$, we want the analytic continuation of the sum involving,
\begin{align}
    \mu_1(n) = \frac{1}{2}\sum_{p=1}^{n-1}\frac{p^2}{\lb \sin^2\lb \frac{\pi p}{n}\rb \rb^\Delta}
\end{align}
For this higher moment, we will consider the Fourier modes of the generating function,
\begin{align}
    \frac{e^{-2i\kappa u }}{(\sin^2(u))^\Delta}
\end{align}
We can take two derivatives with respect to $\kappa$ to obtain the desired sum, and set $\kappa= 0$ in the end. 
By adding this extra source term, we replace $m$ with $m + \kappa$ in eq.~[\ref{eq: Fourier coeff of csc^2Delta}], and obtain,
\begin{align}
    g(m,\kappa) = c_\Delta \frac{\Gamma(1-2\Delta) \Gamma(m + \kappa + \Delta)}{\Gamma(m+\kappa + 1 - \Delta)}
\end{align}
Then the sum is given by,
\begin{align}
    \mu_1(n) &= -\frac{1}{4}\frac{1}{2}\p_\kappa^2 \left[ \sum_{m\geq 1} \lb ng(n(m+\kappa)) - g(m+\kappa) \rb+ \sum_{m\geq 0}\lb ng(-nm+n\kappa) - g(-m-\kappa) \rb \right] \nn \\
    & = -\frac{1}{8} \left[ \sum_{m\geq 1}\lb n^3 g''(n(m+\kappa)) - g''(m+\kappa) \rb +\sum_{\geq 0}\lb n^3g''(-nm+n\kappa) -g''(-m-\kappa) \rb  \right].
\end{align}
% where in the second step, we used $g(-m) = g(m)$ for integer $m$.
Then, we take the derivative with respect to $n$, and take the limit $n=1$ and $\kappa = 0$, along with using the reflection $g(-m) = g(m)$ for integer $m$, to obtain,
\begin{align}
    \p_n \mu_1 \big|_{n=1,\kappa = 0} = -\frac{1}{8} \left[ \sum_{m\geq 1} \lb 3g''(m) + mg'''(m) \rb  + \sum_{m\geq 0}\lb  3g''(m) + mg'''(m) \rb\right]
\end{align}
Making use of the following integral representation of the analytically continued Fourier modes, 
\begin{align}
    g(m+\kappa) &= A_\Delta \int_0^1 dt \, t^{m+\kappa + \Delta - 1} (1-t)^{-2\Delta} \\
    g'(m+\kappa) &= A_\Delta \int_0^1 dt \, t^{m+\kappa + \Delta - 1} (1-t)^{-2\Delta} \ln t \\ 
    g''(m+\kappa) &= A_\Delta \int_0^1 dt \, t^{m+\kappa + \Delta - 1} (1-t)^{-2\Delta} (\ln t)^2 \\ 
    g'''(m+\kappa) &= A_\Delta \int_0^1 dt \, t^{m+\kappa + \Delta - 1} (1-t)^{-2\Delta} (\ln t)^3 \\ 
\end{align}
We obtain,
\begin{align}
    \p_n \mu_1(n)\big|_{n=1} &= -\frac{A_\Delta}{8} \Big( \sum_{m\geq 1} 3 \int_0^1 dt \, t^{m+\kappa + \Delta - 1}(1-t)^{-2\Delta} (\ln t)^2  + \sum_{m\geq 1} m\int_0^1 dt \, t^{m+\kappa + \Delta - 1}(1-t)^{-2\Delta} (\ln t)^3 \\
    & \qquad \qquad  \sum_{m\geq 0 }3 \int_0^1 dt \, t^{m+\kappa + \Delta - 1}(1-t)^{-2\Delta} (\ln t)^2  +  \sum_{m\geq 1} m\int_0^1 dt \, t^{m+\kappa + \Delta - 1}(1-t)^{-2\Delta} (\ln t)^3
    \Big)
\end{align}
Doing the sum over $m$, we get,
\begin{align}
    \p_n \mu_1(n)\big|_{n=1} &= -\frac{A_\Delta}{8} \Bigg( 3\int_0^1 dt \, \frac{t^{\kappa + \Delta} (\ln t)^2}{(1-t)^{2\Delta + 1}} + 2\int_0^1  dt \, \frac{t^{\kappa + \Delta }(\ln t)^3}{(1-t)^{2\Delta + 2}} +  3\int_0^1 dt \, \frac{t^{\kappa + \Delta-1} (\ln t)^2}{(1-t)^{2\Delta + 1}} \Bigg)
\end{align}
We use the following integrals,
\begin{align}
    \int_0^1 dt \, t^{a-1} (1-t)^{b-1}(\ln t)^2 &= B(a,b) \left[ \lb \psi(a) - \psi(a+b) \rb^2 + \lb \psi_1(a) - \psi_1(a+b) \rb  \right] \\ 
    \int_0^1 dt \, t^{a-1} (1-t)^{b-1}(\ln t)^3 &= B(a,b)\lb \psi(a) - \psi(a+b) \rb \left[\lb \psi(a) - \psi(a+b)  \rb^2 + \lb \psi_1 (a) - \psi_1(a+b) \rb \right] \nn \\ 
    & +B(a,b) \left[ 2 \lb \psi(a) - \psi(a+b)  \rb \lb \psi_1(a) - \psi_1(a+b)  \rb + \lb \psi_2(a) - \psi_2(a+b)\rb\right]
\end{align}
to obtain,
\begin{align}
    \p_n \mu_1(n)\big|_{n=1,\kappa = 0} &= -\frac{A_\Delta}{8} \Big[ 3B(1+\Delta, -2\Delta) \left( \lb \psi(\Delta+1) - \psi(-\Delta+1)\rb^2  + \lb \psi_1(\Delta+1) - \psi_1(-\Delta+1) \rb  \right) \Big] \nn \\
    &\qquad - \frac{A_\Delta}{8} \Big[ 3B(\Delta, -2\Delta) \lb  \lb \psi(\Delta) - \psi(-\Delta) \rb^2  + \lb \psi_1(\Delta) - \psi_1(-\Delta) \rb \rb \Big] \nn \\
    & \qquad -\frac{A_\Delta}{8}\Bigg[ 2B(\Delta + 1,-2\Delta -1) \Big[ \lb \psi(\Delta + 1) - \psi(-\Delta)\rb \times \nn \\ 
    &\qquad \qquad \lb \lb \psi(\Delta + 1) - \psi(-\Delta) \rb^2 + \lb \psi_1(\Delta + 1) - \psi_1(-\Delta) \rb   \rb   \nn \\
    &\qquad \qquad +2 \lb \psi(\Delta + 1) - \psi(-\Delta)\rb \lb \psi_1(\Delta + 1) - \psi_1(-\Delta) \rb + \lb \psi_2(\Delta + 1) - \psi_2(-\Delta)\rb \Big] \Bigg]
\end{align}
We use the following identities of Beta function, digamma, trigamma and poly-gamma function, 
\begin{align}
    B(1+\Delta, -2\Delta) &= -B(\Delta, -2\Delta)\\ 
    B(\Delta +1,-2\Delta -1)&= -\frac{\Delta}{2\Delta +1} B(\Delta, -2\Delta)\\
    \psi(\Delta) - \psi(-\Delta) &= -\pi \cot(\pi \Delta) - \frac{1}{\Delta} \\ 
    \psi(1+\Delta) - \psi(1-\Delta) &= -\pi \cot(\pi \Delta) + \frac{1}{\Delta} \\ 
    \psi(1+\Delta) - \psi(-\Delta) &= -\pi \cot(\pi \Delta) \\ 
    \psi_1(\Delta) - \psi_1(-\Delta) &= 2\psi_1(\Delta) - \pi^2 \csc^2(\pi \Delta) -\frac{1}{\Delta^2} \\
    \psi_1(1+\Delta) - \psi_1(1-\Delta) &= 2\psi_1(\Delta) - \pi^2 \csc^2 (\pi \Delta) - \frac{1}{\Delta^2} \\
    \psi_1(1+\Delta) - \psi_1(-\Delta) &= 2\psi_1(\Delta) -\pi^2 \csc^2(\pi\Delta) -\frac{2}{\Delta^2}\\
    \psi_2(\Delta) - \psi(-\Delta) &= -2\pi^3 \csc^2 (\pi \Delta \cot(\pi \Delta)) -\frac{2}{\Delta^3} \\
    \psi_2(1+\Delta) - \psi_2(1-\Delta) &= -2\pi^3 \csc^2 (\pi \Delta) \cot(\pi \Delta) + \frac{2}{\Delta^3} \\ 
    \psi_2(1+\Delta)- \psi_2(-\Delta) &= -2\pi^3 \csc^2 (\pi \Delta)\cot(\pi \Delta)
\end{align} 
Plugging these in the previous equation to obtain,
\begin{align}
    \p_n \mu_1 \big|_{n=1} &= -\frac{A_\Delta}{8} B(\Delta, -2\Delta)\Bigg( 3  \Big[ \lb -\pi \cot(\pi \Delta) - \frac{1}{\Delta} \rb^2 - \lb -\pi \cot(\pi \Delta) + \frac{1}{\Delta} \rb^2\Big] \Bigg) \nn \\ 
    &\qquad +\frac{A_\Delta}{4}\frac{\Delta}{2\Delta+1} B(\Delta, -2\Delta) \Big(  C(C^2 + D) + 2CD + F \Big) \\
    & = T_1 + T_2 \\
\end{align}
where 
\begin{align}
    P &= \pi \cot(\pi \Delta)\\
    Q&= \pi^2 \csc^2(\pi \Delta)\\
    C &= \psi(1+\Delta) - \psi(-\Delta) = -P \\
    D & = \psi_1(1+\Delta) - \psi_1(-\Delta) = 2\psi_1(\Delta ) - Q -\frac{2}{\Delta^2} \\ 
    F &= \psi_2(1+\Delta) - \psi_2(-\Delta) = -2PQ
\end{align}
Then the second term ($T_2$) in these new variables is given by,
\begin{align}
    T_2 &= -P(P^2 + D) -2PD - 2PQ \nn \\
    &= -P^3 -3PD -2PQ \nn \\
    &= -P^3  -3P(2\psi_1(\Delta) - Q - \frac{2}{\Delta^2}) - 2PQ\nn \\
    &= -P^3 -6P\psi_1(\Delta) + PQ + \frac{6P}{\Delta^2} \nn \\
    &= P\left[ \pi^2 -6\psi_1(\Delta) + \frac{6}{\Delta^2} \right]
\end{align}
where in the last step we used $Q^2 = \pi^2 + P^2$. 
Finally we get,
\begin{align}
    \p_n \mu_1 \big|_{n=1, \kappa = 0} &= -\frac{3A_\Delta B(\Delta,-2\Delta)\pi \cot(\pi\Delta)}{2\Delta} + \frac{A_\Delta B(\Delta,-2\Delta) \Delta \pi \cot(\pi \Delta)}{4(2\Delta +1)} \lb \pi^2 -6\psi_1(\Delta) + \frac{6}{\Delta^2} \rb \nn \\
    &= \frac{A_\Delta B(\Delta,-2\Delta)\pi \cot(\pi\Delta)}{8}\left[ -\frac{12}{\Delta} + \frac{2\Delta }{2\Delta + 1} \lb \pi^2 -6\psi_1(\Delta) + \frac{6}{\Delta^2} \rb  \right]
\end{align}
\subsection{\texorpdfstring{$\Delta = 1$ case} {Delta = 1 case}}
We can also check that it does not blow at $\Delta = 1$, which is of interest as it corresponds to massless free boson theory. 
We use the following identity,
\begin{align}
    B(\Delta, -2\Delta)\Big|_{\Delta \to N} = (-1)^N \frac{(N-1)! (N)!}{2(2N)!}.
\end{align}
to get $B(1,-2) = -1/4$. Hence,
\begin{align}
    \p_n \mu_1 (\Delta = 1) \big|_{n=1, \kappa = 0} &= \frac{1}{8} \left[ -12 + \frac{2}{3}\lb \pi^2 -6\psi_1(1) + 6 \rb \right] 
\end{align}
We use the following series expansion of poly-gamma function,
\begin{align}
    \psi_m(z) = (-1)^{m+1} (m!) \sum_{k=0}^\infty \frac{1}{(z+k)^{m+1}}, \qquad \qquad m\geq 1
\end{align}
to get, 
\begin{align}
    \p_n \mu_1  \big|_{n=1, \kappa = 0} (\Delta = 1)&= -1
\end{align}
Similarly, we can check $\p_n \mu_1 \big|_{n=1, \kappa = 0}(\Delta = 2)  = 0$. 
\section{\texorpdfstring{$S_2'(2)$}{ S2'(2)}}
We want to compute the derivative of $S_2(n)$ at $n=2$. Thus we require,
\begin{align}
    \p_n S_n |_{n=2} &= \frac{1}{2}\sum_{m\geq 1} \lb g(2m) + 2m g'(2m) \rb + \frac{1}{2}\sum_{m\geq 0} \lb g(-2m) -2m g'(-2m) \rb \\
    &= \frac{1}{2}\sum_{m\geq 1} \lb g(2m) + 2m g'(2m) \rb + \frac{1}{2}\sum_{m\geq 0} \lb g(2m) +2m g'(2m) \rb
\end{align}
where in the last line, we used the reflection property ($g(-m) = g(m)$) at integer $m$. Then using the integral representation of $g_{\Delta}(m)$ and $g'_\Delta (m)$, we get,
\begin{align}
    \p_n S_n |_{n=2} = \frac{A_\Delta}{2}\left[ \int_0^1 dt \, (1-t)^{-2\Delta} \lb \sum_{m\geq 1}t^{2m+\Delta - 1}(1+2m \ln t) + \sum_{m\geq 0} t^{2m+\Delta - 1}(1+2m \ln t)\rb\right]
\end{align}
We use the following geometric sums,
\begin{align}
    \sum_{m\geq 1} t^{2m} &= \frac{t^2}{1-t^2},\\
    \sum_{m\geq 1} 2mt^{2m}&= \frac{2t}{(1-t^2)^2}
\end{align}
to obtain,
\begin{align}
    \p_n S_n |_{n=2} &= \frac{A_\Delta}{2}\Big( \int_0^1 dt \, t^{\Delta +1}\, (1-t)^{-2\Delta - 1}\, (1+t)^{-1} + \int_0^1 dt \, t^{\Delta -1}\, (1-t)^{-2\Delta - 1}\, (1+t)^{-1} \nn \\
    & \qquad +2\int_0^1 dt \, t^{\Delta }\, (1-t)^{-2\Delta - 1}\, (1+t)^{-1}\ln t \Big)
\end{align}
We use the following identities,
\begin{align}
    \int_0^1 dt \, t^{a-1} (1-t)^{b-1} (1+t)^{-c} &= B(a,b)\, _2F_1(a,c;a+b;-1)\\
    \int_0^1 dt \, t^{a-1} (1-t)^{b-1} (1+t)^{-c} \ln t &= \p_z \lb B(z,b)\, _2F_1(z,c;z+b;-1) \rb|_{z=a}
\end{align}
Denote 
\begin{align}
    I(a,b,c) = B(a,b) \, _2F_1(a,c;a+b;-1) \\
\end{align}
we finally obtain,
\begin{align}
    \p_nS_2(n)|_{n=2} = \frac{A_\Delta}{2}\lb I(\Delta + 2,-2\Delta,1) + I(\Delta,-2\Delta,1) +2\p_aI(a,-2\Delta,1)|_{a = \Delta +1}\rb.
\end{align}

\section{\texorpdfstring{Is there an OPE expansion for $S(\rho_0\|\rho_\Phi)?$}{Is there an OPE expansion for S(rho0 || rho Phi}}

where the first factor comes the normalization of the density matrices, and we have used $\sin(\pi x)=\sin(\pi(1-x))$. We also used the following OPE for the $\Phi$ and $\Phi^*$ with identity close to it,
\begin{align}
    \Phi(z_{1,n}^-) I(z_{0,n}^+) &\sim \Phi(z_{0,n}^+) + \cdots 
    % e^{i\pi (1/n-1/2)}(2\zeta_n)\p \Phi(z_{0,n}^+)+\cdots
    \nn\\
    \Phi^*(z_{n-1,n}^+)I(z_{0,n}^-) &\sim \Phi^*(z_{0,n}^-)+ \cdots 
    % +e^{-i\pi (1/n-1/2)}(2\zeta_n)\p \Phi^*(z_{0,n}^-)\cdots
\end{align}
We compute the two and three point functions above \
% JM{(one should then also consider the two point function of the descendant with the identities, but those are higher order)}
\begin{align}
 % & \braket{V_\beta(z^+_{n-1,n})V_{-\beta}(z^-_{1,n})}=\lb 2\sin\lb\frac{\pi(1+\ep)}{n}\rb\rb^{-4h_\beta}\nn\\
 &\braket{\Phi^*(z^+_{0,n})\Phi(z^-_{0,n})}=\lb 2\sin\lb\frac{\pi(1-\ep)}{n}\rb\rb^{-2\Delta_\Phi}e^{-3\pi i s_\Phi}\\
  &\braket{\Phi(z^+_{n-1,n})\Phi^*(z^-_{1,n})}=\lb 2\sin\lb\frac{\pi(1+\ep)}{n}\rb\rb^{-2\Delta_\Phi}e^{-3\pi i s_\Phi}
% &\braket{V_\beta(z^+_{n-1,n})V_{-\beta}(z^-_{1,n})\mO_L(e^{2\pi i k/n})}=\nn\\
% &C_\beta\lb 2\sin\lb\frac{\pi(1+\ep)}{n}\rb\rb^{-4h_\beta+\Delta_L}\lb 2\sin\lb\frac{\pi(2k+1+\ep)}{2n}\rb\rb^{-\Delta_L}\lb 2\sin\lb\frac{\pi(2k-1-\ep)}{2n}\rb\rb^{-\Delta_L}\nn\\
% &=C_\beta\lb 2\sin\lb\frac{\pi(1+\ep)}{n}\rb\rb^{-4h_\beta+\Delta_L}\lb2\cos(\pi(1+\ep)/n)-2\cos(2\pi k/n) \rb^{-\Delta_L}
\end{align}
\begin{align}
\braket{\Phi(z^-_{0,n})\Phi^*(z^+_{0,n})\mO_L(e^{i\pi(2k+1) /n})}&=C_\Phi\lb 2\sin\lb\frac{\pi(1-\ep)}{n}\rb\rb^{-2\Delta_\Phi+\Delta_L}\lb 2\sin\lb\frac{\pi(k+1-\ep/2)}{n}\rb\rb^{-\Delta_L}\times\nn\\
&\qquad \qquad \qquad \qquad\lb 2\sin\lb\frac{\pi(k+\ep/2)}{n}\rb\rb^{-\Delta_L}e^{-3\pi i s_\Phi}e^{-i\pi s_L\lb (2k+1)/n+1/2)\rb}\nn\\
&=C_\Phi\lb 2\sin\lb\frac{\pi(1-\ep)}{n}\rb\rb^{-2\Delta_\Phi+\Delta_L}\times\nn \\
&\qquad\lb2\cos(\pi(1-\ep)/n)-2\cos(\pi (2k+1)/n) \rb^{-\Delta_L}e^{-3\pi i s_\Phi}e^{-i\pi s_L\lb (2k+1)/n+1/2)\rb}
\end{align}

\begin{align}
\braket{\Phi(z^+_{n-1,n})\Phi^*(z^-_{1,n})\mO_L(e^{i\pi(2k+1) /n})}&=e^{3\pi i s_\Phi/2}e^{-i\pi (s_L+s_\Phi)\lb (2k+1)/(2n)+1/2)\rb} \times \nn\\
&C_\Phi^2\lb 2\sin\lb\frac{\pi(1+\ep)}{n}\rb\rb^{\Delta_L}\times\lb 2\sin\lb\frac{\pi(2k+ 1))}{2n}\rb\rb^{-2\Delta_L}\nn\\
% &=C_\Phi\lb 2\sin\lb\frac{\pi(1-\ep)}{n}\rb\rb^{-2\Delta_\Phi+\Delta_L}\times\nn \\
% &\qquad\lb2\cos(\pi(1-\ep)/n)-2\cos(\pi (2k+1)/n) \rb^{-\Delta_L}e^{-3\pi i s_\Phi}e^{-i\pi s_L\lb (2k+1)/n+1/2)\rb}
\end{align}


Therefore, in the small $\ep$ expansion we have
% \begin{align}
% \braket{\tilde{\mO}_{0|\beta}^{(n-1)}}&=D_n(\ep)^{n-1}\lb\frac{\sin(\pi\ep/n)}{\sin(\pi(1+\ep)/n)}\rb^{4h_\beta}\times\nn\\
% &\lb 1+C_\beta^2\lb\frac{\sin(\pi\epsilon/n)}{\sin(\pi(1+\epsilon)/n)}\rb^{\Delta_L}\sum_{k=0}^{n-2} \lb 2\cos(\pi(1+\ep)/n)-2\cos(2\pi k/n) \rb^{-\Delta_L} +\cdots\rb\ .
% \end{align}
\begin{align}
\braket{\tilde{\mO}_{0|\Phi}^{(n-1)}}=&\lb \frac{n \sin(\pi\ep/n)}{\sin(\pi\ep)}\rb^{-2\Delta_\Phi (n-1)}\lb\frac{\sin(\pi\ep/n)}{\sin(\pi(1-\ep)/n))}\rb^{2\Delta_\Phi}
% \lb\frac{1}{2\sin(\pi(1+\ep)/n)}\rb^{4h_\beta}
\times\nn\\
&\Bigg( 1+C_\Phi\lb\sin(\pi\epsilon/n)\sin(\pi(1-\epsilon)/n)\rb^{\Delta_L}\times \nn \\
&\sum_{k=1}^{n-2}\lb \frac{1}{\sin(\pi (k+1-\ep/2)/n) \sin (\pi (k+\ep/2) /n)}\rb^{-\Delta_L} +\cdots\Bigg)\ .
\end{align}
\begin{align}
\braket{\tilde{\mO}_{0|\Phi}^{(n-1)}}=&\lb \frac{n \sin(\pi\ep/n)}{\sin(\pi\ep)}\rb^{-2\Delta_\Phi (n-1)}\lb\frac{\sin(\pi\ep/n)}{\sin(\pi(1+\ep)/n))}\rb^{2\Delta_\Phi}
% \lb\frac{1}{2\sin(\pi(1+\ep)/n)}\rb^{4h_\beta}
\times\nn\\
&\Bigg( 1+C_\Phi^2 e^{4\pi i s_\Phi} \frac{\lb \sin \lb \pi (1+\ep)/n\rb \rb^{2\Delta_\Phi + \Delta_L}}{\lb \sin (\pi \ep / n)\rb^{2\Delta_\Phi - \Delta_L}}\times \sum_{k=1}^{n-2} \frac{e^{i \pi (s_L - s_\Phi)((2k+1)/(2n))}}{\lb \sin\lb \frac{ \pi(2k+1)}{2n}\rb\rb^{2\Delta_L } } +\cdots\Bigg) \\
=& \lb \frac{n \sin(\pi\ep/n)}{\sin(\pi\ep)}\rb^{-2\Delta_\Phi (n-1)}\lb\frac{\sin(\pi\ep/n)}{\sin(\pi(1+\ep)/n))}\rb^{2\Delta_\Phi}
% \lb\frac{1}{2\sin(\pi(1+\ep)/n)}\rb^{4h_\beta}
\times\nn\\
&\Bigg( 1+C_\Phi^2 e^{4\pi i s_\Phi} \lb\frac{ \sin \lb \pi (1+\ep)/n\rb }{ \sin (\pi \ep / n)}\rb^{2\Delta_\Phi}(\sin(\pi(1+\ep)/n)\sin (\pi \ep/n))^{ \Delta_L} \sum_{k=1}^{n-2} \frac{1}{\lb \sin\lb \frac{ \pi(2k+1)}{2n}\rb\rb^{2\Delta_L} } +\cdots\Bigg)
\end{align}
If we did not have any phase term in the sum in the last line of the above equation, then we have
\begin{align}
    \sum_{k=0}^{n-1} \frac{1}{\sin^2 \lb \pi (2k+1) / (2n)\rb^{\Delta}}= \frac{2}{2n}\sum_{k=1}^{2n-1} \frac{1}{\lb \sin^2\lb \frac{\pi j }{2n}\rb \rb^{\Delta}} -\frac{2}{n}S_2(n,\Delta)
\end{align}
Here we wrote the sum over the odd sites to be the sum over all $2n-1$ sites, and subtract the sum over even sites, which is $2S_2(n)/n$. Note that we have to subtract the terms for $k=1,n-1$ as well. 
Thus we can simplify the sum,
\begin{align}
    \sum_{k=1}^{n-2} \frac{1}{\lb \sin\lb \frac{ \pi(2k+1)}{2n}\rb\rb^{2\Delta_L } } = 2^{2\Delta_L}\lb \frac{S_2(2n,\Delta_L)}{n} - \frac{2S_2(n,\Delta_L)}{n} - \frac{2}{\lb \sin(\pi/(2n))\rb^{2\Delta_L}}\rb 
\end{align} 
The correction to the divergent piece comes when we take the derivative of $\log$ of the things inside parentheses. 
We have the divergent piece in the denominator as the common factor. When the derivative hits the summation, we get,
\begin{align}
    S_2'(2) -S_2(2) - 2S_2'(1) 
\end{align}
While the derivative hitting the term outside the sum gives,
\begin{align}
  &  \lb \lb\frac{ \sin \lb \pi (1+\ep)/n\rb }{ \sin (\pi \ep / n)}\rb^{2\Delta_\Phi}(\sin (\pi (1+\ep)/n)\sin (\pi \ep/n))^{ \Delta_L}\rb'\Bigg|_{n=1} \\&= (-1)^{2\Delta_\Phi+\Delta_L}(\pi \ep)^{2\Delta_L} \lb (2\Delta_\Phi + \Delta_L) \lb -\frac{1}{\ep} + \frac{\pi^2 \ep}{3} + \cdots \rb  - 2\Delta_L\lb 1-\frac{\pi^2 \ep^2}{3} \cdots\rb\rb
    % &\simeq  (-1)^{3\Delta_\Phi+1}\lb 3\Delta_\Phi + \ep (\Delta_L + \Delta_\Phi)\rb \lb \pi \ep\rb^{\Delta_L + \Delta_\Phi} \lb \frac{1}{\ep} - \frac{\pi^2 \ep}{3} + \cdots \rb 
\end{align}
So this piece is divergent, but is suppressed from the leading divergent piece. Similarly, the other piece gives,
\begin{align}
     S_2'(2,\Delta_L) -1 -2\frac{\sqrt{\pi } \Gamma(\Delta_L + 1)}{4\Gamma (\Delta_L + 3/2)} 
\end{align}
where we used $2^{2\Delta_L}S_2(2) = 1$
We compute $S_2'(2)$ in the appendix,
\begin{align}
    &S_2'(2) = \frac{2^{2\Delta}\sin(\pi \Delta)}{\pi} \lb I(\Delta +2,-2\Delta,1) + I(\Delta ,-2\Delta,1) +2\p_aI(a,-2\Delta,1)|_{a=\Delta} \rb\\
    & I (a,b,c) = \int_0^1 dt \, t^{a-1}(1-t)^{b-1} (1+t)^{-c} = B(a,b) \,_2F_1(a,c;a+b;-1) \\
    & \p_a I (a,b,c) = \int_0^1 dt \, t^{a-1}(1-t)^{b-1} (1+t)^{-c} \ln t = \p_a \lb B(a,b) \,_2F_1(a,c;a+b;-1)\rb 
\end{align}
As long as $S_2'(2)$ is finite, these corrections are suppressed.
Finally, we obtain the relative entropy,
\begin{align}
    S(\rho_0\|\rho_\Phi) &= - 2\Delta_\Phi \lb \frac{1}{\ep} - \frac{\pi^2\ep}{3} +\cdots \rb \nn \\&+ (-1)^{2\Delta_\Phi + \Delta_L}C_\Phi^2 (\pi \ep)^{2\Delta_L } \Bigg( (2\Delta_\Phi + \Delta_L  ) \lb \frac{1}{\ep} - \frac{\pi^2 \ep}{3}+\cdots\rb+2\Delta_L \lb 1-\frac{\pi^2 \ep^2}{3}\rb \nn \\
    &- \lb  S_2'(2) +1 - 2\frac{\sqrt \pi \Gamma (\Delta +1)}{4\Gamma (\Delta + 3/2)} \rb\Bigg)
\end{align}

We need to take the logarithm followed by a $\partial_n$ at $n\to 1$. It is straightforward to see
\begin{align}
&  \lb \p_n\log \lb D_n(\ep)^{n-1}\rb\rb_{n\to 1}=0\nn\\
% & \JM{\lb \p_n\log \lb \frac{n \sin(\pi\ep/n)}{\sin(\pi\ep)}\rb^{-2\Delta_\Phi (n-1)}\rb_{n \to 1}}=0 \nn\\
&\lb \p_n\log \lb \frac{\sin(\pi\epsilon/n)}{\sin(\pi(1+\epsilon)/n)}\rb\rb_{n\to 1}=\pi\cot(\pi\ep)=\frac{1}{\ep}-\frac{\pi^2 \ep}{3}+O(\ep^3)\nn\\
&\frac{\sin(\pi\epsilon/n)}{\sin(\pi(1+\epsilon)/n)}=\lb\frac{\pi \ep}{n\sin(\pi/n)} \rb
\lb 1-\frac{\cot(\pi/n)\pi \ep}{n}+O(\ep^2)\rb
\end{align}
Therefore,
\begin{align}
  &S(\rho_0\|\rho_\beta)=4h_\beta\lb\frac{1}{\ep}-\frac{\pi^2\ep}{3}\rb- C_\beta^2 \p_n\mathcal{T}(n)|_{n\to 1}+O(\ep^3)\nn\\
  &\mathcal{T}(n)=\lb\frac{\sin(\pi\epsilon/n)}{\sin(\pi(1+\epsilon)/n)}\rb^{\Delta_L}\sum_{k=0}^{n-2} \lb 2\cos(\pi(1+\ep)/n)-2\cos(2\pi k/n) \rb^{-\Delta_L}\nn\\
  &=\lb\frac{\pi \ep}{4n}\rb^{\Delta_L}T(n)\lb 1-O(\ep)\rb\nn\\
  &T(n)=\sum_{k=0}^{n-2}\lb \sin\lb\frac{\pi}{n}\rb\sin\lb\frac{(2k+1)}{2n}\rb\sin\lb \frac{(2k-1)}{2n}\rb\rb^{-\Delta_L}
\end{align}
\begin{align*}
  &S(\rho_0\|\rho_\Phi)=\JM{2\Delta_\Phi \lb\frac{1}{\ep}- \sum_{k=1}^{\infty} \frac{2^{2k}|B_{2k}|} {(2k)!}\ep^{2k-1} 
  % \frac{\pi^2\ep}{3}
  \rb - C_\Phi \p_n\mathcal{T}(n)|_{n\to 1}}\nn\\
  &\mathcal{T}(n)=\JM{\lb\sin(\pi\epsilon/n)\sin(\pi(1-\epsilon)/n)\rb}^{\Delta_L}\JM{\sum_{k=1}^{n-2} \lb \frac{1}{\sin(\pi (k+1-\ep/2)/n) \sin (\pi (k+\ep/2) /n)}\rb^{-\Delta_L}}\nn\\
  &=\JM{\lb\sin(\pi\epsilon/n)\sin(\pi(1-\epsilon)/n)\rb}^{\Delta_L}T(n)\lb 1-O(\ep)\rb\nn\\
  &T(n)=\JM{\sum_{k=1}^{n-2}\lb \frac{1}{\sin(\pi (k+1-\ep/2)/n) \sin (\pi (k+\ep/2) /n)}\rb^{-\Delta_L}}
  % {\sin\lb\frac{\pi}{n}\rb}
  % \rb^{-\Delta_L}}
\end{align*}
Therefore,
\begin{align}
  S(\rho_0\|\rho_\beta)&=4h_\beta\lb\frac{1}{\ep}-\frac{\pi^2\ep}{3}\rb - C_\beta^2 (\pi \ep)^{\Delta_L} T'(1)+\cdots
\end{align}
\begin{align*}
  S(\rho_0\|\rho_\beta)&=\JM{2\Delta_{\Phi}\lb\frac{1}{\ep}
  - \sum_{k=1}^{K} \frac{2^{2k}|B_{2k}|} {(2k)!}\pi^{2k}\ep^{2k-1}
  %\frac{\pi^2\ep}{3}
  \rb - C_\Phi  \mathcal{T}'(1)}+\cdots
\end{align*}
where we have used the fact that at $n=1$ there should be no contribution coming from OPE expansions, so $T(1)=0$.

\JM{Again, with the use of the OPE in eq.~[\ref{eq: general OPE with identity}], we obtain}
\begin{align}
 S(\rho_0\|\rho_\beta)&=\JM{(1-\delta_{\Psi \Phi})\lb 2\Delta_{\Phi}\lb\frac{1}{\ep}
  - \frac{\ep}{3}
  \rb - C_\Phi  \mathcal{T}'(1)+\cdots\rb}
\end{align}

This makes it clear that the divergence of $S(\rho_0\|\rho_\beta)$ comes from the $1/(1-x)$ term in the first term with residue that is universally fixed to be $4h_\beta$. The comparison with the relative entropy $S(\rho_0\|\rho_\beta)$ of the coherent states suggests that the term $T(n)$ does not contribute to the relative entropy, i.e. its analytic continuation vanishes at this order.

\section{Renyi entropy of excited states of Vertex operators}
In \cite{Alcaraz_2011}, they define 
\begin{align}
    F_{\Psi}^n(x) = \frac{\text{tr}(\rho_\Psi^n)}{\text{tr}(\rho_0^n)},
\end{align}
where $\Psi$ is some excited state. 
In the $x \to 0 $ limit, they show for any excited state generated by the operator $\Psi$,
\begin{align}
    F_\Psi^{(n)}(x) \sim 1+ \frac{h+\bar{h}}{3} \lb\frac{1}{n}-n\rb (\pi x)^2 + \mO(x^{2\Delta_{L}}).
\end{align}
assuming that $\Delta_L>1$.
This piece comes from the $D_n(x)$, without considering the contribution from any two point function of primaries from OPEs. This can be shown by considering,
\begin{align}
    F^{(n)}_\Psi(x) &= \frac{\text{tr}(\rho_\Psi^n)}{\text{tr}(\rho_0^n)} \\
    &= D_n(x)^{-n} \\
    &= \lb\frac{\sin(\pi x)}{n\sin(\pi x/n)}\rb^{2n\Delta_\Psi} \\
    &\simeq \lb\frac{\pi x -\frac{(\pi x)^3}{6} + \cdots } {n \lb \frac{\pi x }{n} - \frac{(\pi x)^3}{6n^3} + \cdots\rb  }\rb^{2n\Delta_\Psi}\\
    &= \lb \frac{1 - \frac{(\pi x)^2}{6}}{1- \frac{(\pi x)^2}{6n^2}}  \rb^{2n\Delta_\Psi} \\
    &\simeq 1+ 2n\Delta_\Psi \frac{(\pi x)^2}{6}\lb\frac{1}{n^2} - 1\rb + \cdots \\
    &= 1+ \frac{h + \bar{h}}{3} \lb\frac{1}{n}-n\rb (\pi x)^2 + \cdots
\end{align}
When $\Delta_L = 1$, we have to take the higher order terms into account which will contribute as it is $\mO(x^2)$. This is the case with massless free bosons. 
In particular, using the 2n-point function of the coherent states, they show that 
\begin{align}
    F^{(n)}_{\Psi[r,s]}(x) = 1, \qquad \forall n,r,s. 
\end{align} 
where $\Psi[r,s] = e^{i (\alpha_+ \phi + \alpha_- \phi)}$
is the vertex operator. This shows that the Renyi entropy of any order $n$ for coherent states is the same as that of the ground state. 
Now, using the replica trick, we can show that for vertex operators, this contribution coming from current $J = i \p \phi$ exactly cancels the piece coming from $D_n(x)$. The contribution from the two point functions is proportional to $S_2(n,1)$. We compute it below,
\begin{align}
    S_2(n,\Delta) &= \frac{n}{2}\sum_{m\in \mathbb{Z}} \left[ng_\Delta(nm) - g_\Delta(m)\right] 
\end{align}
We take $\Delta \sim 1+\varepsilon$ in order to deal with singularities of Gamma functions (Note that, this $\varepsilon$ has nothing to do with the primary in Ising CFT). Then we have,
\begin{align}
    g_{1+\varepsilon}(m) &= 2^2 \Gamma(-1-2\varepsilon)\frac {\sin(\pi (1+\varepsilon))}{\pi}\frac{\Gamma (m+1)}{\Gamma(m)} 
    \\
    &\simeq 4 \lb \frac{1}{2\varepsilon} \rb \lb \frac{-\pi \varepsilon}{\pi}\rb m \\
    &= -2m .
\end{align}
Similarly, we have
\begin{align}
    g_{1+\varepsilon}(nm) = -2nm.
\end{align}
We can then write the sum of all two point functions to be,
\begin{align}
    S_2(n,1) &= 2^{-2}\frac{n}{2}\sum_{m\in \mathbb{Z}} \left[ n(2nm) + 2m \right]\\
    &= -2^{-2}n(n^2-1) 2\sum_{m\in \mathbb{Z}} m \\
    &= -2^{-2}n(n^2 - 1) \zeta(-1) \\
    &= \frac{n(n^2-1)}{24}
\end{align}
Now adding the contribution from both holomorphic and anti-holomorphic currents, we get the total contribution to be,
\begin{align}
    2C_\Psi^2 \lb\frac{2\pi x}{n}\rb^2 \frac{n(n^2-1)}{24} = C_\Psi^2\frac{(\pi x)^2}{3} \lb n-\frac{1}{n}\rb
\end{align}
For coherent states, $C_\Psi = \alpha$. We finally get, 
\begin{align}
    F_\Psi^{(n)} &= 1+ \frac{\alpha^2}{3} \lb\frac{1}{n}-n\rb (\pi x)^2 + \frac{\alpha^2}{3}\lb n-\frac{1}{n}\rb(\pi x)^2 +  \cdots\\
    &= 1 + \cdots
\end{align}
where the $\cdots$ are the higher order terms coming from higher point functions. We essentially showed the cancellation of the terms of the order of $\mO(x^2)$ explicitly. This could be further extended by computing higher order terms in the 2n-point function. 
\bibliographystyle{JHEP}
\bibliography{biblio}

\end{document}